% ============================================================================
% main.tex — Verified constructions, a classification programme, and
% conditional reductions for Gallai's path decomposition conjecture.
% Driver: inputs the five section files; bibliography in refs.bib.
% ============================================================================
\documentclass[11pt]{article}
\usepackage[margin=1in]{geometry}
\usepackage{amsmath,amssymb,amsthm}
\usepackage{enumitem}
\usepackage{booktabs}
\usepackage{longtable}
\usepackage[colorlinks=true,linkcolor=blue,citecolor=blue,urlcolor=blue]{hyperref}

\newcommand{\pn}{p}
\DeclareMathOperator{\odd}{odd}
\newcommand{\EV}{\mathcal{E}}   % even-degree subgraph of a graph G
\newcommand{\ed}{d_{\mathrm{E}}} % E-degree inside EV(G)
\newcommand{\dEV}{d_{\mathrm{E}}} % E-degree inside EV(G), section-file alias

% Hungarian double-acute accent (\H{o} = U+0151), not predefined in the
% LaTeX kernel; XeTeX renders the precomposed character natively.
\providecommand{\H}[1]{ő}
\newcommand{\Z}{\mathbb{Z}}
\newcommand{\ceil}[1]{\left\lceil #1 \right\rceil}
\newcommand{\floor}[1]{\left\lfloor #1 \right\rfloor}

\theoremstyle{plain}
\newtheorem{theorem}{Theorem}[section]
\newtheorem{lemma}[theorem]{Lemma}
\newtheorem{proposition}[theorem]{Proposition}
\newtheorem{corollary}[theorem]{Corollary}
\newtheorem{conjecture}[theorem]{Conjecture}
\theoremstyle{definition}
\newtheorem{definition}[theorem]{Definition}
\newtheorem{observation}[theorem]{Observation}
\newtheorem{example}[theorem]{Example}
\theoremstyle{remark}
\newtheorem{remark}[theorem]{Remark}

% License note: the paper text is released under CC BY 4.0; the repository
% code and scripts are released under the MIT License (see LICENSE).
\newcommand{\licensernote}{The text of this paper is released under the
Creative Commons Attribution 4.0 International License
(\url{https://creativecommons.org/licenses/by/4.0/}). The accompanying
repository code and verification scripts are released under the MIT License.}

% Convention for claim strength, used throughout:
% "Theorem" = complete proof in this paper; "Proposition/Corollary" likewise.
% Candidate statements are stated as Conjecture or as clearly flagged
% candidate Propositions with their tier sentence attached; machine accounts
% are quoted as accounts, never as theorems.
%

\title{Verified constructions, a classification programme, and conditional
reductions for Gallai's path decomposition conjecture}

\author{Independent Researcher\thanks{Contact: 672178818@qq.com. Parts of the
mathematical and bibliographic work, including automated literature
cross-checking and machine verification of the computational accounts, were
carried out with computer assistance; the division of labour is summarised in
Section~\ref{sec:computation}.}}
\date{October 2026}

\begin{document}
\maketitle

\begin{center}\small\licensernote\end{center}

\begin{abstract}
Gallai conjectured that the edges of every connected graph on $n$ vertices can
be decomposed into at most $\lceil n/2\rceil$ paths. We present a self-contained
study of the conjecture organised around three kinds of results. First, exact
constructions: we prove that the path number of $K_{4,n}$ equals $(n+1)/2$ for
every odd $n\ge 3$, develop a structure theory and a ring-closure theorem that
reduce the remaining even columns of the odd-$n_1$ case of the
Constantinou--Ellinas problem to finitely many base instances (closing the
columns $n_2=6$, $8$ and $10$ except five named bases), and prove a chain
theorem for odd clique chains, giving $p(G)\le\lfloor n/2\rfloor$ for $b\ge2$
blocks and $p(G)\le\lfloor n/2\rfloor-1$ for $b\ge4$, via a new decomposition
lemma for complete graphs of odd order minus an arbitrarily prescribed
Hamilton path. Second, a classification programme for the sharp placement
problem on complete graphs $K_{2k+1}$: we prove the prime-family theorem that
for $n=p+1$ with $p\ge7$ a prime congruent to $3$ mod $4$, every non-infinity
placement regime is dead, together with corrected Fourier identities for the
infinity branch, and we register the evidence structure of the exception-set
candidate, whose uniqueness clause is already in tension at $k=21$ with the
propagation theorem (a withdrawn quantitative registration is corrected to the
paper-level bound $|\mathcal{E}_{21}|\ge1$). Third, conditional reductions and
formalisation: an endpoint-budget framework whose budget conjecture $K$ implies
Gallai's conjecture, a conditional chain registering Gallai against an explicit
eight-item premise ledger, per-domain theorem-form closure of the aligned
valuation criterion for $N\in\{6,8,10,12\}$ with the general-$N$ existence half open,
Lean formalisations of $p(K_{4,n})=(n+1)/2$ (odd $n$), $p(K_{2k+1})=k+1$ and a
criterion-level termination bound, and a machine-account layer for a six-family
covering identity, all quoted with explicit one-way (superset) reading
conventions. Every claim is stated at its verified strength; open problems and
residual gaps are collected at the end.
\end{abstract}

\tableofcontents

% ============================================================================
% sec-intro.tex — Introduction, notation, and the shape of the results.
% Label prefix: int: (all labels owned by this file).
% ============================================================================
\section{Introduction, notation, and the shape of the results}\label{sec:intro}

Let $G$ be a finite connected simple graph on $n$ vertices. A \emph{path decomposition} of
$G$ is a partition of $E(G)$ into paths, and $\pn(G)$ denotes the minimum number of paths in
such a decomposition. In his 1966 Tihany lecture, published as \cite{Gallai68}, Gallai
conjectured that $\pn(G)\le\ceil{n/2}$ for every connected $G$; the conjecture was
communicated by Erd\H{o}s and is tracked as problem \#583 in \cite{EP583}. It remains open,
with $\floor{2n/3}$ the best general upper bound, due to Dean and Kouider \cite{DK00}.

The case analysis of the last six decades has been organized, to a remarkable degree, by the
structure of the subgraph $\EV(G)$ induced by the even-degree vertices, together with bounds
on the degrees of $G$ itself. Lov\'asz \cite{Lo68} showed that $E(G)$ decomposes into at most
$\floor{n/2}$ paths and cycles, and settled the case $|\EV(G)|\le1$; Pyber \cite{Pyb96}
proved that $\EV(G)$ a forest implies $\pn(G)\le\floor{n/2}$, with the general bound
$n/2+O(n^{3/4})$; Favaron and Kouider \cite{FK88} settled Eulerian graphs of maximum
degree~$4$; and Fan \cite{Fan05} proved the conjecture when every block of $\EV(G)$ is
triangle-free with maximum degree at most~$3$. Two strengthenings followed: Botler and
Sambinelli \cite{BS21} proved it for graphs whose $\EV(G)$ is a subgraph of a graph in their
class $\mathcal{G}$ (each block of $\EV(G)$ of maximum degree at most~$3$, at most one
triangular block per component; their main theorem is stated with Gallai's ceiling
$\ceil{n/2}$ rather than the floor of Theorem~\ref{int:thm-fan} below), and Fan, Hou and
Zhou \cite{FHZ23} proved it when $\EV(G)$ contains no pair of triangles and every block of
$\EV(G)$ has maximum degree at most~$3$. A parallel track bounds the degrees of $G$ itself:
$\Delta\le5$ \cite{BP19} and treewidth at most~$3$ \cite{BCLS20}. The class-based programme then accelerated: Zhang, Liu and Hong \cite{ZLH24}
proved the conjecture, with the sharper $\floor{n/2}$ bound, for all $3$-degenerate graphs
except $K_3$ and $K_5^-$; Lu and Yi \cite{LY24} did the same for treewidth at most~$4$ with
exceptions $K_3$, $K_5^-$, $K_5$; and Chen and Wu \cite{CW25} settled it for block graphs.
On the algorithmic side, exact path partitions on subcubic graphs are computable in
polynomial time \cite{MY26}. Two further general results bound the problem from above: for
triangle-free graphs, Chu, Fan and Zhou \cite{CFZ22} proved the global bound
$\pn(G)\le 3n/5$; for the \emph{odd semi-cliques} (Definition~\ref{int:def-osc}), the
question raised by Bonamy and Perrett \cite{BP19} whether $\pn(G)\le\floor{n/2}$ holds for
every connected graph that is not an odd semi-clique has its first substantial progress in
Chu, Fan and Zhou \cite{CFZ26}: if the E-subgraph ($\EV(G)$) is isomorphic to a complete
graph $K_m$ with $m\le15$, then $\pn(G)\le\floor{n/2}+1$. Finally, minimal path
decompositions of complete bipartite graphs were studied by Constantinou and Ellinas
\cite{CE18}, whose algorithms cover three of four parameter cases; we return to the open
case below.

\subsection{Blocks versus global maximum degree}\label{int:subsec-blockvsglobal}

The Fan-line conditions quoted above are \emph{block-level}: they constrain the blocks of
$\EV(G)$ one by one, not $\EV(G)$ as a whole. The distinction is not pedantic; the two
readings are not equivalent, in either direction. In $K_{1,5}$ with one pendant leaf attached
at each of its six vertices, the even-degree subgraph is $\EV(G)\cong K_{1,5}$, which has a
vertex of E-degree~$5$, while every block of $\EV(G)$ is a single edge (maximum degree~$1$):
a large \emph{global} degree can coexist with a Fan-line hypothesis satisfied block by block.
Conversely, a cut vertex of $\EV(G)$ may belong to arbitrarily many blocks each of maximum
degree at most~$3$, so the absence of high E-degree vertices is stronger than the block
condition. Accordingly, the correct reading of the twenty-year-old Fan-line hypothesis is:
\emph{every block of $\EV(G)$ is triangle-free and has maximum degree at most~$3$}
\cite{Fan05}, read through Theorem~2(c) of \cite{BS21}, with the forest case of
Theorem~2(b) of \cite{BS21} going back to Pyber \cite{Pyb96}. The theorems of
Section~\ref{subsec:astar} are new precisely at the block level: they admit a \emph{block} of
$\EV(G)$ of maximum degree at least~$4$.

\subsection{Contributions}\label{int:subsec-contrib}

The paper has four mathematical parts (Parts~I--III being constructive and Part~IV
conditional and machine-backed) and a verification apparatus. We list the contributions,
with pointers to the sections that carry them.

\begin{enumerate}[label=(\roman*),leftmargin=2em]
\item \textbf{The exact path number of $K_{4,n}$ for odd $n$, and a decision procedure for
the even columns of Case~4 of \cite{CE18}} (Section~\ref{subsec:k4odd}). Constantinou and
Ellinas \cite{CE18} partition the parameter pairs of $K_{n_1,n_2}$ into four cases and give
minimal-decomposition algorithms for three of them, leaving open Case~4 (odd $n_1$,
$1\le n_2\le n_1-1$) with an explicit request for its ``exact characteristics''. We close the
$n_2=4$ column completely:
\[
\pn(K_{4,n})=(n+1)/2 \qquad\text{for every odd } n\ge3,
\]
improving the best published bound for the family (the published bound $k+2$, where
$n=2k+1$, from \cite{Pyb96,Fan05,LY24}) by one. The value for odd $n\ge3$ is formalized in Lean as a
single theorem with axiom profile \texttt{[propext, Classical.choice, Quot.sound]} and no
\texttt{sorryAx}; numeric corollaries at $n=5,13,19$ follow from it without independent
gaps. For the remaining even columns we prove a general counting cap, a structure theorem
for optimal decompositions across the whole even range, and a \emph{ring-closure theorem}
(one explicit base decomposition of $K_{2k_0+1,2h}$ into $k_0+1$ paths extends, by an
explicit construction, to every $k\ge k_0+h+1$); the column $n_2=6$ closes completely, the
columns $n_2=8,10$ close except for five named base instances, and the conjecture that
$\pn(K_{n_1,n_2})=\ceil{n_1/2}$ throughout the even range (Conjecture~\ref{conj:con27A} in the
body) becomes a column-by-column decision procedure with $h+1$ base instances per column.

\item \textbf{Odd clique chains: refinements beyond block graphs, with the domain fixed}
(Section~\ref{subsec:chains}). Chen and Wu \cite{CW25} proved Gallai's conjecture for all block
graphs. For the subfamily of \emph{odd clique chains} --- block graphs whose block tree is a
path and whose blocks are $B_i\cong K_{2j_i+1}$ with \emph{all $j_i$ odd}, i.e.\ blocks of
order $\equiv 3 \pmod 4$, $K_3$ allowed --- we prove
\[
\pn(G)\le \floor{n/2}\ \ (b\ge2), \qquad \pn(G)\le \floor{n/2}-1\ \ (b\ge4),
\]
where $b$ is the number of blocks and $n=2\sum_i j_i+1$, along with the exact value
$\pn(G)=\ceil{n/2}$ for a single block. Two qualifications fix the domain precisely. First,
\emph{every} vertex of such a chain has even degree: a non-cut vertex of a block
$K_{2j_i+1}$ has degree $2j_i$, and a cut vertex shared by consecutive blocks has degree
$2j_i+2j_{i+1}$; hence $\EV(G)=G$. Second, the complementary case --- chains containing
blocks of order $1 \pmod 4$ (even $j_i$, e.g.\ $K_5$) --- is \emph{open}: no general
construction is known, and in-chain existence is machine-verified only for the block values
$K_5$ and $K_9$. For $b=3$ the $-1$ refinement is likewise open and reduces to a
distinctness question on the middle block. The local engine is: \emph{for every odd $j$ and
every Hamilton path $P$ of $K_{2j+1}$, the graph $K_{2j+1}-P$ decomposes into $j$ paths};
the base zigzag objects are classical \cite{LP92,Alspach08}, the statement for an
arbitrarily prescribed Hamilton path appears to be new.

\item \textbf{Fan-type theorems admitting a block of $\EV(G)$ of maximum degree $\ge4$}
(Section~\ref{subsec:astar}). The Fan line \cite{Fan05,BS21,FHZ23} has kept its requirement at
the block level for twenty years: every block of $\EV(G)$ has maximum degree at most~$3$.
We prove class theorems that admit a block of $\EV(G)$ of maximum degree at least~$4$:
Theorems~\ref{thm:claAstar} and~\ref{thm:lasep} (if $\EV(G)$ is triangle-free, high E-degrees are compensated by an
odd-degree condition, and an explicit alignment condition holds at one high vertex, then
$\pn(G)\le\floor{n/2}$; and the same bound with no alignment hypothesis
whenever distinct high vertices have no common EV-neighbour), and a two-high-vertex theorem
with a corner-case reduction. An honesty note travels with these results: our witness
family $W(m)$ ($m$ odd, $\ge5$) has $\EV(W(m))\cong K_{2,m}$, a single block of maximum
degree $m$, and has treewidth $5$, degeneracy $5$ and maximum degree at least $m+1$; it
therefore lies \emph{outside} the degeneracy and treewidth classes of \cite{ZLH24,LY24} and
outside every published class we are aware of, so our theorems and the class-based line
\cite{ZLH24,LY24,BCLS20,BP19,CW25} are incomparable rather than competitive. The general
statement is posed as Conjecture~\ref{conj:claAstar}.

\item \textbf{A conditional route: Gallai modulo an explicit premise ledger} (Part~IV,
Sections~\ref{sec:computation}--\ref{sec:open}). We develop an endpoint-budget framework: the
per-vertex \emph{endpoint budget} $\beta(G,x)$ (Definition~\ref{int:def-budget}) and the
\emph{budget conjecture} $K$ that $\beta(G,x)\le\ceil{n/2}$ always. On the proved side we
establish an absorption-cost viewpoint (Remark~\ref{rem:absorb}) unifying three obstruction
phenomena from the literature; that $K$ implies Gallai's conjecture
(Theorem~\ref{int:thm-Kgallai}, proved here), so $K$ is a genuine strengthening; and a
conditional chain (Theorem~\ref{int:thm-ledger}, stated here and carried at pointer level in
the ledger subsection) assembling
$\text{Gallai}\Leftarrow K\Leftarrow\text{degree-two layer of }K\Leftarrow
\text{SC-PH}\Leftarrow\text{a placement conjecture}\Leftarrow\text{(E)+(R)}$, with the
bridge from the valuation line (the criterion of Section~\ref{subsubsec:criteria}) carried as a
parallel, redundant-for-implication route. The chapter's central deliverable is an \emph{explicit
hypothesis ledger}: an eight-item premise set, restated descriptively in
Section~\ref{int:subsec-conditional} below. The premise set mixes three kinds of conjuncts:
external steps flagged as unverified against the journal record (the all-odd support of the
Lov\'asz decomposition theorem as consumed by our proof of
Theorem~\ref{int:thm-Kgallai}), open conjectures internal to the paper (BP-odd, the
odd-order density conjecture, SC-PH), and paper-open items with machine backing. No premise
is claimed to be proved.

\item \textbf{Verification methodology} (Sections~\ref{sec:computation} and the formalization
section, Appendices). All exact values and all decompositions cited in this paper were produced and
re-checked by an independent machine pipeline (a bitmask dynamic program plus a
decomposition verifier asserting that the parts are simple paths and exactly partition the
edge set). The large-scale machine accounts of Part~IV are quoted at their measured
strength, under two standing conventions. First, the \emph{one-way superset convention}:
every machine check of the covering identity is performed in a reading that omits one
definitional branch (the \texttt{v7} branch, whose machine face is frozen because the
legacy candidate-generation source files are lost), giving a conservative superset
$\mathrm{fam}_{\mathrm{noV7}}\subseteq\mathrm{Im}\cap L$; zero counterexamples on the
superset face strictly imply zero counterexamples on the true face, and not conversely (the
residual \texttt{v7}-only mass is nonzero and precisely measured on the tested domains).
Second, the \emph{accounting layers are kept separate}: the nine-domain key tables
(comprising $2{,}412{,}436$ keys and $98{,}602{,}396$ key-weighted positions over
$N\in\{6,\dots,14\}$) are counted key-weighted, each key once, while the per-instance
accounts at $N=8$ and $N=9$ are instance-weighted; the two accountings are never mixed, and
the $N=14$ domain is an affine-normalized domain whose equivalence to exhaustive enumeration
is unproven. Lean certificates carry their trust boundaries verbatim: theorems proved from
first principles are reported with their axiom profiles, and finite certificates checked by
\texttt{native\_decide} carry the kernel-reflection axiom \texttt{Lean.ofReduceBool}
explicitly.
\end{enumerate}

\subsection{Claim-strength declaration}\label{int:subsec-claims}

Six no-claim constraints bind the whole account, including this introduction; we state them
verbatim. We do \emph{not} claim:
\begin{enumerate}[label=\textup{(\arabic*)},leftmargin=2em]
\item closure at $L\ge13$ (machine evidence at larger $N$ is extrapolation evidence, not
proof; only candidate-level status is claimed there);
\item general-$N$ paper closure of the covering identity (the general-$N$ existence half
stays open);
\item that the four-piece (E)+(R) landing is complete;
\item paper closure of the general-$L$ uniform-construction account (the
capacity-starved branch of the channel-or-menu dichotomy programme);
\item full closure of that dichotomy programme;
\item an unconditional theorem-form statement of the criterion of
Conjecture~\ref{conj:clacand}. The Lean theorem \texttt{thmA2663\_upgraded} is a
\emph{criterion-level bound}: under the explicit hypotheses that the target side has the
prescribed size and that $\mathrm{val}\le1$, it proves $g_{13}\le d_2+1$ with axiom profile
\texttt{[propext, Quot.sound]} and no \texttt{sorryAx}; the hypothesis $\mathrm{val}\le1$
itself is \emph{not} proved by the theorem, and a per-instance criterion verified over
finitely many $N$ is not a feasibility-existence proof for general $N$.
\end{enumerate}

\subsection{Notation and preliminaries}\label{int:subsec-notation}

All graphs are finite, simple and undirected. A \emph{path} has at least one edge and no
repeated vertex; a \emph{path decomposition} of $G$ is a set $D$ of paths whose edge sets
partition $E(G)$; $\pn(G)$ is the minimum size of a path decomposition. For a decomposition
$D$ and a vertex $x$ we write $D(x)$ for the number of paths of $D$ having $x$ as an
endpoint. A \emph{block} of $G$ is a maximal $2$-connected subgraph or a bridge (together
with its endpoints); every two blocks share at most one vertex, a \emph{cut vertex}. $G$ is
\emph{$k$-degenerate} if every subgraph has a vertex of degree at most $k$. The
\emph{even-degree subgraph} $\EV(G)$ is the subgraph induced by the vertices of even degree;
for $v\in V(G)$ we write $\dEV(v)$ for the degree of $v$ inside $\EV(G)$, the \emph{E-degree}
of $v$; a vertex with $\dEV(v)\ge4$ is called \emph{high}. Note that $\dEV(v)\ge4$ implies
in particular that $v$ has even degree in $G$.

We collect the three folklore lower bounds used throughout; proofs are included for
completeness.

\begin{lemma}[odd-degree lower bound; endpoint parity]\label{int:lem-LB1}
Let $D$ be a path decomposition of $G$. Then:
\begin{enumerate}[label=(\roman*),leftmargin=2em]
\item $2|D|\ge \mathrm{odd}(G)$, where $\mathrm{odd}(G)$ is the number of odd-degree
vertices; in particular $\pn(G)\ge \mathrm{odd}(G)/2$;
\item every odd-degree vertex is an endpoint of at least one path of $D$;
\item for every vertex $v$, $D(v)\equiv \deg_G(v)\pmod 2$.
\end{enumerate}
\end{lemma}

\begin{proof}
For a path $P$ and a vertex $v$, the degree of $v$ in $P$ is $0$, $1$ or $2$, and it is $1$
exactly when $v$ is an endpoint of $P$. Summing over $D$,
$\deg_G(v)=\sum_{P\in D}\deg_P(v)\equiv D(v)\pmod 2$, which gives (iii); (ii) is (iii) for
odd $\deg_G(v)$; and (i) follows by charging each odd-degree vertex to one path ending at
it, each path being charged at most twice.
\end{proof}

\begin{lemma}[edge-count lower bound]\label{int:lem-LB2}
If $G$ is connected with $n\ge2$ vertices and $m$ edges, then
$\pn(G)\ge \ceil{m/(n-1)}$.
\end{lemma}

\begin{proof}
Each path has at most $n-1$ edges, so $k$ paths cover at most $k(n-1)$ edges.
\end{proof}

\begin{lemma}[disconnected sum]\label{int:lem-sum}
If $G$ has components $G_1,\dots,G_c$ on $n_1,\dots,n_c$ vertices, then
$\sum_i \floor{n_i/2}\le \floor{n/2}$.
\end{lemma}

\begin{proof}
Immediate from $\sum_i \floor{n_i/2}\le \floor{\sum_i n_i /2}$ for integers.
\end{proof}

The three standard tools from the literature, stated in the form used later:

\begin{theorem}[Lov\'asz \cite{Lo68}]\label{int:thm-lovasz}
Every $n$-vertex graph (not necessarily connected) has a decomposition of its edge set into
at most $\floor{n/2}$ paths and cycles.
\end{theorem}

\begin{theorem}[Pyber \cite{Pyb96}]\label{int:thm-pyber}
If $\EV(G)$ is a forest, then $\pn(G)\le \floor{n/2}$.
\end{theorem}

\begin{theorem}[Fan \cite{Fan05}]\label{int:thm-fan}
If every block of $\EV(G)$ is triangle-free with maximum degree at most $3$, then
$\pn(G)\le \floor{n/2}$.
\end{theorem}

\begin{definition}[odd semi-clique; {\cite{BP19}}]\label{int:def-osc}
An \emph{odd semi-clique} is a graph obtained from the complete graph on $2k+1$ vertices by
deleting at most $k-1$ edges.
\end{definition}

\begin{observation}[density obstruction]\label{int:obs-density}
An odd semi-clique on $n=2k+1$ vertices has
$m\ge \binom{2k+1}{2}-(k-1)=2k^2+1$ edges, so by Lemma~\ref{int:lem-LB2}
$\pn\ge \ceil{(2k^2+1)/2k}=k+1=\ceil{n/2}$; in particular no odd semi-clique admits a
$\floor{n/2}$-path decomposition.
\end{observation}

Finally we record the classical Walecki partition of the complete graph of odd order, in
the form used in Section~\ref{subsec:chains}; the base objects are classical
\cite{LP92,Alspach08}. Fix $j\ge 1$ and let $V=\{\infty\}\cup \mathbb{Z}_{2j}$. For
$k\in\mathbb{Z}_{2j}$ let $C_k$ be the cycle
\[
C_k \;=\; \infty,\, k,\, k+1,\, k-1,\, k+2,\, k-2,\;\dots,\, k+(j-1),\, k-(j-1),\, k+j,\, \infty
\]
(indices mod $2j$); equivalently the finite vertices of $C_k$ are $p_0=k$ and
$p_s=k+(-1)^{s+1}\ceil{s/2}$ for $s=1,\dots,2j-1$.

\begin{lemma}[Walecki partition; classical, cf.\ \cite{LP92,Alspach08}]\label{int:lem-walecki}
For $k=0,1,\dots,j-1$ the cycles $C_k$ are edge-disjoint and
$\bigsqcup_{k=0}^{j-1}E(C_k)=E(K_{2j+1})$.
\end{lemma}

\begin{proof}
Each $C_k$ is a cycle: the finite vertices $k, k\pm1,\dots,k\pm(j-1),k+j$ are $2j$ distinct
elements covering $\mathbb{Z}_{2j}$. We show every edge lies in exactly one $C_k$, $k$ taken
mod $j$. Write $\delta(x,y)$ for the circular distance of $x,y\in\mathbb{Z}_{2j}$, i.e.\ the
minimum of $\pm(x-y)$ mod $2j$ in $\{0,\dots,j\}$.

Reading off consecutive differences in the sequence $p_s$ (as above,
$p_{s+1}-p_s=\pm(s+1)$), one checks that the finite edges of $C_k$ are exactly the pairs
$\{k+t,\,k-t\}$ for $t\in\mathbb{Z}_{2j}\setminus\{0,j\}$ and the pairs
$\{k+t,\,k-t+1\}$ for $t\in\mathbb{Z}_{2j}$, together with the two $\infty$-edges
$\{\infty,k\}$, $\{\infty,k+j\}$. (In each of the two forms an edge is produced twice, once
from each endpoint, so it is counted once; $t=0$ and $t=j$ would give loops in the first
form and are excluded.) Hence $C_k$ has, among finite edges, two edges of each circular
distance $1,\dots,j-1$ and one edge of distance $j$.

Now fix an edge $e=\{x,y\}$ of $K_{2j+1}$ and let $d=\delta(x,y)$.
\begin{itemize}[leftmargin=1.5em]
\item If $y=x+j$ (the diameter case, $d=j$): for $j$ odd, $e$ has the form
$\{k+t,k-t+1\}$, which forces $2k\equiv 2x+d-1\pmod{2j}$, a unique solution mod $j$ (the
right side is even); for $j$ even, $e$ has the form $\{k+t,k-t\}$ with $2t\equiv d=j$,
forcing $2k\equiv 2x+j\pmod{2j}$, again a unique solution mod $j$. So $e$ lies in exactly
one $C_k$.
\item If $d<j$ is even: $e=\{x,x+d\}=\{k+t,k-t\}$ forces $2k\equiv 2x+d\pmod{2j}$, unique
mod $j$.
\item If $d<j$ is odd: $e=\{x,x+d\}=\{k+t,k-t+1\}$ forces $2k\equiv 2x+d-1\pmod{2j}$,
unique mod $j$.
\item If $e=\{\infty,x\}$: $e\in C_k$ iff $x\in\{k,k+j\}$, i.e.\ $k\equiv x$ or $x+j$
(mod $2j$), a unique class mod $j$.
\end{itemize}
In each case $e$ lies in exactly one of $C_0,\dots,C_{j-1}$, and each $C_k$ has $2j+1$
edges, so the $j(2j+1)=\binom{2j+1}{2}$ edges are partitioned.
\end{proof}

\subsection{Notation for the machine chapters}\label{int:subsec-machnotation}

The machine chapters (Section~\ref{sec:computation} and the formalization material of
Section~\ref{sec:open}) work with finite
combinatorial instances; we fix the notation once here, transcribed so that the chapter
statements are readable without further apparatus.

\begin{definition}[instances]\label{int:def-instance}
An \emph{instance} is $I=(N,P,T,d)$: $\mathbb{Z}_N:=\mathbb{Z}/N\mathbb{Z}$;
$P\subseteq\mathbb{Z}_N$ with $|P|=p$; $T=(t_1,t_2,t_3)$, three pairwise distinct elements
of $\mathbb{Z}_N$; $d=(d_1,d_2,d_3)$ with $d_i\ge1$; and the capacity surplus
$\beta:=p-d_1-d_2-d_3\ge1$ (capacity admission: $d_1\le p-[t_1\in P]$). The shifts are
$\delta_{12}:=t_1-t_2$, $\delta_{13}:=t_1-t_3$, $\delta_{23}:=t_2-t_3$ (mod $N$). The
instance is \emph{aligned} (\emph{so1}) when, with $g:=\gcd(\delta_{12},N)$ and
$L:=\mathrm{ord}(\delta_{12})$, one has $\delta_{13}\equiv k\,\delta_{12}\pmod N$ for the
canonical $k\in[2,L)$; equivalently, $t_1,t_2,t_3$ lie on one coset (``line'') of
$\langle\delta_{12}\rangle$.
\end{definition}

\begin{definition}[legal pairs and the value functional]\label{int:def-legalval}
For $C_1\subseteq P\setminus\{t_1\}$ put $B_{12}:=C_1-\delta_{12}$ and
$Q(C_1):=P\setminus(C_1\cup(B_{12}\cap P)\cup\{t_2\})$. A \emph{legal pair} is
$(C_1,C_2)$ with $C_1\subseteq P\setminus\{t_1\}$, $|C_1|=d_1$, $C_2\subseteq Q(C_1)$,
$|C_2|=d_2$. With $F_{13}:=(C_1-\delta_{13})\cap P$ and $F_{23}:=(C_2-\delta_{23})\cap P$,
the \emph{value} of a legal pair is the three-term count
\[
\mathrm{val}(C_1,C_2)
 :=\bigl|F_{13}\setminus(C_1\cup C_2)\bigr|
  +\bigl|F_{23}\setminus(C_1\cup C_2\cup F_{13})\bigr|
  +\bigl|\{t_3\}\cap P\setminus(C_1\cup C_2)\bigr| .
\]
The \emph{min-val} of the instance is the minimum of $\mathrm{val}$ over all legal pairs.
\end{definition}

The \emph{so1 (aligned) family} is the family of aligned instances; it is the family on
which the paper layer of the conditional chain is formulated, and the machine accounts of
Section~\ref{sec:computation} sweep exactly the statement ``every so1 instance with $\beta\ge2$
has $\min\mathrm{val}\le1$'' on the explicitly tested domains (the sweep domain $\beta\ge1$
contains the $\beta\ge2$ slice). The candidate form of that statement is registered as
Conjecture~\ref{conj:clacand} in Section~\ref{subsubsec:criteria}; its general-$N$ existence
half is Conjecture~\ref{int:conj-gapex}.

\subsection{Two reinsertion lemmas from the literature}\label{int:subsec-fanlemmas}

The Fan-type arguments of Section~\ref{subsec:astar} and the budget framework of
Section~\ref{int:subsec-conditional} consume two deletion--reinsertion lemmas, which we
state in the form used throughout; both follow the restatements in \cite{BS21}
(Lemmas~6 and~7 there), and we record our provenance boundaries with them.

\begin{lemma}[addible edge; {\cite{Fan05},\cite{BS21}}]\label{int:lem-fan6}
Let $uv\in E(G)$ and let $D'$ be a path decomposition of $G-uv$. If
$D'(v)>\ |\{w\in N_{G-uv}(u) : D'(w)=0\}|$, then $G$ has a path decomposition $D$ with
$|D|=|D'|$, $D(u)=D'(u)+1$, $D(v)=D'(v)-1$, and $D(w)=D'(w)$ for all other $w$. One says that
$uv$ is \emph{addible towards $u$}.
\end{lemma}

\begin{remark}[transcription, downgrade and machine evidence]\label{int:rem-fan6}
The proof of Lemma~\ref{int:lem-fan6} in \cite{Fan05} could not be consulted; the statement
follows the restatement in \cite{BS21} (Lemma~6) and has not been independently checked
against the original. A strengthening of the lemma that additionally requires the
underlying re-arrangement to preserve the endpoint profile of $D'$ pointwise is false in
general (a five-vertex counterexample is machine-certified); the lemma as stated imposes no
such requirement and has no known counterexample, and the Lean formalization used in
Section~\ref{sec:computation} carries the four-clause form obtained by deleting the final
pointwise clause $D(w)=D'(w)$ for the remaining vertices --- no application in this paper
consumes that clause. The machine evidence is direction-sensitive and has been
direction-corrected: the bidirectional rescan is complete at $n\le6$, all
$1{,}701{,}141{,}150$ ordered instances at $n=6$ (the fixed-orientation half exactly
reproducing the historical census of $850{,}570{,}575$), with zero violations
(Section~\ref{com:subsec-pipeline}).
\end{remark}

\begin{lemma}[star absorption; {\cite{Fan05},\cite{BS21}}]\label{int:lem-fan7}
Let $B=\{uw_1,\dots,uw_h\}$ be a set of $h$ edges incident to $u$, let $D'$ be a path
decomposition of $G-B$ in which every neighbour of $u$ \emph{in $G$} has $D'\ge 1$, and let
$uw$ be a specified edge of $B$. Then there is a set $A\subseteq B$ with $uw\in A$ and
$|A|\ge\ceil{h/2}$ such that the edges of $A$ can be added back one by one, each by
Lemma~\ref{int:lem-fan6} towards $u$, preserving $|D|$ and increasing $D(u)$ by one at each
step.
\end{lemma}

\begin{remark}\label{int:rem-fan7}
The hypothesis of Lemma~\ref{int:lem-fan7} is read over the neighbours of $u$ \emph{in
$G$}, so that it includes the star vertices $w_i$ themselves, which are not neighbours of
$u$ in $G-B$; this stronger reading is the one verified in every application below. As with
Lemma~\ref{int:lem-fan6}, the proof in \cite{Fan05} (Lemma~3.6 there) could not be
consulted, and the transcription from the restatement in \cite{BS21} (Lemma~7) has not been
independently checked against the original.
\end{remark}

\subsection{The absorption-cost accounting and the budget conjecture}\label{int:subsec-budget}

Several proofs in the literature --- Lov\'asz' deletion--reinsertion, Fan's star absorption,
and the analysis of minimal counterexamples --- fail at the same type of point: one needs a
decomposition of a reduced graph in which \emph{specified vertices hold sufficiently many
path endpoints}, and the count of available endpoints is not controlled by the usual
hypotheses. The following elementary accounting fact, quoted from the project's
verification work, quantifies the obstruction class once and is used only as a viewpoint in
this paper.

\begin{remark}[absorption cost]\label{rem:absorb}
Let $G$ be a graph, $D$ a path decomposition of $G$, $v$ a vertex not incident with any
edge of $G$ ($v$ may be a new vertex), and $N=\{w_1,\dots,w_d\}\subseteq V(G)$ a set of $d$
distinct vertices; let $G^+=G+\{vw: w\in N\}$. Build the bipartite graph $B(D)$ with left
side the $d$ edges $vw$ of the star and right side the paths of $D$, where $vw$ is adjacent
to $P$ iff $P$ has $w$ (or $v$) as an endpoint; let $\mu^*$ be a maximum matching of
$B(D)$. Within the operation class --- (i) attach a star edge $vw$ to the $w$-end of a path
of $D$ ending at $w$; (ii) pair two star edges $vw,vw'$ into the single path $w$--$v$--$w'$;
(iii) let a star edge be a single-edge path --- the minimum number of new paths required to
extend $D$ to a decomposition of $G^+$ is exactly $\ceil{(d-\mu^*)/2}$, and the resulting
$D(v)$ can be made at least $\mu^*$. The equality holds within the operation class and is
\emph{not} a lower bound for arbitrary decompositions of $G^+$. Three obstruction
phenomena in the literature are instances of ``the cost exceeds the budget'': the
Lov\'asz-style vertex reinsertion gap (the unproved strengthening of
Theorem~\ref{int:thm-lovasz} to path-only decompositions with endpoint control), the parity
wall for even-sized stars in the Fan engine (a star of even size leaves its centre even, so
no endpoint at the centre is forced), and the simplicity wall of Lemma~\ref{int:lem-fan6} (a
host endpoint may exist but be unusable because the path through it already contains the
reinsertion target).
\end{remark}

\subsection{The shape of the conditional route}\label{int:subsec-conditional}

The conditional chapters measure, at each vertex, the one quantity that the classical
deletion--reinsertion arguments leave uncontrolled: the number of path endpoints a
decomposition can be made to hold at a specified vertex.

\begin{definition}[endpoint budget]\label{int:def-budget}
For a connected graph $G$ and an even-degree non-isolated vertex $x$ of $G$, the
\emph{endpoint budget} $\beta(G,x)$ is the minimum cardinality of a path decomposition $D$
of $G$ in which $x$ is an endpoint of at least two paths, i.e.\
$\beta(G,x):=\min\{|D| : D \text{ a path decomposition of } G,\ D(x)\ge2\}$. The
\emph{budget conjecture} $K$ is the statement that $\beta(G,x)\le\ceil{n/2}$ for every
connected $G$ on $n$ vertices and every even-degree non-isolated vertex $x$.
\end{definition}

\begin{theorem}\label{int:thm-Kgallai}
Conjecture~$K$ implies Gallai's conjecture: every connected $n$-vertex graph $G$ has
$\pn(G)\le\ceil{n/2}$.
\end{theorem}

\begin{proof}
If $G$ has an even-degree non-isolated vertex $x$, then $K$ provides a decomposition with
at most $\ceil{n/2}$ paths and we are done. Otherwise $G$ is all-odd, so $n$ is even, and
Theorem~\ref{int:thm-lovasz} gives a decomposition into $p$ paths and $c$ cycles with
$p+c\le n/2$. In a mixed path-cycle decomposition the parity identity of
Lemma~\ref{int:lem-LB1}(iii) still holds (cycles contribute even degree at each vertex), so
every vertex is an endpoint of at least one path and $2p\ge n$; hence $p\ge n/2$, so
$p=n/2$, $c=0$, and the decomposition is pure, giving $\pn(G)\le n/2=\ceil{n/2}$.
\end{proof}

Thus $K$ is a strengthening of Gallai's conjecture, not a consequence of it. The remaining
links of the conditional chain --- from the degree-two layer of $K$ through the odd
semi-clique pair-hosting conjecture to the placement problems --- are assembled descriptively
below and carried at ledger level in Section~\ref{sta:subsec-limitations}; here we restate,
descriptively, the eight-item premise set behind the conditional-chain claim of
Section~\ref{int:subsec-contrib}; each link's verification level, premises and residual gaps
are registered at pointer level in the limitations subsection.

\begin{theorem}[conditional chain; descriptive registration]\label{int:thm-ledger}
There is a chain of implications, each available at paper level or registered at pointer
level with its verification status recorded, leading from the conjunction of the eight-item
premise set listed in this subsection (with the valuation criterion of
Conjecture~\ref{conj:clacand} entering, if at all, only through the bridge premises as a
parallel, redundant-for-implication route) to the conclusion
$\pn(G)\le\ceil{n/2}$ for every connected $G$ on $n$ vertices. The theorem is the
\emph{existence of the conditional route}, not the truth of any premise: every premise is an
open conjecture, an unverified external step, or a paper-open item with machine backing,
none is claimed proved, and no unconditional progress on Gallai's conjecture is claimed.
\end{theorem}

Each item below is an \emph{open} conjecture, an \emph{unverified external step}, or a
\emph{paper-open item with machine backing}; none is claimed proved.
\begin{itemize}[leftmargin=2em,itemsep=0.1em]
\item \emph{(external, flagged unverified against the journal record)} the all-odd
branch of the Lov\'asz decomposition theorem (Theorem~\ref{int:thm-lovasz}) as consumed by
the proof that the budget conjecture implies Gallai's conjecture: if $G$ is all-odd, the
theorem gives $p$ paths and $c$ cycles with $p+c\le n/2$, and the parity identity forces
$p=n/2$, $c=0$.
\item \emph{(external-origin, open)} BP-odd: every connected graph $H$ on an odd
number $2k+1$ of vertices with $|E(H)|\le 2k^2$ satisfies $\pn(H)\le k$ --- the odd-order,
tightness-special case of Question~1.1 of \cite{BP19}.
\item \emph{(internal, open)} the odd-order density conjecture at $k\ge4$: every
tight connected graph on $2k+1$ vertices ($\pn=\ceil{(2k+1)/2}$) is an odd semi-clique;
proved at $k\le3$ by exhaustive census, machine-certified at $k=4$ and supported by directed
search at $k=5$, open at $k\ge6$.
\item SC-PH: every odd semi-clique on $2k+1$ vertices hosts any two specified
non-isolated vertices on distinct paths of a $(k+1)$-path decomposition.
\item \emph{(open)} the surviving layers of the budget conjecture $K$: the
degree-two cut-vertex branch (computational level, awaiting theorization) and the
degree-$\ge4$ general-layer residual (a recorded dead zone of $42$ instances, discharged at
witness level only). \textbf{This premise is open as stated}: the recorded instances have
witness certificates, but no general paper proof exists, and the premise must be quoted as
an open hypothesis wherever the chain is used.
\item \emph{(paper-open, machine-backed)} general solvability of the placement
problem on the $(E)$-side, cyclic domain included: zero counterexamples at $k\le7$ over all
$188$ types (and machine accounts beyond), paper proof open.
\item \emph{(paper-open)} the $(R)$-side on the $3f\ge k$ domain: computational at
$k\le7$ (zero counterexamples), paper proof open.
\item \emph{(bridge, open)} the valuation-line bridge conditions: three registered
open conjectures of the difference-board analysis, quoted here descriptively and carried at
pointer level in the ledger.
\end{itemize}

We record explicitly what the ledger is and is not. The chain above is \emph{not} assembled
into a single master theorem of the form ``premises $\wedge$ criterion $\Rightarrow$ Gallai''
in this paper: the bridge premises (the last three items) are registered at pointer level
only, their full definitions residing in the project archive, so the paper certifies their
provenance rather than their re-adjudicability. What the ledger delivers is checkability:
any future resolution of any single premise propagates along the described chain through
paper-level implications whose verification status is recorded. No premise is claimed to be
within reach of the machine layer except as recorded, and the conditional route must not be
cited as progress towards Gallai's conjecture beyond its registered content.

\subsection{Reproducibility}\label{int:subsec-repro}

The machine pipeline behind every computational claim in this paper --- the exact-value
computations, the decomposition verifiers, the per-instance sweeps of
Section~\ref{sec:computation}, and the Lean development reported there and in
Section~\ref{sec:open} --- will be released together with the paper: the verification scripts, the certificate and
ledger files, the Lean sources with their axiom profiles, and the per-theorem verification
status table. The release makes every machine-checked number in this paper independently
re-runnable and every Lean theorem independently re-compilable.

% end sec-intro.tex

% ============================================================================
% sec-constructions.tex — Constructions: the exact path number of K_{4,n}
% (odd n), the ring-closure machinery, and odd clique chains.
% ===========================================================================
\section{Constructions: \texorpdfstring{$K_{4,n}$}{K(4,n)} and the ring
closure; odd clique chains}\label{sec:constructions}

\subsection{The exact path number of \texorpdfstring{$K_{4,n}$}{K(4,n)} for
odd \texorpdfstring{$n$}{n}}\label{subsec:k4odd}

Let $K_{4,n}$ have bipartition $A=\{a_0,a_1,a_2,a_3\}$ and
$B=\{b_0,\dots,b_{n-1}\}$. For $n$ odd, $n=2k+1$, the four vertices of $A$
have odd degree and the vertices of $B$ have degree $4$, so
$\mathrm{odd}(K_{4,n})=4$ and Lemma~\ref{int:lem-LB1} gives only $\pn\ge2$: the
odd case is genuinely harder than the even one, where the edge bound alone
is tight. Write $n'=n+4=2k+5$ for the total order; Gallai's target is
$\ceil{n'/2}=k+3$, and the best published bound is $k+2$, since
$\EV(K_{4,n})$ is the independent set $B$, a forest, so
Theorem~\ref{int:thm-pyber} applies, and the same bound follows from
\cite{Fan05} and from \cite{LY24}, as $\mathrm{tw}(K_{4,n})=4$. Note that
\cite{ZLH24} does not apply: $K_{4,n}$ has minimum degree $4$ and is not
$3$-degenerate.

For context, Constantinou and Ellinas \cite{CE18} partition the parameter
pairs $(n_1,n_2)$ of $K_{n_1,n_2}$ into four cases and give
minimal-decomposition algorithms for three of them; in their Case~2 (even
$n_1$, $1\le n_2\le n_1-1$) the minimum is $n_1/2$. Their Case~4 (odd
$n_1$, $1\le n_2\le n_1-1$) is left open, with the explicit request that
its ``exact characteristics need to be first understood''. Our main theorem
settles the $n_2=4$ column of Case~4;
Section~\ref{subsec:ring} settles the $n_2=6$ column, proves a general
structure theorem and a ring-closure theorem for all even columns, and
reduces the remaining columns to finitely many base instances.

\subsubsection{Lower bound}

\begin{lemma}\label{lem:concap}
Every path in $K_{4,n}$ uses at most $8$ edges if $n\ge5$, at most $6$ edges
if $n=3$, and at most $2$ edges if $n=1$. Consequently, for odd $n\ge3$,
$\pn(K_{4,n})\ge(n+1)/2$, and $\pn(K_{4,1})\ge2$.
\end{lemma}

\begin{proof}
Paths in a bipartite graph alternate between $A$ and $B$. A path containing
$\alpha$ vertices of $A$ and $\beta$ of $B$ has $e=\alpha+\beta-1$ edges
with $|\alpha-\beta|\le1$. Since $|A|=4$ we have $\alpha\le4$ and
$\beta\le\alpha+1\le5$, so $e\le8$ for $n\ge5$ (where $|B|\ge5$); for
$n=3$, $\beta\le3$ gives $e\le6$. With $m=4n$ edges partitioned into $p$
paths, $8p\ge4n$ gives $p\ge n/2$, i.e.\ $p\ge(n+1)/2$ for odd $n\ge5$; and
$6p\ge12$ gives $p\ge2=(3+1)/2$ for $n=3$. For $n=1$ the graph is a
$4$-star and every path through the centre uses at most two edges, so
$p\ge\ceil{4/2}=2$.
\end{proof}

\subsubsection{Structure of optimal decompositions}

\begin{theorem}\label{thm:conk4struct}
Let $n=2k+1\ge3$ and let $D=\{P_1,\dots,P_{k+1}\}$ be a path decomposition
of $K_{4,n}$ with $k+1$ paths. Then:
\begin{enumerate}[label={\rm(\arabic*)},leftmargin=2em]
\item every path of $D$ contains all four vertices of $A$; each $a\in A$ is
an endpoint of exactly one path of $D$ and an internal vertex of exactly
$k$ paths;
\item the multiset of path lengths (edge counts) is one of
$[7,7,7,7,\,8^{k-3}]$ ($k\ge3$), $[6,7,7,\,8^{k-2}]$ ($k\ge2$), or
$[6,6,\,8^{k-1}]$ ($k\ge1$; for $k=1$ this reads $[6,6]$);
\item a $6$-edge path has vertex type $4A3B$ (both endpoints in $A$), a
$7$-edge path has type $4A4B$ (one endpoint in each class), and an
$8$-edge path has type $4A5B$ (both endpoints in $B$).
\end{enumerate}
\end{theorem}

\begin{proof}
(1) Fix $a\in A$, of degree $2k+1$. Let $r_a$ be the number of paths of $D$
in which $a$ is an endpoint and $s_a$ the number in which $a$ is internal.
Counting the degree of $a$ path by path, $r_a+2s_a=2k+1$, so $r_a$ is odd,
$r_a\ge1$, and $r_a+s_a=(2k+1+r_a)/2\le k+1$ (there are only $k+1$ paths)
forces $r_a\le1$. Hence $r_a=1$ and $s_a=k$. So each $a\in A$ is an
endpoint of exactly one path and internal in exactly $k$ paths, i.e.\ $a$
occurs in all $k+1$ paths.

(2)+(3) Assume first $n\ge5$. By Lemma~\ref{lem:concap} each path has
$e_i\le8$ edges, and $\sum_ie_i=4n=8k+4$, while $(k+1)\cdot8=8k+8$; hence
the deficits $\delta_i=8-e_i$ sum to $4$. A path with $\alpha$ vertices of
$A$ and $\beta$ of $B$, $\alpha+\beta=e+1$ and $|\alpha-\beta|\le1$, has an
endpoint in $A$ exactly when $\alpha=\beta+1$. Enumerating the
possibilities with $\alpha\le4$: an $8$-edge path ($\delta=0$) has type
$(\alpha,\beta)=(4,5)$ and $0$ endpoints in $A$; a $7$-edge path
($\delta=1$) has type $(4,4)$ and exactly $1$; a $6$-edge path
($\delta=2$) has type $(4,3)$ or $(3,4)$, so at most $2$; a $5$-edge path
($\delta=3$) has type $(3,3)$, exactly $1$; a $4$-edge path ($\delta=4$)
has type $(3,2)$ or $(2,3)$, at most $2$. (Larger deficits cannot occur
since $\sum\delta_i=4$.) Let $f(\delta)$ denote these maxima $0,1,2,1,2$.
By (1) the total number of endpoints lying in $A$ is exactly
$\sum_{a}r_a=4$, while it is at most $\sum_if(\delta_i)$; over all
partitions of $\sum\delta_i=4$ the maximum of $\sum_if(\delta_i)$ is $4$,
attained exactly at the profiles $(2,2)$, $(2,1,1)$, $(1,1,1,1)$ (all other
partitions give $\le2$). Hence equality holds throughout: the deficit
multiset is one of the three listed, and each path attains its
$f(\delta_i)$, which for $\delta=2$ forces the type $(4,3)$; this gives
(2) and (3). For $n=3$ the cap is $6$: two paths with $\sum e_i=12$ force
both to be $6$-edge paths of type $(4,3)$ with $2+2=4$ endpoints in $A$,
consistent with all three claims.
\end{proof}

In particular, in any minimal decomposition each $a\in A$ has its $2k+1$
incident edges partitioned into $k$ ``pairs'' (the two $B$-edges flanking
$a$ inside a path) plus one ``single'' (the edge of the path on which $a$
is an endpoint), and the four resulting pair-plus-single systems each
partition $B$. This is the language in which we audit the constructions
below, and it answers, for the column $n_2=4$, the request of \cite{CE18}
for the exact characteristics of Case~4.

\subsubsection{Upper bound: a closed-form construction for large
\texorpdfstring{$n$}{n}}

\begin{proposition}\label{prop:conk4large}
For every $k\ge6$, writing $n=2k+1$, the graph $K_{4,n}$ has a path
decomposition into $k+1$ paths.
\end{proposition}

\begin{proof}
Split $B$ into a \emph{core} $S=\{s_0,\dots,s_6\}$ of size $7$ and a
\emph{rim} $R=\{r_0,\dots,r_{m-1}\}$ of size $m=2k-6$ (even, and $m\ge6$);
$|S|+|R|=2k+1=n$. We use four \emph{gadget paths} of $7$ edges each, on
$S$ only:
\begin{align*}
G_1&=a_0\,s_0\,a_1\,s_1\,a_2\,s_2\,a_3\,s_3, &
G_2&=a_1\,s_2\,a_0\,s_3\,a_2\,s_0\,a_3\,s_1,\\
G_3&=a_2\,s_4\,a_1\,s_5\,a_3\,s_6\,a_0\,s_1, &
G_4&=a_3\,s_4\,a_0\,s_5\,a_2\,s_6\,a_1\,s_3,
\end{align*}
and $k-3$ \emph{cyclic paths} of $8$ edges each, on $R$ only, with indices
mod $m$:
\[
Q_t \;=\; b_{r(2t)}\,a_0\,b_{r(2t+1)}\,a_1\,b_{r(2t+2)}\,a_2\,b_{r(2t+3)}\,a_3\,b_{r(2t+4)},
\qquad t=0,\dots,k-4 .
\]
\emph{The paths are legal.} Each $G_i$ is an alternating sequence of $4$
distinct $A$-vertices and $4$ distinct elements of $S$, hence a simple path
of $7$ edges. Each $Q_t$ uses five consecutive indices $2t,\dots,2t+4$ mod
$m$; since $m=2k-6\ge6>5$ these are distinct, so $Q_t$ is a simple path of
$8$ edges.

\emph{The gadget paths partition $A\times S$.}
For each $a\in A$, list the pairs of $S$-edges flanking $a$ (or the single
edge, when $a$ is an endpoint) in each of the four $G_i$:
\begin{align*}
a_0&:\ \{s_0\};\ \{s_2,s_3\};\ \{s_6,s_1\};\ \{s_4,s_5\} &
&\text{partition } S;\\
a_1&:\ \{s_0,s_1\};\ \{s_2\};\ \{s_4,s_5\};\ \{s_6,s_3\} &&
\text{partition } S;\\
a_2&:\ \{s_1,s_2\};\ \{s_3,s_0\};\ \{s_4\};\ \{s_5,s_6\} &&
\text{partition } S;\\
a_3&:\ \{s_2,s_3\};\ \{s_0,s_1\};\ \{s_5,s_6\};\ \{s_4\} &&
\text{partition } S.
\end{align*}
Each pair $\{s',s''\}$ corresponds to the two edges $as'$, $as''$ used by
that path at the internal position of $a$, and each singleton to the
endpoint edge. Since the four groups partition $S$ for every $a$, every
edge of $A\times S$ is used by exactly one $G_i$.

\emph{The cyclic paths partition $A\times R$.}
For fixed $a_c$, $c\in\{0,1,2,3\}$, the path $Q_t$ uses at $a_c$ the pair
of edges $\{a_cr(2t+c),\,a_cr(2t+c+1)\}$ (indices mod $m$). As $t$ runs
over $0,\dots,k-4$, the even residues $2t$ cover all $k-3$ even elements
of $\mathbb{Z}_m$ and the odd residues $2t+1$ cover all $k-3$ odd
elements; shifting by $c$, each element of $\mathbb{Z}_m$ occurs exactly
once, in exactly one pair, across the $k-3$ paths. Hence the $m=2(k-3)$
edges $a_cR$ are partitioned, for each $c$ independently.

\emph{Count.} $4+(k-3)=k+1$ paths, as claimed.
\end{proof}

\subsubsection{Small cases}

\begin{proposition}\label{prop:conk4small}
Explicit decompositions into $k+1$ paths exist for $n=3,5,7,9,11$
($k=1,\dots,5$). The path lists below are machine-verified to partition
$E(K_{4,n})$; write $b_j$ for the $j$-th vertex of $B$.
\begingroup\sloppy
\begin{itemize}[leftmargin=1.5em]
\item[{$k=1$ ($n=3$; lengths $[6,6]$):}]
$a_0b_0a_1b_1a_2b_2a_3$;\quad
$a_1b_2a_0b_1a_3b_0a_2$.
\item[{$k=2$ ($n=5$; $[6,7,7]$):}]
$a_0b_0a_1b_1a_2b_2a_3b_3$;\quad
$a_1b_2a_0b_4a_3b_0a_2b_3$;\quad
$a_2b_4a_1b_3a_0b_1a_3$.
\item[{$k=3$ ($n=7$; $[7,7,7,7]$):}]
$a_0b_0a_1b_1a_2b_2a_3b_3$;\quad
$a_1b_2a_0b_3a_2b_0a_3b_1$;\quad
$a_2b_4a_1b_5a_3b_6a_0b_1$;\quad
$a_3b_4a_0b_5a_2b_6a_1b_3$.
\item[{$k=4$ ($n=9$; $[7,7,7,7,8]$):}]
$b_0a_0b_1a_1b_2a_2b_3a_3b_4$;\quad
$a_0b_3a_1b_0a_2b_1a_3b_2$;\quad
$a_1b_4a_0b_5a_2b_6a_3b_0$;\quad
$a_2b_7a_1b_5a_3b_8a_0b_2$;\quad
$a_3b_7a_0b_6a_1b_8a_2b_4$.
\item[{$k=5$ ($n=11$; $[7,7,7,7,8,8]$):}]
$b_0a_0b_1a_1b_2a_2b_3a_3b_4$;\quad
$b_2a_0b_3a_1b_0a_2b_1a_3b_5$;\quad
$a_0b_4a_1b_5a_2b_6a_3b_0$;\quad
$a_1b_6a_0b_7a_2b_8a_3b_2$;\quad
$a_2b_9a_1b_7a_3b_{10}a_0b_5$;\quad
$a_3b_9a_0b_8a_1b_{10}a_2b_4$.
\end{itemize}
\endgroup
The reason $k=4,5$ are listed separately is that the rim $R$ of
Proposition~\ref{prop:conk4large} is too small there; the single $8$-edge
paths absorb the shortfall. All lists conform to the length profiles of
Theorem~\ref{thm:conk4struct}.
\end{proposition}

\subsubsection{The main theorem and its corollaries}

\begin{theorem}\label{thm:conk4odd}
For every odd $n\ge3$, $\pn(K_{4,n})=(n+1)/2=\ceil{n/2}$. Also
$\pn(K_{4,1})=2$.
\end{theorem}

\begin{proof}
Combine the lower bound of Lemma~\ref{lem:concap} with the upper bounds of
Proposition~\ref{prop:conk4large} ($k\ge6$) and
Proposition~\ref{prop:conk4small} ($k\le5$). For $n=1$, the lower bound of
Lemma~\ref{lem:concap} matches the two-path decomposition
$a_0b_0a_1 \mid a_2b_0a_3$.
\end{proof}

\begin{corollary}[Gallai comparison]\label{cor:conk4gallai}
For odd $n\ge3$, $K_{4,n}$ has $n+4=2k+5$ vertices and Gallai's conjecture
demands $\pn\le k+3$; in fact $\pn=k+1$, with slack $2$. There is no
counterexample here.
\end{corollary}

\begin{proof}
Immediate from Theorem~\ref{thm:conk4odd}: the value $(n+1)/2=k+1$ is two
fewer than the $k+3$ paths the conjecture allows.
\end{proof}

\begin{corollary}[comparison with published bounds]\label{cor:conk4compare}
The best published bound for this family was $k+2$: it follows from
Pyber's theorem \cite{Pyb96} (equally from \cite{Fan05}, and from
\cite{LY24} since $\mathrm{tw}(K_{4,n})=4$). Theorem~\ref{thm:conk4odd}
improves it by one and shows that the $\floor{n'/2}$ bound of the
treewidth-$4$ theorem \cite{LY24} is not tight on this family.
\end{corollary}

\begin{proof}
Immediate: Theorem~\ref{int:thm-pyber} (equally \cite{Fan05,LY24}) gives $k+2$
for the family, while Theorem~\ref{thm:conk4odd} attains $k+1$; the
$\floor{n'/2}=k+3$ bound of \cite{LY24} exceeds the exact value by $2$.
\end{proof}

\begin{remark}[unified formula for the whole family]\label{rem:conk4all}
Combining Theorem~\ref{thm:conk4odd} with the known even-case values
(Case~2 of \cite{CE18} for even $n\ge6$, and the small values
$\pn(K_{4,2})=2$, $\pn(K_{4,4})=3$, which we also re-verified by machine),
the path number of $K_{4,n}$ is
\[
\pn(K_{4,n})=\ceil{n/2}+[n\le 4]\quad\text{for all } n\ge1,
\]
i.e.\ $\ceil{n/2}$ except for the three small values
$\pn(K_{4,1})=\pn(K_{4,2})=2$ and $\pn(K_{4,4})=3$.
\end{remark}

\subsection{The remaining columns of Case 4}\label{subsec:columns}

For odd $n_1$, Case~4 of \cite{CE18} is the range $1\le n_2\le n_1-1$. The
state of affairs by column:
\begin{center}\small
\begin{tabular}{c l l}
\hline
$n_2$ & $\pn(K_{n_1,n_2})$, $n_1$ odd & status \\
\hline
$1$ & $\ceil{n_1/2}$ & star decomposition, classical\\
$2$ & $(n_1+1)/2$ & counting cap (as in Lemma~\ref{lem:concap}) $+$ \cite{Pyb96}\\
$3$ & $(n_1+3)/2$ & Lemma~\ref{int:lem-LB1} $+$ \cite{Pyb96}\\
$4$ & $(n_1+1)/2=\ceil{n_1/2}$ & Theorem~\ref{thm:conk4odd} (this paper)\\
$6$ & $(n_1+1)/2$ for $n_1\ge7$ & Corollary~\ref{cor:concol6}
(Sec.~\ref{subsec:ring})\\
even $\ge8$ & $\ceil{n_1/2}\ \le\ \pn\ \le\ \floor{(n_1+n_2)/2}$ & bases
settled for $n_2=8,10$; rest open (Sec.~\ref{subsec:ring})\\
\hline
\end{tabular}
\end{center}

\begin{conjecture}\label{conj:con27A}
For all odd $n_1$ and even $n_2$ with $2\le n_2<n_1$:
$\pn(K_{n_1,n_2})=\ceil{n_1/2}$.
\end{conjecture}

The evidence: the universal counting cap gives $\pn\ge\ceil{n_1/2}$ in the
whole even range (Theorem~\ref{thm:concapgen} below; each path uses at
most $2n_2$ edges, and $m=n_1n_2$); the even analogue (Case~2 of
\cite{CE18}) has exactly the value $n_1/2=\ceil{n_1/2}$; and
$\pn(K_{5,2})=\pn(K_{5,4})=3=\ceil{5/2}$ by exact computation. The
structural tools developed in the next subsection reduce the conjecture to
finitely many instances per column, of which all but five named ones are
now settled.

\subsection{Beyond the fourth column: ring closure and the state of
Conjecture~\ref{conj:con27A}}\label{subsec:ring}

In this subsection $n_1=2k+1$ is odd and $n_2=2h$ is even with
$k\ge h\ge1$ --- exactly the range of Conjecture~\ref{conj:con27A}. We
write $A=\{a_0,\dots,a_{2h-1}\}$ for the side of size $n_2$ (each of its
vertices has degree $2k+1$, odd) and $B=\{b_0,\dots,b_{2k}\}$ for the side
of size $n_1$ (each of degree $2h$, even); note that this swaps the roles
of $A$ and $B$ relative to the earlier parts of this section, where the
$4$-vertex side was called $A$. Conjecture~\ref{conj:con27A} predicts
$\pn(K_{2k+1,2h})=k+1$.

\begin{theorem}[general counting cap]\label{thm:concapgen}
For all $k\ge h\ge1$, $\pn(K_{2k+1,2h})\ge k+1$.
\end{theorem}

\begin{proof}
Paths in a bipartite graph alternate between the two sides. A path
containing $\alpha$ vertices of $A$ and $\beta$ vertices of $B$ has
$\alpha+\beta-1$ edges with $|\alpha-\beta|\le1$; since $\alpha\le2h$ we
get $\beta\le\alpha+1\le2h+1$, so every path has at most
$2h+(2h+1)-1=4h$ edges. If $p$ paths partition the $m=(2k+1)\cdot2h$
edges, then $4hp\ge2h(2k+1)$, i.e.\ $p\ge\ceil{(2k+1)/2}=k+1$.
\end{proof}

\begin{theorem}[structure of optimal decompositions, general $h$]
\label{thm:constructgen}
Let $k\ge h\ge1$ and let $D$ be a path decomposition of $K_{2k+1,2h}$ with
exactly $k+1$ paths. Then:
\begin{enumerate}[label={\rm(\arabic*)},leftmargin=2em]
\item every path of $D$ contains all $2h$ vertices of $A$; each $a\in A$
is an endpoint of exactly one path of $D$ and an internal vertex of
exactly $k$ paths;
\item every path has $e\in\{4h-2,4h-1,4h\}$ edges, and writing
$\delta=4h-e\in\{2,1,0\}$, $\delta$ equals the number of endpoints of the
path lying in $A$;
\item with $x,y,z$ the numbers of paths of deficit $\delta=2,1,0$,
respectively, we have $x+y+z=k+1$ and $2x+y=2h$; the total number of
path-ends lying in $B$ equals $2z+y=2(k-h+1)$;
\item every $b\in B$ is an endpoint of an even number of paths of $D$.
\end{enumerate}
\end{theorem}

\begin{proof}
(1) Fix $a\in A$, of degree $2k+1$; let $r_a$ be the number of paths of $D$
ending at $a$ and $s_a$ the number having $a$ as an internal vertex.
Counting the degree of $a$ path by path gives $r_a+2s_a=2k+1$, so $r_a$ is
odd and $r_a\ge1$; since $a$ occurs in $r_a+s_a$ of the $k+1$ paths,
$r_a+s_a\le k+1$, i.e.\ $(2k+1+r_a)/2\le k+1$, forcing $r_a=1$ and
$s_a=k$. So $a$ occurs in all $k+1$ paths, and $\sum_{a\in A}r_a=2h$.

(2)+(3) Every path contains all $2h$ vertices of $A$ and alternates, so if
it contains $\beta$ vertices of $B$ then $|\beta-2h|\le1$; accordingly its
two endpoints lie both in $A$ (when $\beta=2h-1$), one in each class
($\beta=2h$), or both in $B$ ($\beta=2h+1$), and
$e=2h+\beta-1\in\{4h-2,4h-1,4h\}$ with $\delta=4h-e$ equal to the number
of $A$-endpoints. Summing over $D$, $\sum e=|E|=2h(2k+1)$, so
$\sum\delta=4h(k+1)-2h(2k+1)=2h$. Counting $A$-endpoints two ways gives
$2x+y=\sum_{a\in A}r_a=2h$; combined with $x+y+z=k+1$ this yields
$z=k+1+x-2h$ and $2z+y=2(k+1)-2h=2(k-h+1)$.

(4) For $b\in B$, of degree $2h$, the count of (1) reads $r_b+2s_b=2h$, so
$r_b$ is even.
\end{proof}

\begin{theorem}[ring closure]\label{thm:conringclosure}
Let $h\ge1$ and suppose that for some $k_0\ge h$ the graph
$K_{2k_0+1,2h}$ has a path decomposition $D_0$ into $k_0+1$ paths. Then
for every $k\ge k_0+h+1$ the graph $K_{2k+1,2h}$ has a path decomposition
into $k+1$ paths; consequently $\pn(K_{2k+1,2h})=k+1$ for all
$k\ge k_0+h+1$, by Theorem~\ref{thm:concapgen}.
\end{theorem}

\begin{proof}
Write $m=2(k-k_0)$, an even number satisfying $m\ge2h+2$, in particular
$2h+1<m$. Split the $B$-side of $K_{2k+1,2h}$ into
$B_0=\{b_0,\dots,b_{2k_0}\}$, of size $2k_0+1$, which carries $D_0$, and
$m$ \emph{fresh} ring vertices $b_{2k_0+1},\dots,b_{2k}$; write
$r(i)=2k_0+1+(i\bmod m)$, so that the ring vertices are
$\{b_{r(i)} : i\in\mathbb{Z}_m\}$. The new decomposition consists of the
paths of $D_0$ unchanged (they use only edges into $B_0$) and
$k-k_0=m/2$ \emph{ring paths}, one for each $t\in\{0,\dots,m/2-1\}$:
\[
Q_t\;=\;b_{r(2t)}\,a_0\,b_{r(2t+1)}\,a_1\,b_{r(2t+2)}\,a_2\,\cdots\,a_{2h-1}\,b_{r(2t+2h)}
\qquad(\text{all ring indices mod } m).
\]
\emph{The paths are legal.} $Q_t$ uses the $2h+1$ consecutive ring indices
$2t,\dots,2t+2h$; since $2h+1<m$ these are distinct, so $Q_t$ is a simple
path with $4h$ edges, all of them edges between $A$ and the ring.

\emph{The ring paths partition $A\times R$, where $R$ is the ring.} Fix
$a_c$, $c\in\{0,\dots,2h-1\}$. The path $Q_t$ uses at $a_c$ the pair of
edges $\{a_cb_{r(2t+c)},\,a_cb_{r(2t+c+1)}\}$. As $t$ runs over
$0,\dots,m/2-1$, the index $2t+c$ runs over all elements of
$\mathbb{Z}_m$ congruent to $c$ modulo $2$, and $2t+c+1$ over the
complementary class (the two parity classes are well defined because $m$
is even); hence the $m/2$ pairs are pairwise disjoint with union
$\mathbb{Z}_m$, i.e.\ each of the $m$ edges from $a_c$ to the ring is used
exactly once, and this for each $c$ independently.

\emph{Ring degrees.} A ring vertex $b_{r(j)}$ occurs at $a_c$ exactly when
$j\in\{2t+c,\,2t+c+1\}$ for some $t$, and for each fixed $c$ there is
exactly one such $t$; so each ring vertex occurs once at each $a_c$, i.e.\
exactly $2h$ times in all --- its degree --- each time through a different
edge.

\emph{Count.} The edges used by $D_0$ lie in $A\times B_0$ and those of
the ring paths in $A\times R$, disjoint sets; $D_0$ partitions
$A\times B_0$, of size $2h(2k_0+1)$, by hypothesis, and the ring paths
partition $A\times R$, of size $2hm$, by the paragraph above. The total
number of paths is $(k_0+1)+(k-k_0)=k+1$, so this is a decomposition of
$K_{2k+1,2h}$ into $k+1$ paths, and Theorem~\ref{thm:concapgen} makes the
value exact.
\end{proof}

\begin{corollary}[the $n_2=6$ column]\label{cor:concol6}
For every odd $n_1\ge7$, $\pn(K_{n_1,6})=(n_1+1)/2$; equivalently,
$\pn(K_{2k+1,6})=k+1$ for all $k\ge3$.
\end{corollary}

\begin{proof}
By Theorem~\ref{thm:concapgen} it suffices to exhibit $(k+1)$-path
decompositions for the four bases $k=3,4,5,6$. For $k=3$ ($K_{7,6}$, with
$A=\{a_0,\dots,a_5\}$, $B=\{b_0,\dots,b_6\}$) the following four paths, of
deficits $[2,2,1,1]$ in the sense of Theorem~\ref{thm:constructgen}(3),
are machine-verified to partition the $42$ edges (Appendix~\ref{com:app-comp}):
\begin{center}\small
$a_0b_0a_1b_1a_2b_2a_3b_3a_4b_4a_5$;\quad
$a_1b_2a_0b_3a_2b_0a_3b_5a_5b_6a_4$;\\
$a_2b_4a_0b_5a_1b_6a_3b_1a_4b_0a_5b_2$;\quad
$a_3b_4a_1b_3a_5b_1a_0b_6a_2b_5a_4b_2$.
\end{center}
For $k=4,5,6$ ($K_{9,6}$, $K_{11,6}$, $K_{13,6}$) machine-verified
$(k+1)$-path witnesses are archived, each with the length profile forced
by the ledger of Theorem~\ref{thm:constructgen}(3) ($\sum\delta=6$), and
each checked by two independent implementations. The base $k_0=3$, fed
into Theorem~\ref{thm:conringclosure}, covers every $k\ge7$; the pipeline
additionally re-derived the ring construction explicitly for $k=7,\dots,31$
and obtained overlapping closure chains from the bases $k_0=4,5$
(Appendix~\ref{com:app-comp}).
\end{proof}

\paragraph{The state of Conjecture~\ref{conj:con27A}.}
Theorem~\ref{thm:conringclosure} upgrades Conjecture~\ref{conj:con27A}
from a statement about infinitely many graphs to a
\emph{column-by-column decision procedure}: in the column $n_2=2h$,
Theorem~\ref{thm:concapgen} shows that it suffices to exhibit, for each of
the $h+1$ \emph{base} values $k\in[h,2h]$, one explicit $(k+1)$-path
decomposition of $K_{2k+1,2h}$; the base $k_0=h$ then closes the whole
tail $k\ge2h+1$ for free, and every further base lowers the closure line.
Applying the procedure as far as the current machine pipeline reaches
gives:

\begin{center}\small
\begin{tabular}{c p{6.4cm} p{3.2cm}}
\hline
$n_2$ & status of $\pn(K_{n_1,n_2})$, $n_1$ odd & remaining bases \\
\hline
$2$ & $(n_1+1)/2$; proved (counting cap $+$ \cite{Pyb96}; the pipeline also
holds a constructive closure chain from the base $k_0=1$) & none \\
$4$ & $(n_1+1)/2$; Theorem~\ref{thm:conk4odd} & none \\
$6$ & $(n_1+1)/2$ for $n_1\ge7$; Corollary~\ref{cor:concol6} & none \\
$8$ & bases $k=4,5,6$ exact: $\pn(K_{9,8})=5$, $\pn(K_{11,8})=6$,
$\pn(K_{13,8})=7$; tail $k\ge9$ closed by ring closure & $k=7,8$ \\
$10$ & bases $k=5,6,7$ exact: $\pn(K_{11,10})=6$, $\pn(K_{13,10})=7$,
$\pn(K_{15,10})=8$; tail $k\ge11$ closed & $k=8,9,10$ \\
even $\ge12$ & open; reduced to $h+1$ bases per column & $k\in[h,2h]$ \\
\hline
\end{tabular}
\end{center}

In total this adds ten exact values beyond Theorem~\ref{thm:conk4odd}
(four in the column $n_2=6$, three in the column $n_2=8$, three in the
column $n_2=10$), every one attained by an explicitly archived and
independently verified witness, and two constructively closed tails
($k\ge9$ for $n_2=8$ from the base $K_{9,8}$, $k\ge11$ for $n_2=10$ from
the base $K_{11,10}$). No instance in the range is known to violate
Conjecture~\ref{conj:con27A}.

What remains open are exactly the five middle-window bases $K_{15,8}$,
$K_{17,8}$, $K_{17,10}$, $K_{19,10}$, $K_{21,10}$. On these the
deterministic search pipeline has burnt more than $14$ million search
nodes (for the $n_2=8$ pair) and about $0.9$ million nodes (for the
$n_2=10$ triple), over every ledger-feasible deficit multiset of
Theorem~\ref{thm:constructgen}(3), without a hit and without exhausting a
single schedule; these negatives are truncations, not exclusions. The
following observation records the data-level regularities; it is a
\emph{computational observation}, not a theorem.

\begin{observation}[middle-window cost wall; computational]
\label{obs:concostwall}
The following regularities of the deterministic, fixed-seed search
pipeline are data, not proved statements.
(1)~\emph{Hit-cost trajectory.} Within a column, the first hit at level
$k$ is cheap for $k\le 2h-3$ --- $5,7,43$ nodes at $k=4,5,6$ in the column
$n_2=8$, and $6,13,182$ nodes at $k=5,6,7$ in the column $n_2=10$ --- and
then jumps off a cliff inside the \emph{middle window} $k\in[2h-2,2h]$: at
$k=2h-1$ in the column $n_2=8$ (i.e.\ $K_{15,8}$) the strongest schedule
burnt more than $4.8\times10^6$ nodes without a hit, and at $k=2h-2$ in
the column $n_2=10$ (i.e.\ $K_{17,10}$) more than $1.19\times10^5$ nodes;
no schedule was run to exhaustion. The two columns share the same cost
anatomy --- cheap head bases, a free tail by
Theorem~\ref{thm:conringclosure}, a hard middle window --- with the cliff
at $k=2h-1$ for $n_2=8$ and, on current evidence one notch earlier, at
$k=2h-2$ for $n_2=10$.
(2)~\emph{Deficit-multiset signal.} Every hit at $k\le 2h-3$ lies in a
$z=0$ multiset of Theorem~\ref{thm:constructgen}(3), at the first schedule
tried; the unique $z=0$ multisets of the window --- $(x,y,z)=(0,8,0)$ at
$K_{15,8}$ and $(1,8,0)$ at $K_{17,10}$ --- are precisely the types that
resisted the largest search effort, which makes them the strongest current
candidates for non-realizable ledgers in the finitary question ``which
ledger-feasible multisets $(x,y,z)$ occur in some $(k+1)$-path
decomposition?''. The only proved non-realizable multiset so far is
$(x,y,z)=(3,0,1)$, i.e.\ deficit list $[2,2,2,0]$, at $K_{7,6}$: a single
$BB$-type path has two distinct $B$-endpoints $u\ne v$, each forced to be
an endpoint of an even, hence at least $2$, number of paths
(Theorem~\ref{thm:constructgen}(4)), contradicting $2z+y=2$.
\end{observation}

\subsection{Odd clique chains: refinements beyond block graphs}
\label{subsec:chains}

A \emph{block graph} is a graph whose blocks are complete graphs. Chen and
Wu \cite{CW25} proved Gallai's conjecture for all block graphs. Here we
study \emph{odd clique chains}: block graphs whose blocks are
$B_i\cong K_{2j_i+1}$ with $j_i$ odd (equivalently, blocks of order
$3 \bmod 4$), and whose block tree is a path (every block has at most two
cut vertices). We write $b$ for the number of blocks and
$n=2\sum_ij_i+1$ for the order; note $n$ is odd and
$\floor{n/2}=\sum_ij_i$. Blocks that are $K_3$ ($j_i=1$) are allowed.

The engine is the following statement about complete graphs of odd order
minus a Hamilton path.

\begin{theorem}\label{thm:conL0}
Let $j\ge1$ be odd and let $P_0$ be any Hamilton path of $K_{2j+1}$. Then
$K_{2j+1}-P_0$ has a path decomposition into $j$ paths, each with exactly
$2j-1$ edges. Consequently $\pn(K_{2j+1}-P_0)=j$, the value of
Lemma~\ref{int:lem-LB2}.
\end{theorem}

\begin{proof}
For $j=1$ the graph $K_3-P_0$ is a single edge, a decomposition into
$j=1$ paths, and the value $j$ matches Lemma~\ref{int:lem-LB2}; so assume
$j\ge3$ odd. The automorphism group of $K_{2j+1}$ is transitive on
Hamilton paths, so it suffices to exhibit one Hamilton path $P^*$ such
that $K_{2j+1}-P^*$ decomposes into $j$ paths; a general $P_0$ is then the
image of $P^*$ under a permutation of vertices, which transports the
decomposition.

Write $V=\{\infty\}\cup\mathbb{Z}_{2j}$, $\beta=(j+1)/2$. Let
$C_0,\dots,C_{j-1}$ be the Walecki cycles of Lemma~\ref{int:lem-walecki}. For
a vertex $v=k+t$ of $\mathbb{Z}_{2j}$, its two neighbours inside $C_k$ are
read off from the sequence: they are $\{\infty,k+1\}$ if $t=0$;
$\{k+j+1,\infty\}$ if $t=j$; and $\{k-t+1,\,k-t\}$ otherwise. The two
\emph{corners} at a designated \emph{centre} $c$ of $C_k$ are the two
edges of $C_k$ incident to $c$.

We designate one centre in each cycle and one arm-pair $\{y_i,y_i+1\}$ per
corner, as follows (all arithmetic mod $2j$):
\begin{enumerate}[label=\arabic*.,leftmargin=2em]
\item in $C_{j-1}$: centre $c_1=2j-1$ (a vertex of type $t=j$ there),
arm-pair $\{0,\infty\}$ via the reading above: the neighbours of $2j-1$ in
$C_{j-1}$ are $\{2j,\infty\}=\{0,\infty\}$;
\item in $C_{\beta-1}$: centre $c_2=\beta-1$ (type $t=0$), arm
$\{\infty,\beta\}$;
\item for $i=3,\dots,j$: in $C_{k_i}$, where $k_i=i-3$ if $3\le i\le\beta+1$
and $k_i=i-2$ if $\beta+2\le i\le j$, take centre $c_i=2k_i-y_i$ with arm
$\{y_i,y_i+1\}$, $y_i=\beta+i-3$.
\end{enumerate}
\emph{Corner legality.} Cases 1 and 2 are immediate from the neighbour
reading. For case 3, write $c_i=k_i+t_i$ with
$t_i=k_i-y_i\equiv-\beta\equiv(3j-1)/2$ for $i\le\beta+1$ and
$t_i\equiv 1-\beta\equiv(3j+1)/2$ for $i\ge\beta+2$. Both values lie in
$\mathbb{Z}_{2j}\setminus\{0,j\}$ (they equal $0$ never, and equal $j$
only in the void case $j=1$, $i\ge\beta+2=3$), so the neighbour reading
gives $N_{C_{k_i}}(c_i)=\{k_i-t_i+1,\,k_i-t_i\}=\{y_i+1,\,y_i\}$, as
required.

\emph{The prescribed Hamilton path.} Let
\[
P^* \;=\; 0,\;2j-1,\;\infty,\;\beta-1,\;\beta,\;c_3,\;\beta+1,\;c_4,\;\beta+2,\;\dots,\;
c_i,\;\beta+i-2,\;c_{i+1},\;\dots,\;c_j,\;\beta+j-2 .
\]
The edge set of $P^*$ is the union of the $j$ corner arm-pairs (adjacent
corners share the arm vertex $a_i=\beta+i-2$), i.e.\ $2j$ edges. Its
vertex set is
$\{0,\infty\}\cup\{\beta,\dots,\beta+j-2\}\cup\{c_1,\dots,c_j\}$: the arm
interval $[\beta,\beta+j-2]=[(j+1)/2,(3j-3)/2]$ is disjoint from the
centre set
$\{1,\dots,(j-3)/2\}\sqcup\{(j-1)/2\}\sqcup\{(3j-1)/2,\dots,2j-2\}\sqcup\{2j-1\}$,
every centre interval below $(j+1)/2$ resp.\ above $(3j-3)/2$ being
separated from the arm interval (for $j=3$ the first centre interval is
empty), and all $2j+1$ listed values are distinct and cover $V$. Hence
$P^*$ is a Hamilton path with $2j$ edges.

\emph{The partition.}
The designated centre of $C_{j-1}$ is $c_1$, of $C_{\beta-1}$ is $c_2$,
and of $C_{k_i}$ is $c_i$ for $i\ge3$; the index map $i\mapsto k_i$
satisfies
$\{j-1,\beta-1\}\cup\{0,\dots,\beta-2\}\cup\{\beta,\dots,j-2\}=\{0,\dots,j-1\}$
bijectively, so each cycle loses exactly one corner and the corners are
edge-disjoint. By Lemma~\ref{int:lem-walecki},
\[
E(K_{2j+1})\;=\;E(P^*)\;\sqcup\bigsqcup_{k=0}^{j-1}\bigl(E(C_k)\setminus\{\text{the two corner
edges at } c_{k}\}\bigr),
\]
and each $C_k$ minus the two edges at its centre is a path on the
remaining $2j$ vertices with $2j-1$ edges. This gives $j$ paths
decomposing $K_{2j+1}-P^*$. Finally
$\pn(K_{2j+1}-P_0)\ge\ceil{j(2j-1)/2j}=j$ by Lemma~\ref{int:lem-LB2}, so the
value is exactly $j$.
\end{proof}

\begin{remark}
The construction is a Walecki-type zigzag, and the base decomposition
objects are classical \cite{LP92,Alspach08}. The statement for an
\emph{arbitrarily prescribed} Hamilton path, as needed to chain blocks of
order $3 \bmod 4$ (the $1 \bmod 4$ case is the even-$j$ residual, left
open below and in Section~\ref{sec:open}), appears to be new: the
Oberwolfach problem and its variants concern $2$-regular factors and do
not cover path factors, and we could not find this statement in the
Hamilton decomposition literature. The even-$j$ residual is recorded as an
open problem in Section~\ref{sec:open}. The formula was verified by an
independent computer check for all odd $j\le121$; this anchor is a
certificate, not part of the proof.
\end{remark}

The chaining argument needs the following accounting of the decomposition
produced by Theorem~\ref{thm:conL0}.

\begin{lemma}[slot accounting]\label{lem:conslots}
Let $D=\{Q_1,\dots,Q_j\}$ be a decomposition of $K_{2j+1}-P_0$ into $j$
paths of $2j-1$ edges each (as in Theorem~\ref{thm:conL0}), where the ends
of $P_0$ are $w,w'$. Then:
\begin{enumerate}[label={\rm(\roman*)},leftmargin=2em]
\item each $Q_i$ misses exactly one vertex, and the $j$ missed vertices
are distinct;
\item the missed vertices all lie among the internal vertices of $P_0$ (the
even-degree vertices of $K_{2j+1}-P_0$), and $w,w'$ are never missed;
\item each of $w,w'$ is an endpoint of exactly one path of $D$; every
internal vertex $v$ of $P_0$ is either missed (endpoint of no path) or an
endpoint of exactly two paths; and the total number of path-ends is $2j$.
\end{enumerate}
\end{lemma}

\begin{proof}
Each $Q_i$ has $2j-1$ edges, hence spans $2j$ of the $2j+1$ vertices and
misses exactly one. For $v\in V$, let $k_v$ be the number of paths missing
$v$ and $s_v$ its total endpoint count; let $a_v$ be the number of paths
containing $v$ internally, so that the paths containing $v$ are exactly
$a_v+s_v$ and $j=k_v+a_v+s_v$. Lemma~\ref{int:lem-LB1}(iii) gives
$s_v\equiv\deg(v)\pmod 2$, where $\deg_{K_{2j+1}-P_0}(v)=2j-1$ for
$v\in\{w,w'\}$ and $2j-2$ otherwise; and $2\,a_v+s_v=\deg(v)$. If $v$ is
internal to $P_0$ and $k_v=0$: $2a_v+s_v=2j-2$ and $a_v+s_v=j$ give
$s_v=2$; if $k_v=1$: $a_v+s_v=j-1$ gives $s_v=0$. If $v=w$ and $k_v=1$:
$2(j-1-s_v)+s_v=2j-1$ is absurd; so $k_v=0$, $a_v+s_v=j$,
$2a_v+s_v=2j-1$ give $s_v=1$. This proves (ii) and (iii). For (i): each
$Q_i$ misses exactly one vertex, so it suffices that no vertex is missed
twice. But $a_v+s_v=j-k_v$ and $2a_v+s_v=\deg(v)$ give
$s_v=2(j-k_v)-\deg(v)$; for $k_v\ge2$ this is
$s_v\le 2(j-2)-\deg(v)\le 2j-4-(2j-2)=-2<0$ (and $=-3$ at $w,w'$),
impossible. Hence $k_v\le1$ for every $v$, and the $j$ missed vertices are
distinct. The total $2j$ of path-ends is automatic.
\end{proof}

We also need the elementary splicing operation. If $P$ and $P'$ are paths
with $V(P)\cap V(P')=\{w\}$ and $w$ is an endpoint of both, then
$P\cdot P'$ (the concatenation) is a path, and replacing $P,P'$ by
$P\cdot P'$ reduces the size of a decomposition by one, with all other
path-end counts unchanged except that $w$ loses two endpoints. We call
this a \emph{splice at $w$}.

\begin{theorem}\label{thm:conchain}
Let $G$ be a connected graph whose blocks are $B_1,\dots,B_b$ with
$B_i\cong K_{2j_i+1}$, all $j_i$ odd, and whose block tree is a path. Let
$n=2\sum_ij_i+1$. Then:
\begin{enumerate}[label={\rm(\roman*)},leftmargin=2em]
\item if $b=1$, $\pn(G)=\ceil{n/2}=\sum_ij_i+1$;
\item if $b\ge2$, $\pn(G)\le\floor{n/2}=\sum_ij_i$;
\item if $b\ge4$, $\pn(G)\le\floor{n/2}-1=\sum_ij_i-1$.
\end{enumerate}
\end{theorem}

\begin{proof}
(1) For $b=1$, $G=K_{2j+1}$: Lemma~\ref{int:lem-LB2} gives
$\pn\ge\ceil{j(2j+1)/2j}=j+1$; conversely a Hamilton path $P$ together
with the $j$-path decomposition of $K_{2j+1}-P$ from Theorem~\ref{thm:conL0}
gives $\pn\le j+1$.

Assume $b\ge2$ and order the blocks along the block tree as
$B_1,\dots,B_b$, with $B_i\cap B_{i+1}=\{w_i\}$; the $w_i$ are distinct
cut vertices, $B_1,B_b$ are end blocks with one cut each, and each middle
block $B_i$ ($2\le i\le b-1$) has cut set $W_i=\{w_{i-1},w_i\}$.

\emph{Per-block decompositions.} In each block $B_i$, choose a Hamilton
path $P_i$ whose endpoint set is $W_i$ (for an end block, one endpoint is
its cut vertex and the other is any non-cut vertex). By
Theorem~\ref{thm:conL0}, relabelled, $B_i-P_i$ decomposes into $j_i$ paths
$Q_{i,1},\dots,Q_{i,j_i}$. By Lemma~\ref{lem:conslots} applied to
$(B_i,P_i)$: each cut vertex of $B_i$ is an endpoint of $P_i$ and of
\emph{exactly one} $Q_{i,\cdot}$, and no other path ends at a cut vertex.
Thus each block offers exactly two path-ends at each of its cut vertices,
one from $P_i$ and one from a single $Q_{i,\cdot}$; call the latter the
\emph{$Q$-end of $B_i$ at $w$}.

\emph{Splicing.} At a cut vertex $w_i$ joining $B_i$ and $B_{i+1}$, any
path of $B_i$ ending at $w_i$ and any path of $B_{i+1}$ ending at $w_i$
meet in exactly $\{w_i\}$ (distinct blocks share only their cut vertex),
so they can be spliced. Each splice reduces the path count by one.
Starting from the disjoint union of the per-block decompositions we have
$\sum_i(j_i+1)=\sum_ij_i+b$ paths.

(2) Perform the following $b$ splices:
\begin{itemize}[leftmargin=1.5em]
\item at $w_1$: splice $P_1@w_1$ with $P_2@w_1$, and splice the $Q$-end of
$B_1$ at $w_1$ with the $Q$-end of $B_2$ at $w_1$ (two splices; all four
paths exist and are distinct within their blocks, the two of $B_1$ being
$P_1$ and its unique cut-$Q$, the two of $B_2$ being $P_2$ and its unique
$Q$-end at $w_1$);
\item at each $w_i$, $2\le i\le b-1$: splice the (currently available)
$w_i$-end of the $P$-chain accumulated so far with $P_{i+1}@w_i$ (one
splice each).
\end{itemize}
The $P$-chain after the splices at $w_1,\dots,w_i$ is a single path
contained in $B_1\cup\dots\cup B_{i+1}$ and ending at $w_{i+1}$: indeed
each splice consumes one endpoint of the chain (the one at the cut just
crossed) and one endpoint of $P_{i+1}$, retaining $P_{i+1}$'s other
endpoint, which is $w_{i+1}$. All splices are between paths intersecting
only at their common cut vertex, hence legal. Total splices:
$2+(b-2)=b$, so
\[
\pn(G)\;\le\;\sum_i(j_i+1)-b\;=\;\sum_ij_i\;=\;\floor{n/2}.
\]

(3) For $b\ge4$, add one further splice at the far end $w_{b-1}$: splice
the $Q$-end of $B_{b-1}$ at $w_{b-1}$ with the $Q$-end of $B_b$ at
$w_{b-1}$. The $Q$-end of $B_{b-1}$ at $w_{b-1}$ is still available: in
step (2) the block $B_{b-1}$ only lost the endpoint $P_{b-1}@w_{b-2}$
(consumed by the $P$-chain splice at $w_{b-2}$), not any $Q$-end at
$w_{b-1}$; and $B_b$'s two $w_{b-1}$-paths ($P_b$ and its cut-$Q$) had
only $P_b@w_{b-1}$ consumed. The two spliced paths lie in different blocks
and meet only at $w_{b-1}$. Total splices: $b+1$, giving
\[
\pn(G)\;\le\;\sum_i(j_i+1)-(b+1)\;=\;\sum_ij_i-1\;=\;\floor{n/2}-1. \qedhere
\]
\end{proof}

\begin{remark}[comparison with \cite{CW25} and the computational record]
\label{rem:conchaincomp}
\begin{enumerate}[label={\rm(\arabic*)},leftmargin=2em]
\item Theorem~\ref{thm:conchain}(ii)--(iii) are strict refinements of
\cite{CW25} on this subfamily: block graphs in general are only granted
$\ceil{n/2}$ there, while odd clique chains admit $\floor{n/2}$, and
$\floor{n/2}-1$ from four blocks on. The odd clique chains are far from
arbitrary block graphs: \emph{every} vertex has even degree (a non-cut
vertex of a block $K_{2j_i+1}$ has degree $2j_i$, and a cut vertex has
degree $2j_i+2j_{i+1}$; hence $\EV(G)=G$), which is what the extra
splicing exploits.
\item For $b=3$ the extra splice of (iii) needs the $Q$-ends of the middle
block at its two cuts to be \emph{distinct} paths (for $b\ge4$ this is not
required, by the proof above). We record that the chain of three $K_5$'s
--- here $j=2$, so this example lies outside the domain of
Theorem~\ref{thm:conchain} and illustrates only the distinctness question
above --- admits a $5$-path decomposition, i.e.\
$\pn\le\sum j_i-1=5$, exhibited by:
\begin{center}\small
$a$--$b$--$c$--$d$--$w_1$--$e$--$f$--$g$--$w_2$--$h$--$i$--$j$--$k$;\quad
$d$--$a$--$c$--$w_1$--$f$--$w_2$--$e$;\quad
$e$--$g$--$w_1$--$w_2$--$k$--$h$--$j$;\quad
$a$--$w_1$--$b$--$d$;\quad
$k$--$i$--$w_2$--$j$,
\end{center}
with blocks $B_1=\{w_1,a,b,c,d\}$, $B_2=\{w_1,w_2,e,f,g\}$,
$B_3=\{w_2,h,i,j,k\}$; each line is a simple path and the five paths
partition the $30$ edges. We note that this contradicts the value
$\pn=6$ for this graph recorded in the computational archive of the
underlying project; the record is refuted by the printed decomposition
(Appendix~\ref{com:app-comp}). The general $b=3$ case of the $-1$ refinement
is open and reduces to the distinctness question above.
\item The even-$j$ residual of Theorem~\ref{thm:conL0} (chains containing
blocks of order $4t+1$) is open; existence is verified by computer for
blocks $K_5$ and $K_9$ within chains, but no general construction is known
(Section~\ref{sec:open}).
\end{enumerate}
\end{remark}

\subsection{Formalization and independent-recheck status of this section}
\label{subsec:conformal}

Two statements of this section are fully formalized in Lean, with no
admitted axioms beyond the standard three: for every odd $n\ge3$,
$p(K_{4,n})=(n+1)/2$ (declaration \texttt{pNumK4n\_odd\_all}; hypothesis
$3\le n$ with $n$ odd; an axiom audit returns exactly
\texttt{[propext, Classical.choice, Quot.sound]}), and for every $k\ge1$,
$p(K_{2k+1})=k+1$ (declaration \texttt{pCompleteOdd}, via the edge-count
lower bound and the explicit Walecki-type construction). The odd clique
chain theorem (Theorem~\ref{thm:conchain}) and its engine
(Theorem~\ref{thm:conL0}) have no Lean formalization and rest on paper
proofs; the machine anchors cited above cover the underlying
decompositions, not the chaining argument.

An independent recheck script, written from scratch and run for this
paper, verified: the ring-closure construction of
Theorem~\ref{thm:conringclosure} over all $280$ parameter pairs
($h=1,\dots,14$ and, per $h$, the $20$ even ring sizes $m$ from $2h+2$ to
$2h+41$; simplicity, edge-exactness and the $4h$-edge length profile at
each pair), and the $K_{4,n}$ construction of
Proposition~\ref{prop:conk4large} over all $95$ parameter values
$k=6,\dots,100$ (asserting $k+1$ paths, simplicity, and the exact
$A\times B$ edge partition with multiplicity $1$). These are
construction-level checks: they cover the printed constructions and the
small-case path lists of Proposition~\ref{prop:conk4small} and
Corollary~\ref{cor:concol6}, and do not re-verify the search-based
negatives of Observation~\ref{obs:concostwall} or the splicing argument of
Theorem~\ref{thm:conchain}.

% ============================================================================
% sec-classification.tex — Fan-type theorems; the placeability classification
% programme on K_{2k+1}; criterion-level results.
% ===========================================================================
\section{Fan-type theorems, the placeability classification programme, and
criterion-level results}\label{sec:classification}

\subsection{Fan-type theorems beyond the block-level degree bound}
\label{subsec:astar}

Throughout this subsection $G$ is a connected graph, $n=|V(G)|$, and
$\dEV(v)$ denotes the E-degree of $v$. Recall that a vertex with
$\dEV\ge 4$ has even degree in $G$; call such a vertex \emph{high}, and
call a vertex with $1\le \dEV\le 3$ \emph{low}. All published results in
the Fan line \cite{Fan05,BS21,FHZ23} require the blocks of $\EV(G)$ to have
maximum degree at most $3$ (the triangle conditions of the three differ:
\cite{Fan05} asks every block of $\EV(G)$ to be triangle-free, \cite{FHZ23}
that $\EV(G)$ contain no pair of triangles, and the class $\mathcal{G}$ of
\cite{BS21} allows one triangular block per component). This block
condition is implied by, but does not imply, the absence of high vertices:
a cut vertex of $\EV(G)$ may belong to arbitrarily many blocks of maximum
degree at most $3$ (in $K_{1,5}$ with one pendant leaf attached at each of
its six vertices, $\EV(G)=K_{1,5}$ has a vertex of E-degree $5$ while every
block is a single edge); what is new below is therefore a block of
$\EV(G)$ of maximum degree at least $4$, not a high vertex as such.
Theorem~\ref{thm:claAstar} allows high vertices, under an odd-degree
compensation condition (A2) together with an alignment condition (A4) at
one high vertex; Theorem~\ref{thm:lasep} removes the alignment condition
when distinct high vertices have no common EV-neighbour, and
Remark~\ref{rem:claCO} weakens separation to a parity condition on the
profiles of common EV-neighbours; Theorem~\ref{thm:clatwohigh} treats the
case of exactly two high vertices that share EV-neighbours, and
Corollary~\ref{cor:clatwohigh} collects the at-most-two-high layer,
leaving only a corner open there, subject to five simultaneous constraints
at every common EV-neighbour of the two high vertices; the band-free
tightness theorem (Remark~\ref{rem:clagap}) empties this corner on the
constructed family proper whenever the counting gap is at most two.

\begin{theorem}[alignment at a high vertex]\label{thm:claAstar}
Let $G$ be a connected graph such that
\begin{enumerate}[label={},leftmargin=2.6em]
\item[(A1)] $\EV(G)$ is triangle-free;
\item[(A2)] every high vertex $v$ has odd $\dEV(v)$, and every EV-neighbour
$w$ of $v$ satisfies $\dEV(w)\le 3$;
\item[(A4)] there is a high vertex $u$ such that every even vertex
$x\notin\{u\}\cup N_{\EV}(u)$ satisfies
$\dEV(x)-|N_{\EV}(x)\cap N_{\EV}(u)|\le 3$, where $N_{\EV}(x)$ denotes the
set of EV-neighbours of $x$.
\end{enumerate}
Then $\pn(G)\le \floor{n/2}$.
\end{theorem}

Here $N_{\EV}(x)$ denotes the set of EV-neighbours of $x$, i.e.\ its
neighbours inside $\EV(G)$.\footnote{The numbering skips (A3): a condition
so numbered in an earlier draft of this material was absorbed into the
separation hypothesis of Theorem~\ref{thm:lasep}, and the remaining labels
were kept unchanged.}

\begin{proof}
Let $F=\{uv : v\in N_{\EV}(u)\}$ be the $E$-star at $u$, $|F|=k$, and let
$\tilde G=G-F$. We first analyse the parity and structure of $\tilde G$.

\emph{Step 1 (structure of $\tilde G$).}
Every edge of $F$ joins two even-degree vertices of $G$, so odd-degree
vertices of $G$ lose no edges and stay odd. The vertex $u$ loses $k$ edges,
an odd number, so $u$ is odd in $\tilde G$. Each $v\in N_{\EV}(u)$ loses
exactly one edge, so $v$ is odd in $\tilde G$. Every other vertex keeps its
degree. Hence
$V(\EV(\tilde G))=\{x\notin\{u\}\cup N_{\EV}(u) : \deg_G(x)\text{ even}\}$,
and since $\EV(\tilde G)$ is an induced subgraph of $\EV(G)$, it is
triangle-free by (A1).

Moreover every $x\in V(\EV(\tilde G))$ has $\dEV$ in $\tilde G$ equal to
$\dEV(x)-|N_{\EV}(x)\cap N_{\EV}(u)|\le 3$ by (A4). Thus \emph{every}
vertex of $\EV(\tilde G)$ has E-degree at most $3$, so every block of
$\EV(\tilde G)$ is triangle-free of maximum degree at most $3$. Applying
Theorem~\ref{int:thm-fan} to each component of $\tilde G$ (the hypotheses are
component-wise) and summing with Lemma~\ref{int:lem-sum},
\begin{equation}\label{eq:clatildebound}
\pn(\tilde G)\;\le\;\textstyle\sum_c \floor{n_c/2}\;\le\;\floor{n/2},
\end{equation}
where $n_c$ are the component orders. Fix a decomposition $\tilde D$ of
$\tilde G$ with $|\tilde D|\le\floor{n/2}$.

\emph{Step 2 (hosts).}
The neighbours of $u$ in $\tilde G$ are precisely its odd-degree neighbours
in $G$; they are untouched and odd in $\tilde G$, so by
Lemma~\ref{int:lem-LB1}(ii) each has $\tilde D\ge 1$.

\emph{Step 3 (reinsertion of the addible part).}
Apply Lemma~\ref{int:lem-fan7} with $B=F$, $h=k$, and specified edge $uv_1$
(any fixed element of $F$). The hypothesis of Lemma~\ref{int:lem-fan7} is
satisfied: every neighbour of $u$ has $\tilde D\ge 1$ by Step 2 (this
holds whether the hypothesis is read over the neighbours of $u$ in
$\tilde G$ or in $G$, since the star vertices $v\in N_{\EV}(u)$ are odd in
$\tilde G$, each having lost a single edge, and have $\tilde D\ge 1$ by
Lemma~\ref{int:lem-LB1}(ii)). We obtain $A\subseteq F$ with $|A|\ge\ceil{k/2}$
which can be added back towards $u$, producing a decomposition $D_1$ of
$\tilde G\cup A$ with $|D_1|=|\tilde D|$ and
$D_1(u)\ge \tilde D(u)+\ceil{k/2}\ge 1+\ceil{k/2}$, the bound
$\tilde D(u)\ge1$ coming from Lemma~\ref{int:lem-LB1}(ii) applied to $u$,
which is odd in $\tilde G$.

\emph{Step 4 (reinsertion of the remainder).}
Let $L=F\setminus A$; $|L|=\floor{k/2}=(k-1)/2$. Process the edges of $L$
one at a time. Suppose $j-1$ edges of $L$ have been reinserted, giving a
decomposition $D_j$ of
$\tilde G\cup A\cup(\text{first } j-1 \text{ edges of } L)$ with
\[
D_j(u)\;\ge\;1+\ceil{k/2}-(j-1).
\]
Let $uv\in L$ be the next edge. We apply Lemma~\ref{int:lem-fan6} \emph{towards
$v$}, i.e.\ with the roles of $u$ and $v$ in the statement exchanged: the
condition required is
\[
D_j(u)\;>\;\bigl|\{w\in N_{G_j}(v) : D_j(w)=0\}\bigr|,
\]
where $G_j$ is the current graph, and the operation then yields a
decomposition $D_{j+1}$ of $G_j+uv$ with $D_{j+1}(v)=D_j(v)+1$,
$D_{j+1}(u)=D_j(u)-1$, and $|D_{j+1}|=|D_j|$.

We bound the right-hand side. If $D_j(w)=0$ then $w$ has even degree in
$G_j$ by Lemma~\ref{int:lem-LB1}(iii). We classify the neighbours of $v$ and
their endpoint counts. Odd-degree vertices of $G$ are never incident with
a star edge, keep their degree, and remain odd throughout; they always
have $D\ge1$. A star vertex $v'\in N_{\EV}(u)$, $v'\ne v$, is \emph{not a
neighbour of $v$ at all}: the vertices $v$ and $v'$ are both even-degree
vertices of $G$, so the edge $vv'$ would close the triangle $u,v,v'$ in
$\EV(G)$, contradicting (A1). (This is just as well: a star vertex whose
edge has been reabsorbed towards $u$ by Lemma~\ref{int:lem-fan7} may well have
$D=0$, each absorption step lowering $D(v')$ by one.) The vertex $u$
\emph{is} a neighbour of $v$ and has regained at least $\ceil{k/2}\ge3$
edges, so $D_j(u)\ge3>0$. The only remaining neighbours of $v$ that were
even in $G$ and lie outside $\{u\}\cup N_{\EV}(u)$ are the members of
$N_{\EV}(v)\setminus\{u\}$, of which there are $\dEV(v)-1\le2$. Hence
\[
\bigl|\{w\in N_{G_j}(v) : D_j(w)=0\}\bigr|
\;\le\; \bigl|\{w\in N_{\EV}(v) : w\ne u\}\bigr|
\;=\;\dEV(v)-1\;\le\;2,
\]
using that $v$ is low, i.e.\ $\dEV(v)\le 3$ by (A2). At the last step
$j=(k-1)/2$ we have
$D_j(u)\ge 1+\ceil{k/2}-(j-1)=1+(k+1)/2-(k-3)/2=3>2$, and for earlier steps
the budget is only larger. So each edge of $L$ is addible towards its star
vertex $v$, free of charge.

After all $k$ edges are reinserted we obtain a decomposition of $G$ with
$|D|=|\tilde D|\le\floor{n/2}$, completing the proof.
\end{proof}

The alignment condition is not the only route through a high vertex; the
following theorem needs no alignment hypothesis at all, at the price of a
separation condition between high vertices.

\begin{theorem}[separated high vertices]\label{thm:lasep}
Assume {\rm(A1)}, {\rm(A2)}, and moreover that any two distinct high
vertices of $G$ have no common EV-neighbour. Then $\pn(G)\le\floor{n/2}$.
\end{theorem}
\begin{proof}
We prove the stronger statement, for all graphs in the class connected or
not, by strong induction on $|E(G)|$. If there is no high vertex, every
E-degree is at most $3$ and Theorem~\ref{int:thm-fan} applies component-wise,
so Lemma~\ref{int:lem-sum} gives the bound. Otherwise pick a high vertex $u$
and let $k=\dEV(u)\ge 4$, which is odd by {\rm(A2)}; let $F$ be the E-star
at $u$ and $\tilde G=G-F$. Every star vertex has odd degree in $\tilde G$
(it loses the single edge $uv$), and $u$ has odd degree in $\tilde G$; the
even-degree vertices of $\tilde G$ are precisely the even-degree vertices
of $G$ outside $\{u\}\cup N_{\EV}(u)$, and for such a vertex $x$ its
E-degree in $\tilde G$ equals
$\dEV(x)-|N_{\EV}(x)\cap N_{\EV}(u)|$; when $x$ is high the subtraction
term is $0$, because distinct high vertices have no common EV-neighbour.
Hence $\tilde G$ again satisfies {\rm(A1)}, {\rm(A2)} and the separation
condition, so the induction hypothesis gives a decomposition $\tilde D$ of
$\tilde G$ with $|\tilde D|\le\floor{n/2}$. The neighbours of $u$ in
$\tilde G$ are its odd-degree neighbours in $G$, so each has
$\tilde D\ge 1$ (Lemma~\ref{int:lem-LB1}(ii)); the star vertices are odd in
$\tilde G$, each having lost a single edge, so they too have
$\tilde D\ge1$ by Lemma~\ref{int:lem-LB1}(ii), and the hypothesis of
Lemma~\ref{int:lem-fan7} holds read over $G$; Lemma~\ref{int:lem-fan7} returns an
addible set $A\subseteq F$ with $|A|\ge\ceil{k/2}$ towards $u$, giving
$D_1(u)\ge 1+\ceil{k/2}$. Each edge of $L=F\setminus A$, $|L|=\floor{k/2}$,
is then reinserted towards its star vertex $v$ via Lemma~\ref{int:lem-fan6}:
the vertices $w$ with $D_j(w)=0$ are even in $G_j$, and as in the proof of
Theorem~\ref{thm:claAstar} they are confined to
$N_{\EV}(v)\setminus\{u\}$, of size at most $\dEV(v)-1\le 2$; at the last
step $D_j(u)\ge 3>2$. This rebuilds a decomposition of $G$ with
$\floor{n/2}$ paths.
\end{proof}

\begin{remark}\label{rem:claCO}
The proof of Theorem~\ref{thm:lasep} goes through, verbatim up to one
parity check, with separation replaced by the following weaker condition:
for any two distinct high vertices $x,z$ and any set $S$ of high vertices
disjoint from $\{x,z\}$, the number of common EV-neighbours $w$ of $x$ and
$z$ with $N_{\EV}(w)\cap H\cap S=\emptyset$ is even (equivalently: for
each pair $x,z$, the common EV-neighbours are split by their ``profile''
$\pi(w)=N_{\EV}(w)\cap H$ into classes of even size). Indeed this
condition is again inherited after a star deletion, because a vertex of
$\EV(\tilde G)$ is precisely an even vertex seeing no member of
$\{u\}\cup N_{\EV}(u)$, so the counts for $\tilde G$ with $S'$ are the
original counts for $S=\{u\}\cup S'$.
\end{remark}

Separation and its profile relaxation still leave open the case of high
vertices that do share EV-neighbours. The next theorem settles the case of
exactly two high vertices except for a corner recorded afterwards.

\begin{theorem}[two high vertices]\label{thm:clatwohigh}
Assume {\rm(A1)} and {\rm(A2)}, let $u,x$ be the \emph{only} high vertices
of $G$, and suppose $t_{ux}:=|N_{\EV}(u)\cap N_{\EV}(x)|\ge1$. If either
\begin{enumerate}[label={\rm(\roman*)},leftmargin=2em]
\item some common EV-neighbour $v_1$ of $u$ and $x$ has $\dEV(v_1)=2$, or
\item $x$ has a private EV-neighbour $y^*$, i.e.\
$y^*\in N_{\EV}(x)\setminus N_{\EV}(u)$, with $\dEV(y^*)=1$,
\end{enumerate}
then $\pn(G)\le\floor{n/2}$.
\end{theorem}

\begin{proof}
Let $F$ be the E-star at $u$, $k=\dEV(u)$, and delete one further edge
$e^*$ at $x$: $e^*=xv_1$ in~(i), $e^*=xy^*$ in~(ii), giving
$\tilde G=G-F-e^*$. In $\tilde G$ the vertices $u$ and $x$ are odd, and so
is every former star vertex (each lost the single edge $uv$); in~(ii) the
far endpoint $y^*$ of $e^*$ is odd as well (it lost $xy^*$). The remaining
even vertices --- in~(i) the vertex $v_1$ (which lost both $uv_1$ and
$xv_1$ and stays even) together with the far even vertices and the private
EV-neighbours of $x$; in~(ii) the far even vertices other than $y^*$ and
the private EV-neighbours of $x$ other than $y^*$ --- all have E-degree at
most $3$ in $\tilde G$: the private ones by {\rm(A2)} at $x$, the far ones
because $u$ and $x$ are the only high vertices, and $v_1$ because
$\dEV(v_1)=2$. Hence $\EV(\tilde G)$ is triangle-free, being an induced
subgraph of $\EV(G)$, of maximum E-degree at most $3$, so $\tilde G$ lies
in the class of Theorem~\ref{int:thm-fan} and has, component-wise with
Lemma~\ref{int:lem-sum}, a decomposition $\tilde D$ with
$|\tilde D|\le\floor{n/2}$. In~(i) reinsert $e^*$ towards $v_1$
(Lemma~\ref{int:lem-fan6}: $\tilde D(x)\ge1$ since $x$ is odd, while every
neighbour of $v_1$ in $\tilde G$ is odd except $x$ --- the even-degree
neighbours of $v_1$ in $G$ are exactly $u$ and $x$ --- so the passing
count is $0$); this activates $v_1$. In~(ii) no activation is needed:
every star vertex has lost exactly one edge and is odd. Then
Lemma~\ref{int:lem-fan7} (hypothesis read over $G$, as in the proof of
Theorem~\ref{thm:claAstar}) returns $A\subseteq F$ with $|A|\ge\ceil{k/2}$
addible towards $u$, and each edge of $L=F\setminus A$ is reinserted
towards its star vertex as in the proof of Theorem~\ref{thm:claAstar}
(passing count $\le\dEV(v)-1\le2$, budget $\ge3$ at the last step). In~(ii)
the edge $xy^*$ is reinserted towards $y^*$, whose neighbours except $x$
are all odd, so this step is legal at any position of the sequence. All
$k+1$ deleted edges are restored with $|D|=|\tilde D|\le\floor{n/2}$.
\end{proof}

\begin{corollary}\label{cor:clatwohigh}
If $G$ satisfies {\rm(A1)}--{\rm(A2)} and has at most two high vertices,
then $\pn(G)\le\floor{n/2}$, unless $G$ has exactly two high vertices $u,x$
with $t_{ux}\ge1$, every common EV-neighbour of $u$ and $x$ of E-degree
$3$, no private EV-neighbour of E-degree $1$ on either side, and such
that, for \emph{every} common EV-neighbour $v$ of $u$ and $x$, writing $z$
for the third EV-neighbour of $v$, $F$ for the E-star at $u$ and
$\tilde G=G-F-\{xv\}$, all of the following hold:
$\pn(\tilde G)=\floor{n/2}$; $\beta(\tilde G,x)=\floor{n/2}+1$ (cf.\
Section~\ref{int:subsec-budget}); $z$ is buried (i.e.\ $D(z)=0$) in every
optimal decomposition of $\tilde G$; in every optimal decomposition
$\tilde D$ of $\tilde G$, either $\tilde D(v)=0$, or $\tilde D(v)\ge2$ and
at least $\tilde D(v)$ private EV-neighbours of $x$ have $\tilde D=0$; and
$\tilde G$ has at most one component of odd order.
\end{corollary}
\begin{proof}
$t_{ux}=0$ is Theorem~\ref{thm:lasep} (with only two high vertices the
separation condition reduces to $t_{ux}=0$). The two alternatives are
Theorem~\ref{thm:clatwohigh}. For the remaining graphs fix a common
EV-neighbour $v$, write $z$ for its third EV-neighbour and
$\tilde G=G-F-\{xv\}$; then $\EV(\tilde G)$ again has maximum E-degree
$\le3$, and each of the following five routes closes the case for this
$v$, the star at $u$ being absorbed afterwards exactly as in the proof of
Theorem~\ref{thm:clatwohigh}:
\begin{enumerate}[label={\rm(\alph*)},leftmargin=2em]
\item $\pn(\tilde G)\le\floor{n/2}-1$; take an optimal decomposition
$\tilde D$. If $\tilde D(x)\ge3$ we are on route~(b). Otherwise
$\tilde D(x)=1$ (it is odd by Lemma~\ref{int:lem-LB1}(iii), $x$ being odd in
$\tilde G$), and since $\deg_{\tilde G}(x)=\deg_G(x)-1\ge5$, the unique
path ending at $x$ continues through $x$; cutting it at $x$ adds one path
(at most $\floor{n/2}$ in all) and raises the endpoint count at $x$ to $3$
--- route~(b);
\item a decomposition $\tilde D$ with $|\tilde D|\le\floor{n/2}$ and
$\tilde D(x)\ge3$: then $xv$ is addible towards $v$, the passing count at
$v$ being at most $1$ ($z$ is the only neighbour of $v$ in $\tilde G$ that
can be dead; all other neighbours are odd);
\item such a decomposition with $\tilde D(x)=1$ and $\tilde D(z)\ge1$: the
passing count at $v$ drops to $0$;
\item such a decomposition with $\tilde D(v)\ge2$ larger than the number of
dead private EV-neighbours of $x$: reinsert $xv$ towards $x$ (the passing
count at $x$ sees only private EV-neighbours of $x$, the other neighbours
of $x$ being active), so that $v$ remains active and the full star,
including $uv$, is absorbable;
\item $\tilde G$ has at least two components of odd order: the
component-wise Theorem~\ref{int:thm-fan}, summed with Lemma~\ref{int:lem-sum},
gives route~(a).
\end{enumerate}
The budget arithmetic of the reinsertion skeleton behind routes (a)--(d),
and the formal statement that the exception of this corollary is exactly
the simultaneous failure of all five routes at \emph{every} common
EV-neighbour, are given in Appendix~\ref{com:app-band}.
The stated exception is exactly the case in which, for every common
EV-neighbour $v$, none of the five routes applies.

\emph{Band closure of the corner.} One further narrowing is available in
the constructed family of Remark~\ref{rem:clagap}. Route~(a) can fail
there only when $\tilde G$ carries a \emph{deadlock band} (an even vertex
of $\tilde G$ with more corridor neighbours than terminal units; see
Remark~\ref{rem:clagap} for the definition), because a band-free
$\tilde G$ in the family has
$\pn(\tilde G)=\mathrm{odd}(\tilde G)/2$; and the counting theorem forces
$\tilde G$ to be band-free whenever the counting gap
$\floor{n/2}-\mathrm{odd}(\tilde G)/2$ --- which is at least $1$ in this
situation --- is at most $2$, on the family proper (the constructed family
without cross-private pairs). Hence for every family-proper instance with
counting gap at most $2$, route~(a) applies unconditionally: the corner is
empty at theorem level, not merely computationally; the cross-private
instances are backed at probe level by the machine certificates of
Appendix~\ref{com:app-comp}. For larger gaps the corner is blocked only by the
excess bound recorded in Remark~\ref{rem:clagap}.
\end{proof}

\begin{corollary}\label{cor:clauniquehigh}
If $G$ satisfies (A1) and (A2) and has \emph{exactly one} high vertex,
then $\pn(G)\le\floor{n/2}$.
\end{corollary}

\begin{proof}
Let $u$ be the unique high vertex. Every vertex other than $u$ is low, so
every even vertex $x\notin\{u\}\cup N_{\EV}(u)$ satisfies
$\dEV(x)-|N_{\EV}(x)\cap N_{\EV}(u)|\le\dEV(x)\le 3$; this is condition
(A4) of Theorem~\ref{thm:claAstar}. (Alternatively, with only one high
vertex the separation hypothesis of Theorem~\ref{thm:lasep} holds
vacuously.)
\end{proof}

We now formulate the statement with neither the alignment nor the
separation condition, which is the form investigated in the computational
study behind this paper and which all our computations support.

\begin{conjecture}\label{conj:claAstar}
Let $G$ be a connected graph such that $\EV(G)$ is triangle-free, and
every vertex $v$ with $\dEV(v)\ge 4$ has odd $\dEV(v)$ and only
EV-neighbours $w$ with $\dEV(w)\le 3$. Then $\pn(G)\le\floor{n/2}$.
\end{conjecture}

\begin{remark}[where the general argument stops]\label{rem:clagap}
The class {\rm(A1)--(A2)} is not preserved by the only deletion the engine
supports. Deleting the E-star at a high vertex $u$ turns every star vertex
odd; an even vertex $x$ outside the star keeps its degree but loses the
E-degrees contributed by star vertices, and its E-degree in $\tilde G=G-F$
equals $\dEV(x)-t_x$ with $t_x=|N_{\EV}(x)\cap N_{\EV}(u)|$. If $x$ is a
high vertex with $t_x$ odd and $\dEV(x)-t_x\ge 4$, then $x$ sits in
$\EV(\tilde G)$ at an even E-degree, the class is lost, and neither
Theorem~\ref{int:thm-fan} nor the induction applies to $\tilde G$. Concretely,
in the double star on
$\{u,x,v_1,w_1,\dots,w_4,y_1,\dots,y_4\}$ with edges
$uv_1,uw_1,\dots,uw_4,xv_1,xy_1,\dots,xy_4$ and pendant leaves attaching
so that the eleven listed vertices have even degree, both $u$ and $x$ are
high with $\dEV=5$, $t_{ux}=1$ is odd, and after deleting the star at $u$
the vertex $x$ has E-degree $4$, even. This graph (and its even-order
one-leaf extension, on which our exact programs give $\pn=5\le 11$)
satisfies {\rm(A1)--(A2)} but none of {\rm(A4)}, the separation condition,
or the profile condition of Remark~\ref{rem:claCO}; a minimal
counterexample to Conjecture~\ref{conj:claAstar}, if it exists, must have
a high vertex with such a bad partner at {\rm every} high vertex
(Proposition~\ref{prop:claresidue}). No counterexample is known: exact
dynamic programs verify $\pn\le\floor{n/2}$ on all tested instances
satisfying {\rm(A1)--(A2)} ($166$ instances in the computational study
behind this paper, a further $413$ in Appendix~\ref{com:app-comp}, zero
violations), and the family above sits far from the bound
($\pn=\mathrm{odd}(G)/2$).

The ``one endpoint unit'' obstruction recorded in an earlier draft of this
material was an artefact: the reinsertion bookkeeping it refers to
restores $k+2$ edges after deleting $k+1$, i.e.\ counts the patched edge
$xv_1$ twice (an activation step towards $v_1$ and a final step towards
$x$). Once $x$'s endpoint loss is recognised as cost-free --- $|D|$ is the
only invariant, and the argument of Theorem~\ref{thm:clatwohigh} involves
no induction --- the obstruction disappears for all two-high graphs except
the corner recorded in Corollary~\ref{cor:clatwohigh}.

The exceptional graphs of Corollary~\ref{cor:clatwohigh} are subject to
five simultaneous constraints, of which only the first two were recorded
when the corner was first isolated: the burial of the third EV-neighbour
$z$ and of $v$ (or of enough private neighbours of $x$) in \emph{every}
optimal decomposition of $\tilde G$ (with $\tilde G$ as in
Corollary~\ref{cor:clatwohigh}), and the parity of the components of
$\tilde G$, are each individually sufficient to reopen the reinsertion
argument (the edge $xv$ can then be added towards $v$ at zero passing
cost, or a second deletion $G-\tilde G$ becomes unnecessary). No instance
satisfying even the original two constraints was found: across $207$
constructed instances of the dirty class---including the entire range
where the counting gap $\floor{n/2}-\mathrm{odd}(\tilde G)/2$ attains its
minimum $1$, which the degree capacity of the two far vertices forces to a
single parameter-free shape---a decomposition of $\tilde G$ with at most
$\floor{n/2}-1$ paths was always produced, so the corner never materializes
in computation.

The band-free tightness statement is now a theorem for the \emph{constructed
family} (the two-high graphs satisfying the invariants of our construction
pipeline: no $\EV$-edge joins two far even vertices of $\tilde G$; the
odd-neighbour blocks of $u$ and of $x$ are of leaf type (an odd number of
pendant leaves) or gadget type, possibly mixed; every far vertex of
$\tilde G$ has degree at most $4$; the $207$ instances above and the $125$
decision-level instances of Appendix~\ref{com:app-comp} all lie in it).
Writing, for an even vertex $f$ of $\tilde G$, $Y(f)$ for its
\emph{corridor} neighbours --- the degree-two even vertices adjacent to
both $x$ and $f$ --- and $N(f)$ for its neighbour set in $\tilde G$, call
$f$ a \emph{deadlock band} if it has more corridor neighbours than
terminal units, where the terminal units of $f$ are its remaining incident
edges (in symbols, $|Y(f)|>|N(f)|-|Y(f)|$). If no even vertex is a
deadlock band, then $\tilde G$ admits a \emph{paired} decomposition:
exactly $\mathrm{odd}(\tilde G)/2$ paths with all and only the odd
vertices of $\tilde G$ as endpoints. The proof is a five-step pairing
scheme---self-contained odd-neighbour blocks, forced fusions of the
corridors and of the unit at the common neighbour, an injective pairing of
corridor ends with distinct terminal units at every far vertex (which is
precisely the negation of the band condition), and a sorted offset
matching at $x$ that never joins two corridors of the same far vertex,
possible because a far vertex of the family has degree at most $4$ and
hence at most two corridors. Every resulting trail visits $x$ once and
leaves each far vertex for good, so no trail closes. A single band, on
the other hand, costs exactly one path: spending the free edge of $x$ on a
surplus corridor makes that corridor a double endpoint and lets its
partner ride through $x$ to a terminal, matching the exact value
$\pn(\tilde G)=\mathrm{odd}(\tilde G)/2+1$ computed by dynamic program.
Consequently the corner of Corollary~\ref{cor:clatwohigh} is empty
whenever the counting gap is at most $2$, unconditionally on the
constructed family proper, with the cross-private extension covered at
probe level by the machine certificates of Appendix~\ref{com:app-comp}; for
larger gaps it remains blocked only by the excess bound
$\mathrm{excess}\le\lfloor y/2\rfloor$, where
$\mathrm{excess}:=\pn(\tilde G)-\mathrm{odd}(\tilde G)/2$ and $y$ is the
number of corridors --- verified on all $235$ computed instances of the
family, the provable bound being
$\pn(\tilde G)\le \mathrm{odd}(\tilde G)/2+2B-1$ for $B$ bands, with a
single band forcing exactly $+1$. The restriction to the constructed
family is now known to be sharper than necessary in two directions and
irreducible in a third. Far--far edges of the even subgraph and far
vertices carrying several parity leaves, both excluded by the
construction, can be admitted: the corridor ends at a far vertex $f$
admit an injective pairing with terminal units and far--far ends whenever
$|Y(f)|\le t(f)+d(f)$ (where $t(f)$ is the number of terminal units of
$f$ and $d(f)$ its far--far degree), and the matching at $x$ survives
arbitrary corridor multiplicities because a far vertex of the class has at
most three even neighbours, so a rearrangement argument replaces the
sorted offset. What cannot be admitted are arbitrary odd-neighbour
blocks: a block whose component has exactly two odd vertices admits a
paired decomposition if and only if it admits a Hamilton path between
them, and the cycle-with-one-pendant and double-triangle blocks have
none, so $\tilde G$ can be untight while carrying no band at all---the
band-free tightness statement is a phenomenon of the constructed family,
not of the class, and these cycle blocks are the first known
counterexamples to the general band-free statement. (Same-side pairs of
private neighbours, a third candidate obstruction, are excluded by the
triangle-free condition {\rm(A1)} already.) The counting gap of the dirty
class also assumes that every even vertex lies in the skeleton; blocks
with even interior vertices shift $\floor{n/2}$ and must be recomputed
directly (see Appendix~\ref{com:app-comp} for the machine certificates).
\end{remark}

\begin{proposition}\label{prop:claresidue}
Let $G$ be a minimal counterexample to Conjecture~\ref{conj:claAstar}.
Then every high vertex $u$ of $G$ admits a far high vertex $x$ with
$|N_{\EV}(x)\cap N_{\EV}(u)|$ odd and
$\dEV(x)-|N_{\EV}(x)\cap N_{\EV}(u)|\ge4$. In particular, by
Theorem~\ref{thm:clatwohigh} a minimal counterexample has at least three
high vertices, unless it is one of the two-high exceptional graphs of
Corollary~\ref{cor:clatwohigh} (empty on the constructed family proper
whenever the counting gap is at most two, by the band-free tightness
material of Remark~\ref{rem:clagap}; no instance was found in our
computations, in or beyond the family).
\end{proposition}
\begin{proof}
Read ``minimal'' as minimal with respect to the number of edges, and let
$u$ be a high vertex of $G$. We show that if $u$ has no bad partner, then
$\pn(G)\le\floor{n/2}$, contradicting the choice of $G$; hence $u$ has a
bad partner. Suppose, then, that no far even vertex $x$ satisfies $t_x$
odd and $\dEV(x)-t_x\ge4$, where $t_x=|N_{\EV}(x)\cap N_{\EV}(u)|$. Let
$F$ be the E-star at $u$ and $\tilde G=G-F$. As in Step~1 of the proof of
Theorem~\ref{thm:claAstar}, the even-degree vertices of $\tilde G$ are
precisely the far even vertices of $G$, $\EV(\tilde G)$ is an induced
subgraph of $\EV(G)$ and hence triangle-free, and a far vertex $x$ has
E-degree $\dEV(x)-t_x$ in $\tilde G$. Thus every high vertex of
$\tilde G$ is a far even vertex $x$ of $G$ with $\dEV(x)-t_x\ge4$; by our
assumption $t_x$ is even for such an $x$, and since $\dEV(x)\ge4$ makes
$x$ high in $G$, (A2) gives $\dEV(x)$ odd, so the E-degree
$\dEV(x)-t_x$ of $x$ in $\tilde G$ is odd. Moreover the EV-neighbours $w$
in $\tilde G$ of such an $x$ satisfy
$\dEV(w)-t_w\le\dEV(w)\le3$, the last inequality by (A2) applied to the
high vertex $x$ in $G$. Hence $\tilde G$ satisfies {\rm(A1)} and
{\rm(A2)}. Each component of $\tilde G$ is a connected graph satisfying
{\rm(A1)}--{\rm(A2)} with fewer edges than $G$ (as $u$ is high,
$F\neq\emptyset$), so by minimality each admits a decomposition with at
most $\floor{n_c/2}$ paths, and Lemma~\ref{int:lem-sum} gives a decomposition
$\tilde D$ of $\tilde G$ with $|\tilde D|\le\floor{n/2}$. Now the
reinsertion argument of Steps~3 and~4 of the proof of
Theorem~\ref{thm:claAstar} uses (A4) only through the existence of such a
$\tilde D$: the star vertices are odd in $\tilde G$, so the hypothesis of
Lemma~\ref{int:lem-fan7} holds read over $G$ exactly as in Step~3, and each
edge of $F\setminus A$ is addible towards its star vertex against a
passing count of at most $\dEV(v)-1\le2$ (by (A2)) and a budget of at
least $3$. The full star is restored with $|D|=|\tilde D|\le\floor{n/2}$
paths --- a contradiction.
\end{proof}

We also record a structural consequence of (A1) on odd orders, used below
to place the sharpness family of Proposition~\ref{prop:claoscsharp}
relative to the class.

\begin{proposition}\label{prop:claodddensity}
Let $G$ be connected of odd order $n$ with $\EV(G)$ triangle-free. Then
$m\le (n-1)^2/2$; equivalently, no odd semi-clique has triangle-free
$\EV(G)$.
\end{proposition}
\begin{proof}
Suppose $c=\binom{n}{2}-m\le (n-3)/2$. Let $S$ be the set of even-degree
vertices, $s=|S|$; since $n$ is odd, a vertex has odd degree iff its
complement-degree is odd, so each vertex outside $S$ has positive
complement-degree. Let $t$ be the number of complement edges inside $S$.
Since $G[S]=\EV(G)$ is triangle-free, Mantel's theorem gives
$\binom{s}{2}-t\le\floor{s^2/4}$, i.e.\ $t\ge \binom{s}{2}-\floor{s^2/4}$.
Counting complement degrees over $S$ and over its complement separately,
$2c\ge 2t+(n-s)$, so $2t\le s-3$. If $s$ is even this reads
$s(s-2)/2\le s-3$, i.e.\ $s^2-4s+6\le0$; if $s$ is odd,
$(s-1)^2/2\le s-3$, i.e.\ $s^2-4s+7\le0$; both are impossible. (Note
$s\ge1$: the number of odd-degree vertices is even while $n$ is odd.)
\end{proof}

We record the sharpness of the triangle-free requirement. By
Proposition~\ref{prop:claodddensity} this failure is a density phenomenon:
on odd orders, high density and (A1) are mutually exclusive, so the
following family is forced.

\begin{proposition}\label{prop:claoscsharp}
For every $k\ge 2$, the graph $G_k=K_{2k+1}-(k-1)K_2$ (delete a matching
of size $k-1$) satisfies $\EV(G_k)\cong K_3$, hence violates (A1), and
$\pn(G_k)\ge k+1=\ceil{n/2}>\floor{n/2}$. Moreover $\pn(G_k)=k+1$ for the
tested values $k\in\{2,3,4,5,6,8,10\}$, with explicit decompositions
machine-verified.
\end{proposition}

\begin{proof}
The $k-1$ matched vertices have degree $2k-1$ and the three unmatched
vertices have degree $2k$; the latter are pairwise adjacent, so
$\EV(G_k)$ is a triangle on them. The edge count is
$m=\binom{2k+1}{2}-(k-1)=2k^2+1$, so Lemma~\ref{int:lem-LB2} gives
$\pn\ge\ceil{(2k^2+1)/2k}=k+1$; equivalently $G_k$ is an odd semi-clique
and Observation~\ref{int:obs-density} applies. This proves the necessity
claim for all $k$. For the tested values, explicit decompositions into
$k+1$ paths (obtained from the Walecki cycle decomposition of $K_{2k+1}$,
spreading the matching edges over distinct cycles, cutting each affected
cycle at its matching edge and splitting the untouched cycle into a path
plus a single edge) were verified by machine to partition $E(G_k)$
exactly; for $k=2$ the value $\pn(K_5-M_1)=3$ is also certified by exact
dynamic programming.
\end{proof}

\begin{remark}
The exactness claim in Proposition~\ref{prop:claoscsharp} is stated only
for tested $k$: the existence, for all $k$, of a Walecki-type cycle
decomposition of $K_{2k+1}$ in which a given $(k-1)$-matching is spread
over distinct cycles is plausible but we do not prove it.
\end{remark}

Next, the witness family $W(m)$, which motivates Theorem~\ref{thm:claAstar}:
it lies in no published class we are aware of, yet satisfies all
hypotheses.

\begin{figure}[h]
\centering
\begin{minipage}{0.92\textwidth}\small
For odd $m\ge 5$, let $W(m)$ have vertex set
$\{a_1,a_2\}\cup\{b_1,\dots,b_m\}\cup\{\ell_1,\ell_2\}\cup\{h,h'\}\cup\{x_1,\dots,x_6\}$
and edge set
\[
E(K_{2,m}) \;\cup\; \{a_1\ell_1,\,a_2\ell_2\} \;\cup\;
\{b_1h,\,b_1h'\} \;\cup\; \{hx_1,\,hx_2,\,h'x_1,\,h'x_2\} \;\cup\; E(K_6),
\]
where $K_6$ has vertex set $\{x_1,\dots,x_6\}$. Precisely: $h$ and $h'$
each have neighbourhood $\{b_1,x_1,x_2\}$.
\end{minipage}
\caption{The witness family $W(m)$.}\label{fig:claWm}
\end{figure}

\begin{lemma}\label{lem:claWmEV}
For every odd $m\ge 5$, the graph $W(m)$ of Figure~\ref{fig:claWm}
satisfies $\EV(W(m))\cong K_{2,m}$ exactly, hence satisfies (A1) and (A2)
with $\dEV(a_1)=\dEV(a_2)=m$ and $\dEV(b_j)=2$ for all $j$.
\end{lemma}

\begin{proof}
The degrees are:
$\deg(a_i)=m+1$, even; $\deg(b_j)=2$ for $j\ge2$ and $\deg(b_1)=4$, all
even; $\deg(\ell_i)=1$, $\deg(h)=\deg(h')=3$, and $\deg(x_i)=5+2=7$ for
$i\in\{1,2\}$, $\deg(x_i)=5$ for $i\in\{3,\dots,6\}$, all odd. Hence the
even-degree vertex set is exactly $\{a_1,a_2\}\cup\{b_1,\dots,b_m\}$, and
the induced subgraph is $K_{2,m}$: triangle-free, with $\dEV(a_i)=m$ odd
and all EV-neighbours of $a_i$ of E-degree $2\le3$. Note that the two
degree-$3$ vertices $h,h'$ attach to $b_1$ \emph{in pairs} (both $h$ and
$h'$), keeping $b_1$ even; attaching a single such vertex would flip $b_1$
to odd degree and destroy the structure.
\end{proof}

\begin{proposition}\label{prop:claWmbound}
For every odd $m\ge 5$, $\pn(W(m))\le\floor{n/2}$ with $n=m+12$.
\end{proposition}

\begin{proof}
By Lemma~\ref{lem:claWmEV} the graph satisfies (A1), (A2). We verify (A4)
at $u=a_1$: the only even-degree vertex outside
$\{a_1\}\cup\{b_1,\dots,b_m\}$ is $a_2$, and, since
$N_{\EV}(a_1)=N_{\EV}(a_2)=\{b_1,\dots,b_m\}$,
$\dEV(a_2)-|N_{\EV}(a_2)\cap N_{\EV}(a_1)|=m-m=0\le3$; every vertex
outside this set other than $a_2$ has odd degree, hence is not in
$\EV(\tilde G)$ after the deletion and is irrelevant (its E-degree in
$\tilde G$ is $0$). Thus (A4) holds, and Theorem~\ref{thm:claAstar} gives
$\pn(W(m))\le\floor{(m+12)/2}$. For $m\in\{5,7,9\}$ explicit
decompositions into exactly $\floor{n/2}=8,9,10$ paths were in addition
machine-verified.
\end{proof}

We now compare Theorem~\ref{thm:claAstar} with the published record. The
following two statements make the incomparability precise; both witnesses
have exactly computed path numbers.

\begin{proposition}[escape]\label{prop:claescape}
Let $m\ge5$ be odd and let $W(m)$ be as in Figure~\ref{fig:claWm}. Then
$\EV(W(m))\cong K_{2,m}$, whence Theorem~\ref{thm:claAstar} gives
$\pn(W(m))\le\floor{n/2}$ with $n=m+12$. No published theorem known to us
implies this bound:
\begin{itemize}[leftmargin=1.5em]
\item every member of the Fan line \cite{Fan05,BS21,FHZ23} requires every
block of $\EV(G)$ to have maximum degree at most $3$; the graph
$\EV(W(m))\cong K_{2,m}$ is $2$-connected, hence a single block of maximum
degree $m\ge5>3$, so this requirement fails (the block level, not the
global value of $\Delta(\EV)$, is the discriminating parameter ---
compare $K_{1,5}$ with one pendant leaf at each of its six vertices, where
$\Delta(\EV)=5$ while every block of $\EV$ is a single edge);
\item $W(m)$ contains a $K_6$, so $\mathrm{tw}(W(m))\ge5$ and its
degeneracy is at least $5$, which excludes \cite{LY24} (treewidth $\le4$),
\cite{BCLS20} (treewidth $\le3$) and \cite{ZLH24} ($3$-degenerate);
\item $\Delta(W(m))=\max(m+1,7)\ge7$, which excludes \cite{BP19}
($\Delta\le5$) and \cite{CFL} ($\Delta\le6$ with an independence
condition);
\item $\Delta(W(m))\ge7$, so $W(m)$ is not subcubic and the
polynomial-time algorithm of \cite{MY26} for subcubic graphs does not
apply;
\item $\EV(W(m))\cong K_{2,m}$ contains a $4$-cycle, so it is not a forest
and \cite{Pyb96} does not apply;
\item $W(m)$ is not Eulerian, so \cite{FK88} does not apply; it is not a
block graph, so \cite{CW25} does not apply.
\end{itemize}
For $m\in\{5,7,9\}$ every membership and non-membership claim above,
together with explicit decompositions into $\floor{n/2}$ paths, is
machine-certified.
\end{proposition}

\begin{proposition}[proper inclusion in both directions]\label{prop:claincomp}
The class of Theorem~\ref{thm:claAstar} neither contains nor is contained
in the Botler--Sambinelli class $\mathcal{G}$ \cite{BS21} or the
Fan--Hou--Zhou class \cite{FHZ23}.
\begin{enumerate}[label={\rm(\roman*)},leftmargin=2em]
\item Let $G$ be $K_{2,5}$ with one pendant leaf attached to each of the
two degree-$5$ vertices. Then $\EV(G)\cong K_{2,5}$ is triangle-free, the
two degree-$5$ vertices have odd E-degree $5$ with all E-neighbours of
E-degree $2$, so $G$ satisfies (A1)--(A2); but the graph
$\EV(G)\cong K_{2,5}$ is $2$-connected, hence a single block of maximum
degree $5>3$, so $G$ lies in none of \cite{Fan05,BS21,FHZ23}. Moreover
$\pn(G)=3\le\floor{9/2}$ by exact computation. (For the record, $G$ is
$2$-degenerate, so its bound is also implied by \cite{ZLH24}; this witness
is used only for the incomparability with the Fan line.)
\item Let $G'$ be a triangle $\{t_1,t_2,t_3\}$ together with two vertices
$w_1,w_2$ adjacent only to $t_1$ and to three pendant leaves each, so that
$\EV(G')$ consists of the triangle $\{t_1,t_2,t_3\}$ plus the two bridge
edges $t_1w_1$, $t_1w_2$. Every block of $\EV(G')$ has maximum degree at
most $3$ and its single component contains exactly one block with
triangles, so $G'$ lies in $\mathcal{G}$ and in the Fan--Hou--Zhou class;
yet the triangle in $\EV(G')$ violates (A1), so Theorem~\ref{thm:claAstar}
is silent on $G'$. Note $G'$ is not one of the degenerate exceptions of
the ``Gallai-or-SET'' alternative of \cite{BS21}. Exactly,
$\pn(G')=4\le\floor{11/2}$ by exact computation, while the main theorem of
\cite{BS21} gives only $\pn(G')\le\ceil{11/2}=6$ on its class (the
Botler--Sambinelli bound is stated with Gallai's ceiling, as noted in
Section~\ref{sec:intro}); the Fan--Hou--Zhou membership of $G'$ rests on
the block structure of $\EV(G')$ alone and is unaffected.
\end{enumerate}
\end{proposition}

\begin{remark}[formalization and provenance status of the Fan-type chain]
\label{rem:claformal}
The Lean formalization of the statements of this subsection is in
progress: in the current Lean sources the declarations corresponding to
Theorem~\ref{thm:claAstar}, Theorem~\ref{thm:lasep} and their supporting
reinsertion lemmas (\texttt{astar}, \texttt{astar\_separated},
\texttt{fan6\_use}, \texttt{fan7\_use}, \texttt{fan\_class}) carry
admitted (``\texttt{sorry}'') placeholders --- five in total --- so none
of them may be quoted as machine-verified or formally proved, and the
paper cites the chain as paper proofs with formalization in progress.
One research-technical lemma of the formalization archive
(\texttt{fan6\_nf}, a pointwise-profile strengthening of
Lemma~\ref{int:lem-fan6}) was refuted by a machine counterexample, is sealed
as a research record, and is never cited in this paper; the Lean
formalization of Lemma~\ref{int:lem-fan6} uses the four-clause form obtained
by deleting the final pointwise clause $D(w)=D'(w)$ for the remaining
vertices, and no application in this paper consumes that clause. On the
provenance side: the statements of Lemma~\ref{int:lem-fan6} and
Lemma~\ref{int:lem-fan7} follow the restatement in \cite{BS21} (Lemmas 6 and
7 there); the proofs in \cite{Fan05} could not be consulted and the
transcription has not been independently checked against the original.
The machine evidence is direction-corrected: the bidirectional rescan of
Lemma~\ref{int:lem-fan6} is complete at $n\le6$, all
$1{,}701{,}141{,}150$ ordered instances at $n=6$ (the fixed-orientation
half exactly reproducing the historical census of $850{,}570{,}575$), with
zero violations.
\end{remark}

\subsection{The placeability classification programme on
\texorpdfstring{$K_{2k+1}$}{K(2k+1)}}\label{subsec:echar}

\subsubsection{The layer domain}\label{subsubsec:layerdomain}

Fix the layer domain: the carrier is $K_{2k+1}$ on
$V=\mathbb{Z}_{2k}\cup\{\infty\}$; a \emph{type} is an isomorphism class of
$k$-edge loopless graphs without isolated vertices; and for such a type $F$
we say that $F$ is \emph{sharp at layer $k$} ($\mathrm{sharp}_k(F)$) if
there is an injection $\phi:V(F)\to V$ whose $k$ image-edge classes hit
$\mathbb{Z}_k$ exactly, where the edge classes are given by the closed
form
\begin{equation}\label{eq:claclass}
\mathrm{class}(\{u,v\})=\bigl\lfloor(u+v)/2\bigr\rfloor \bmod k,
\qquad
\mathrm{class}(\{\infty,y\})=y \bmod k .
\end{equation}
A placement realizing $\mathrm{sharp}_k(F)$ is a \emph{witness}. Write
\[
\mathcal{E}_k:=\{F:\ \neg\,\mathrm{sharp}_k(F)\}
\]
for the \emph{exception set} at layer $k$, and $K_n@k$ for the sharpness
question for $K_n$ at layer $k$. The values $k=\binom n2$ are the
\emph{triangular layers}; $K_3@3$, $K_4@6$ and $K_5@10$ are sharp (three
positive controls), so $n\ge6$ is the exact lower bound for exceptions
among complete graphs.

\subsubsection{The first four complete-graph layers}\label{subsubsec:firstlayers}

The exact status of the first four critical layers is as follows; tiers
are quoted verbatim and travel with every use of the numbers.

\begin{enumerate}[leftmargin=2.2em,itemsep=0.2em]
\item \textbf{$K_6@15$ (canonical).} $K_6$ is not sharp at $k=15$, and
$\mathcal{E}_{15}=\{K_6\}$ exactly --- the only layer at which the
uniqueness side is closed. The evidence is four-fold: three independent
depth-first implementations exhaust the same $145{,}923$-node search tree
without witness and without cap, and a direct combinatorial exhaustion of
all $\binom{31}{6}=736{,}281$ six-subsets of the carrier (the image-edge
set of a $K_6$ placement depends only on the six chosen vertices) finds
zero placements with $15$ distinct classes.
\item \textbf{$K_7@21$ (machine universal-closure candidate, not
independently audited).} $K_7$ is not sharp at $k=21$, closed on three
independent lines: an engine exhaustion of $2{,}868{,}963$ nodes without
witness and without cap; a regime census of $3{,}071{,}958$ configurations;
and an infinity-census of $3{,}472{,}896$ configurations (all zero hits).
\item \textbf{$K_8@28$ (machine universal-closure candidate, not
independently audited).} $K_8$ is not sharp at $k=28$: a nine-line census
over $354{,}484{,}845$ configurations with zero hits, zero truncation, and
no sieve dependence (the sieve-excluded regimes were also actually
scanned and are empty).
\item \textbf{$K_9@36$ (machine universal-closure candidate, not
independently audited).} $K_9$ is not sharp at $k=36$: a configuration-level
true enumeration over $18{,}385{,}772{,}680$ configurations (nine regime
lines summing to $\binom{71}{8}=10{,}639{,}125{,}640$ --- the Vandermonde
identity is machine-checked --- plus the infinity line
$\binom{36}{8}\cdot2^8=7{,}746{,}647{,}040$, covered in four shards with an
exact non-overlapping-partition assertion), every line zero hits,
doubly cross-validated against the layered-arithmetic verdict below.
\end{enumerate}

\begin{remark}[a registered prediction that must not be read as a
measurement]\label{rem:claregistered}
An earlier registration predicted for $K_8@28$ ``surviving regimes only
$\{q=0,q=7\}$''; the phrase is a registered prediction, not a measured
residual. The actual census found zero surviving regimes on every line,
and the registered phrase must not be quoted as the census outcome.
\end{remark}

\subsubsection{The characterization candidate and the propagation
bound}\label{subsubsec:echarcand}

The four layers above are exactly what the following candidate predicts.
In its literal form: for every $k\ge1$,
\begin{equation}\label{eq:claechar}
\mathcal{E}_k \;=\; \{\,K_n\ :\ n\ge6,\ \tbinom n2=k\,\},
\end{equation}
equivalently the three-part decomposition: (i) $K_n@\binom n2$ is not
sharp for all $n\ge6$ (sharp at $n=3,4,5$, making $n\ge6$ the exact lower
bound); (ii) $\mathcal{E}_k$ is empty off the triangular layers; (iii) on
triangular layers the exception is exactly the complete graph.

\begin{remark}[registry-layer status of \eqref{eq:claechar}; travels with
any citation]\label{rem:claregistry}
Three registration-layer facts bind every use of \eqref{eq:claechar}.
\emph{First}, two competing candidate shapes are on record --- the literal
shape \eqref{eq:claechar} and the structural shape
``$\mathcal{E}_k$ = the types containing a $K_v$ component'' --- and the
two have never been adjudicated against each other. \emph{Second}, the
literal uniqueness side (iii) stands in tension with the propagation
theorem (a super-type is sharp only if each of its components is
placeable; contrapositively, an unplaceable component makes every
\emph{same-layer} type containing it unsharp): since $K_7$ is unsharp at
layer $21$, at $k=21$ the literal shape would force
$|\mathcal{E}_{21}|=1$, while the structural shape predicts a larger set.
An earlier registration claimed the stronger quantitative form
$|\mathcal{E}_{21}|\ge740{,}226$ via the map $T\mapsto K_7\sqcup T$ on the
$14$-edge types; that claim was withdrawn on review --- $K_7\sqcup T$ has
$21+14=35$ edges and is \emph{not} a layer-$21$ type, so the construction
has no objects in its claimed quantifier domain (and cross-layer
inheritance of unplaceability has known counterexamples in the positive
direction). The honest paper-level reading is
$|\mathcal{E}_{21}|\ge1$ (the type $K_7$ itself, under both shapes), and
uniqueness of the minimal core survives as an open question for $k\ge21$.
\emph{Third}, the candidate
carries mixed candidate/profile load and is not independently audited;
its closed faces are: side (ii) to $k\le16$ (including the $k=17$
full-type sweep of $53{,}394{,}755$ types, every type sharp, tier
machine universal-closure candidate, not independently audited), side (i)
to $n=13$ (tower plus census at $n=6..9$ and the layered-arithmetic
deaths of $K_{11}@55$, $K_{12}@66$, $K_{13}@78$), and side (iii) only at
$k=15$; cumulative over $k=2..17$, $69{,}698{,}732$ types were scanned
with exactly one counterexample ($K_6@15$). Open faces: layers $k\ge18$
(other than $17$ itself), tower layers $n\ge14$, uniqueness of the
minimal core for $k\ge21$, and the prime-family infinity segment
(Section~\ref{subsubsec:primefam}).
\end{remark}

\subsubsection{The prime-family theorem at its non-infinity scope}
\label{subsubsec:primefam}

\begin{theorem}[prime family, non-infinity subdomain]\label{thm:claprimefam}
Let $n=p+1$ with $p\ge7$ a prime congruent to $3 \bmod 4$, and let
$k=\binom n2$ (which is even, divisible by $p$, with $p\mid n-1$). Then
every \emph{non-infinity} regime of $K_n@k$ is dead on paper: the
arithmetic sieve forces $q\equiv eo\equiv0\pmod p$, leaving exactly the
allowed regimes $\{(n,0),(1,p),(p,1)\}$, and each of the three dies by
the order-$2$ criterion, whose feasibility equation reduces to ``$n$ must
be a perfect square'' --- contradicting $p\ge7$.
\end{theorem}

\begin{proof}
Write $e$ and $o$ for the numbers of even- and odd-degree placed vertices,
$q=eo$ for the number of odd--odd pairs, and $\Sigma u$ for the sum of the
half-values $u_i=(y_i\pm k)/2$-type representatives of the finite placed
vertices. Two identities hold for any placement: the total class sum
reads $k(k-1)/2+k\cdot t$ with $t$ the number of half-sums in the upper
half-range, and twice the sum of the half-sums over all pairs equals
$(n-1)\Sigma u-q$ (each odd--odd pair loses one in the halving). Since the
class multiset must be exactly $\mathbb{Z}_k$, summing the classes gives
$(n-1)\Sigma u=k(k-1)+2kt+q$. Reducing mod a prime $p$ with $p\mid k$ and
$p\mid(n-1)$ annihilates the first two terms, forcing $q\equiv0\pmod p$. With
$e+o=n-1=p+1$ we get $q=eo\equiv e(1-e)\pmod p$; since $p$ is prime,
$p\mid e$ or $p\mid e-1$, and $1\le e\le p$ forces $e\in\{1,p\}$ (the
regimes $(1,p),(p,1)$; the case $e=p+1$ i.e.\ $o=0$ is the regime
$(n,0)$, where $q=0$ trivially). These are exactly the three allowed
regimes.

Now $k$ is even, so the order-$2$ (second-character) criterion applies:
writing $x_E=\sum(-1)^{u_i}$ over even points and
$x_O=\sum(-1)^{u_j}$ over odd points, feasibility requires
$x_E^2+2x_Ex_O-x_O^2=e-o$, equivalently
$2x_E^2-(x_E-x_O)^2=e-o$, with $|x_E|\le e$, $|x_O|\le o$,
$x_E\equiv e\pmod2$, $x_O\equiv o\pmod2$.
\emph{Regime $(n,0)$}: $o=0$, so $x_O=0$ and $x_E^2=n$: $n$ must be a
perfect square. \emph{Regime $(1,p)$}: $x_E=1$ (from $x_E\equiv1$,
$|x_E|\le1$), so $(x_O-1)^2=p+1=n$: again $n$ must be a perfect square.
\emph{Regime $(p,1)$}: $x_O=\pm1$, so $(x_E+x_O)^2=n$: the same
conclusion. For $p\ge7$, $n=p+1$ lies strictly between $3^2=9$'s
neighbours $(\lceil\sqrt{8}\rceil)^2$ and the next squares: concretely
$(p-1)/2+1$ is not an integer square root situation --- $p+1$ is even and
$(\tfrac p2)^2<p+1<(p+2)^2$ with no square in between only for
$p\in\{3\}$; since $p\ge7$, $n$ is not a perfect square (if
$n=(m-1)(m+1)$ with $m\ge2$, then $p=(m-1)(m+1)$ is divisible by both
$m-1$ and $m+1$, so $m=2$ and $p=3$, contradiction). All three allowed
regimes die.
\end{proof}

\begin{remark}[scope of Theorem~\ref{thm:claprimefam}; binding]\label{rem:claprimscope}
Three boundary sentences travel with Theorem~\ref{thm:claprimefam}.
\emph{First}, it covers the non-infinity regimes only: the infinity
branch of the prime family is \emph{open} --- there is no paper proof that
infinity-mode cannot be sharp for $n=p+1$ with $p\equiv3\pmod4$, $p\ge7$
(the machine route through the layered arithmetic is blocked at the first
prime layer $d=p$ for $p\ge19$, where the fresh-enumeration volume
$\binom{3p-1}{2p-1}$ exceeds executability; $p=7$ and $p=11$ have been
completed, the latter closing the infinity branch of $K_{12}@66$). The
tower's infinity-mode deaths cover only $n=5..13$. Consequently the
extrapolations registered in the source material --- ``$\mathcal{E}$
contains the infinite family $\{K_{p+1}\}$'' and ``seed propagation
across $n=6..12$'' --- are \emph{withdrawn}: the theorem gives the
non-infinity subdomain only, is a sufficient-family statement and not a
classification, and no infinite-family lower bound for $\mathcal{E}$
follows from it until the infinity branch is closed. \emph{Second}, the
full per-$n$ verdicts require the infinity branch separately even where
the non-infinity side is dead ($n=8$ is closed by census; $n=12$ by the
completed layered enumeration; the remaining prime-family members are
open). \emph{Third}, the machine companion table for $n\le200$ lists $28$
full-death points, of which $23$ have $n-1$ prime (matching the theorem)
and $5$ are composite points $\{112,116,124,160,188\}$ --- the theorem is
a sufficient family, not a classification of the death set.
\end{remark}

\subsubsection{The corrected Fourier identity of the regime-general
character equation}\label{subsubsec:fourier}

The layered arithmetic used above has a closed-form face via the
characters of $\mathbb{Z}_k$. For a nontrivial character $\psi$ of
$\mathbb{Z}_k$, write, for a placement $S$ with finite points parametrized
as $y=2v$ (even) and $y=2w+1$ (odd),
\[
\Psi_E=\textstyle\sum_i \psi(v_i),\quad
\Psi_{E,2}=\textstyle\sum_i \psi(2v_i),\quad
\Psi_O=\textstyle\sum_a \psi(w_a),\quad
\Psi_{O,2}=\textstyle\sum_a \psi(2w_a).
\]

\begin{theorem}[corrected Fourier identity]\label{thm:clafourier}
Let $\psi$ be a nontrivial character of $\mathbb{Z}_k$.
\begin{enumerate}[label={\rm(\roman*)},leftmargin=2em]
\item (non-infinity placements) If a placement is sharp, then
\[
\tfrac12\bigl(\Psi_E^2-\Psi_{E,2}\bigr)
+\psi(1)\cdot\tfrac12\bigl(\Psi_O^2-\Psi_{O,2}\bigr)
+\Psi_E\Psi_O \;=\;0 .
\]
\item (infinity placements) If a placement with $\infty\in S$ is sharp,
then the infinity-class contribution is
\begin{equation}\label{eq:clarho}
\rho_\psi\;=\;\sum_i \psi(y_i \bmod k)\;=\;\Psi_{E,2}\;+\;\psi(1)\cdot\Psi_{O,2},
\end{equation}
i.e.\ both sides are doubled: the odd-side contribution is the doubled
half-value sum $\Psi_{O,2}=\sum_a\psi(2w_a)$, not $\Psi_O$. Equivalently,
the infinity-mode equation takes the ``$i\le j$'' form
\[
\tfrac12\bigl(\Psi_E^2+\Psi_{E,2}\bigr)
+\psi(1)\cdot\tfrac12\bigl(\Psi_O^2+\Psi_{O,2}\bigr)
+\Psi_E\Psi_O \;=\;0 .
\]
\end{enumerate}
\end{theorem}

\begin{proof}
(i) The classes of even--even pairs are
$\mathrm{class}(\{2v_i,2v_j\})=v_i+v_j$, of odd--odd pairs
$\mathrm{class}(\{2w_a+1,2w_b+1\})=w_a+w_b+1$, and of mixed pairs
$\mathrm{class}(\{2v_i,2w_a+1\})=v_i+w_a$ \eqref{eq:claclass}. Summing
$\psi$ over all unordered pairs with distinct classes, using the
algebraic identity
$\sum_{i\ne j}\psi(x_i+x_j)=\Psi^2-\Psi_{2}$ folded to $i<j$ (which
produces the $-\Psi_{E,2}$ resp.\ $-\Psi_{O,2}$ halves and the factor
$\psi(1)$ on the odd side from the $+1$ shift), the sharpness hypothesis
(class multiset $=\mathbb{Z}_k$ exactly) makes the total vanish.

(ii) The infinity edges $\{\infty,y\}$ have classes $y\bmod k$, so their
character contribution is $\sum_i\psi(y_i\bmod k)$. For an even point
$y=2v$ this is $\psi(2v)$; for an odd point $y=2w+1$ it is
$\psi(2w+1)=\psi(2w)\psi(1)$. Summing gives
$\rho_\psi=\Psi_{E,2}+\psi(1)\Psi_{O,2}$ \eqref{eq:clarho}. Substituting
into (i) and merging the doubled contributions --- each half-value sum now
appears both through its square and through its doubled self-sum --- turns
the ``$i<j$'' form into the ``$i\le j$'' form stated.
\end{proof}

\begin{remark}[why the corrected form is quoted; the falsified
predecessor]\label{rem:clafourierfix}
An earlier draft of the identity stated the infinity contribution as
$\rho_\psi=\Psi_{E,2}+\psi(1)\Psi_O$, omitting the odd-side double angle.
That form is false: on the $K_3$ configuration $\{\infty,1,3\}$ at $k=3$
with $\psi(a)=\omega^a$, $\omega^3=1$, $\omega\ne1$ (a sharp placement,
the classes being $1,0,2$), the original form gives the residual
$\omega+2\omega^2=-2-\omega\neq0$, while the corrected form gives $0$.
This was verified independently in exact $\mathbb{Z}[\omega]$ integer
arithmetic, and the corrected identity \eqref{eq:clarho} was verified
identically on all infinity placements at $k=3,6,10$ ($15+220+4{,}845$
placements; the original form fails on $8$, $178$ and $4{,}514$ of them
respectively). Two accompanying sentences bind every citation of the
layered-arithmetic machine results: \emph{all machine results (the tower
counting rules and the censuses) are independent of the character-equa
formula} --- the machine enumeration counts infinity classes by their
actual residues $2a$ and $2a+1$, which is the corrected bookkeeping ---
and the erroneous formula, wherever it existed, appeared only in the
internal source report and in no form in the frozen v2.5 submission.
\end{remark}

\subsubsection{The revival-depth law and the falsification of the
forall-\texorpdfstring{$k$}{k} target}\label{subsubsec:revdepth}

We now turn to the polar-corner family of the placeability programme. For
$k\ge5$ let $T_k=(k-1)\cdot K_2$ (the \emph{polar-corner family}): the
placement engine places the $k-1$ edges one by one, a \emph{state}
$\sigma=(U,C)$ at slot $j$ has $j$ components placed with occupied points
$U$ ($|U|=2j$) and occupied classes $C$ ($|C|=j$), free points
$F=\mathbb{Z}_{2k}\setminus U$, free classes $C_f=\mathbb{Z}_k\setminus C$,
and $\rho:=m-j=(k-1)-j$ unplaced components. The state is \emph{dead} if
$\mathrm{Span}(F)\cap C_f=\emptyset$ (no further $K_2$ can be placed); a
\emph{$d$-revival} of a dead state is the retraction of some $d$ placed
components (a set $W\subseteq U$ of points with $|W|=2d$ and a set
$Y\subseteq C$ of classes with $|Y|=d$) after which the remaining
$\rho+d$ components can all be placed legally; $d^*$ is the minimum $d$
admitting a revival.

\begin{theorem}[revival-depth lower bound]\label{thm:clarevdepth}
Let $T_k=(k-1)\cdot K_2$ and let $\sigma=(U,C)$ be a dead state at slot
$j$ with $\rho=(k-1)-j$ unplaced components. If $\sigma$ has a
$d$-revival, then $\rho+d\le3d$, i.e.\ $\rho\le2d$, i.e.
$d^*\ge\lceil\rho/2\rceil$.
\end{theorem}

\begin{proof}
Five steps, using only the placement semantics and the death identity.

(1) A $d$-revival retracts $d$ placed components; since every component of
$T_k$ is a $K_2$, the retracted point set $W$ has $|W|=2d$, and since each
placed component uses exactly one class, the retracted class set $Y$ has
$|Y|=d$.

(2) After the retraction, the $\rho+d$ components of the suffix must be
placed on the points $F\cup W$ using pairwise-distinct classes from
$C_f\cup Y$; the placements are pairwise point-disjoint with pairwise
distinct classes, i.e.\ they form a rainbow matching of size $\rho+d$ on
the board $(F\cup W,\,C_f\cup Y)$.

(3) Every edge $e$ of this matching either lies entirely in $F$ (an
``$F$-internal'' edge) or meets $W$; this dichotomy is exhaustive.

(4) An $F$-internal edge has its class in $\mathrm{Span}(F)$ (both
endpoints in $F$); but the death identity gives
$\mathrm{Span}(F)\cap C_f=\emptyset$, so its class lies in $Y$. Distinct
matching edges have distinct classes, so the number of $F$-internal edges
is at most $|Y|=d$.

(5) An edge meeting $W$ consumes at least one point of $W$, and matching
edges are disjoint, so their number is at most $|W|=2d$.

Steps (3)--(5) cover all edges of the matching, so
$\rho+d\le d+2d=3d$, i.e.\ $\rho\le2d$.
\end{proof}

\begin{remark}[the companion machine law and the paper-ization gap]
\label{rem:cladstarlaw}
The reverse inequality $d^*\le\lceil\rho/2\rceil$ is a \emph{machine law,
not a theorem}: it holds with zero violations over $4{,}273$ case slots
($4{,}252$ pool dead ends plus $21$ polar minima; deduplicated $4{,}268$),
but no paper proof of the upper-bound direction is known, and the
paper-ization of this law is registered as an open gap. The companion
deficit law $\min\text{-}\mathrm{deficit}(d)=\max(0,\rho-2d)$ is likewise
a machine law.

Consequence (falsification): the long-standing target statements $K_1$ and
the invariant refinement $[\mathrm{INV}\text{-}\mathrm{BT}''\text{-}2]$ in
$\forall k$ form are \emph{false}. The polar-corner minima at
$k=23,\dots,29$ have $\rho=5$, so by Theorem~\ref{thm:clarevdepth} no
$2$-revival exists; the machine channels agree (slot-pair channels $0$
out of $136$ and $0$ out of $153$ candidates, full $(W,Y)$-space sampling
$0$ out of $3{,}000$), and the former $60$-million-node ``double OPEN''
search walls at $k=23/24$ are structural, decided false in constant time
by the four-layer pruning of which Theorem~\ref{thm:clarevdepth} is the
last layer, with $d^*=3$ witnesses verified. The replacement targets are
the stratified statements: $d^*=\lceil\rho/2\rceil$ on the polar family
(lower bound theorem, upper bound machine law).
\end{remark}

\subsubsection{The criterion-level results: stratified statements only}
\label{subsubsec:criteria}

The remaining material of this section concerns criterion-level results:
statements about necessary conditions, their completeness on tested
domains, and their stratified theorem-form closures. A recurring discipline
binds all of them, which we state once here: \emph{a min-val$\le1$-type
criterion verified by per-instance machine scans over finitely many $N$ is
not an existence proof for general $N$}, and no sentence below may be read
as one.

\paragraph{The closed-form criterion family for the packing residual.}
For the $\mathbb{Z}_7$ packing problem at the top layer $\Sigma d=|S|$,
define the power-sum moment conditions $P_T$ ($\bmod\ 7$) and their
integer upgrades $\widetilde P_T$: there exist $C\subseteq S$ with
$|C|=\Sigma d$ and pairwise disjoint key sets $K_i\subseteq
\mathbb{Z}_7\setminus\{0\}$, $|K_i|=d_i$, matching all power sums up to
degree $T$. Necessity is immediate (a feasible packing makes the two sides
the same multiset), and at the top layer with $m=\Sigma d=|S|\le6$,
sufficiency of $P_m$ follows from Newton's identities ($t\le6$ invertible
in $\mathbb{F}_7$) forcing equality of characteristic polynomials, hence
of multisets: \emph{completeness at $T=|S|$} is a paper theorem with
machine backing ($3{,}872$ top-layer points, zero counterexamples). On the
census domain ($120$ orbits $\times\,330$ demand vectors $=39{,}600$
points, exhausted), the second integer moment $\widetilde P_2$ and the
third mod-$7$ moment $P_3$ agree with the oracle on every point ---
$0$ misclassifications out of $39{,}600$, zero false positives on all
$6{,}649$ feasible points --- and the residual core collapses $424\to0$ in
this domain (coverage ladder $1{,}133\to1{,}519\to1{,}557$). Out-of-domain
evidence is a sample, not an enumeration, quoted only at the corrected
scale: $1{,}200{,}000$ draws yield $33{,}018$ valid sample points
($31{,}242$ after deduplication; sampling support $933{,}450$;
$\approx3.5\%$ of support). Two boundary sentences bind any citation: the
$T=3$/$T=2$ closure is a domain-exhaustion statement (the census domain),
not a theorem --- the only paper-level completeness is at $T=|S|\le6$ ---
and the integer moments are not affinely equivariant, so cross-domain
citations use the mod-$7$ version. A convex-unreachability companion
holds in its corrected scope: if an infeasible point is a rational convex
combination of feasible colour-count vectors ($38/38$ explicit
certificates, exactly re-checked), then no criterion that is a relaxation
of the natural fractional LP can exclude it --- its original
``any-rank Chv\'atal--Gomory'' listing clause is refuted by rank-$1$
CG single cuts ($9/38$ of the very certificates) and is never quoted.

\paragraph{The stratified valuation criterion.}
The valuation criterion in its aligned family reads: \emph{for every
instance with the three anchor offsets on a common coset of the
subgroup generated by the first anchor difference, and with capacity slack
$\beta\ge2$, the minimum of the valuation $\mathrm{val}(C_1,C_2)$ over
admissible pairs is $\le1$}; the machine scan domain is $\beta\ge1$. We
register the criterion and its existence half as formal objects:
\begin{conjecture}[valuation criterion, candidate form]\label{conj:clacand}
Every so1 instance with $\beta\ge2$ has $\min\mathrm{val}\le1$.
\end{conjecture}
The statement is a \emph{candidate}: no paper proof exists, its
falsification clause is an explicit so1 instance with $\beta\ge2$ and
$\min\mathrm{val}>1$, and every machine account of this line is a directed
sweep of exactly that clause, with zero hits on the tested domains. The
general-$N$ existence half is carried separately:
\begin{conjecture}[the residual existence gap]\label{int:conj-gapex}
For every $N$ and every so1 instance with $\beta\ge2$,
$\min\mathrm{val}\le1$.
\end{conjecture}
The status of Conjecture~\ref{conj:clacand} is stratified, and only this
stratified form may be cited:
\begin{itemize}[leftmargin=1.5em]
\item \emph{Per-$N$ closure in theorem form} for $N\in\{6,8,10,12\}$: a
three-tier finite search, proved complete and terminating on paper (a
witness-bound theorem gives the per-instance bound $B(I)=d_2+1$ and a
uniform $B(N)=N-2$; a completeness theorem shows the three searched
families cover every witness), with zero residue on
$588/30{,}760/634{,}960/7{,}798{,}903$ instances respectively.
\item \emph{The general-$N$ existence half of
Conjecture~\ref{int:conj-gapex} is open}: whether for all $N$
and all instances with $\beta\ge2$ an admissible pair with
$\mathrm{val}\le1$ exists. The $C_1$ side is closed for all $N$ by a
constructive one-leak theorem --- for every instance and every admissible
$d_1$ there is $C_1$ with at most one leak, by an explicit three-case
construction over the cycle decomposition (full-cycles, partial cycles,
and the cycle containing the first anchor) --- so the obstruction is
entirely on the $C_2$ side (reachability of a near-closed board under the
capacity identity). The phrase ``theorem-level closure of the criterion''
is an explicit non-claim, including at larger $N$.
\end{itemize}

\paragraph{The covering identity as a candidate.}
The strongest structural statement of the valuation line is the following
six-family identity, quoted strictly at candidate tier.

\begin{proposition}[covering identity; paper candidate assembled and
independently re-audited; not mechanized]\label{prop:clacover}
For every instance $I=(N,P,T,d)$ with $P\subseteq\mathbb{Z}_N$,
$|P|\ge4$, $T=(t_1,t_2,t_3)$ pairwise distinct,
$d\in\mathbb{Z}_{\ge0}^3$, $\Sigma d_i=|P|-1$,
$d_i\le|P|-[t_i\in P]$: let $W=P\setminus\{t_1\}$,
$k=d_2+d_3+[t_1\notin P]$ (so $|W|-k=d_1$),
$\delta_{12}=(t_1-t_2)\bmod N$, $\delta_{13}=(t_1-t_3)\bmod N$, and
$L=\{R\subseteq W:\ |R|=k\}$. Let $\mathrm{Im}$ be the set of position-sets
produced by the eight-branch family generator, and let
$C_{12},C_{13},\mathrm{RX},O3_{12},O3_{13},V_7$ be the six covering
families defined via the per-orbit run statistic $\mathrm{rc}_\delta$
(full containment counts $1$), cross-reachability, full-orbit removal, and
the image of the legacy candidates routine. Then
\[
\mathrm{Im}\cap L \;=\; (V_7\cap L)\cup(C_{12}\cap L)\cup(C_{13}\cap L)
\cup(\mathrm{RX}\cap L)\cup(O3_{12}\cap L)\cup(O3_{13}\cap L).
\]
\emph{Corollary (boundary).} $R\in L\setminus\mathrm{Im}$ implies
$\mathrm{rc}_{12}(R)\ge3$, $\mathrm{rc}_{13}(R)\ge3$, not
cross-reachable, not $O3_{12}$, not $O3_{13}$, and $R\notin V_7$.
\end{proposition}

\begin{remark}[tier and riders of Proposition~\ref{prop:clacover}]
\label{rem:clacoverriders}
The identity is a paper candidate (both directions have paper proofs and
the package passed an independent from-scratch re-derivation audit); it is
\emph{not} a mechanized or formalized theorem, and \emph{no disjointness
of the six families is claimed}. The machine-checkable face is the
five-family identity without the $V_7$ term: all machine checks are run in
the \emph{no-$V_7$ reading} (the $V_7$ branch omitted, giving a
conservative superset $\mathrm{fam}_{\mathrm{no}V_7}\subseteq
\mathrm{Im}\cap L$, a one-way implication toward the true face), and the
$V_7$ branch's machine face is frozen --- the legacy candidate-generation
source files were lost with no recoverable copy --- while the $V_7$ term
itself is definitional. The residual $V_7$-only mass is nonzero and
precisely measured: $340$ positions at $N=8$ ($1{,}266$ superset-uncovered
minus $926$ true-face uncovered, both sides measured) and $2{,}267{,}655$
instance-weighted at $N=14$ ($52{,}572{,}256$ minus $50{,}304{,}601$).
The supporting lemma list ($L1$--$L9$) is paper-proved; $L1/L9$ are
definitional ($V_7$ outputs), and the count-difference identity inside
$L8$ counts whole \emph{orbits}, not points (a point-count reading
mismatches in $136/936$ enumerated configurations). One reverse clause of
the run-statistic calculus must be stated one-way: ``$\mathrm{rc}=2$
implies a union of two disjoint cyclic pieces; the converse only gives
$\mathrm{rc}\le2$''. The three micro-lemmas (single interval
$\Rightarrow\mathrm{rc}=1$; two disjoint intervals
$\Rightarrow\mathrm{rc}\le2$; subadditivity) carry paper proofs plus two
rounds of randomized machine checks with zero violations ($12{,}558 /
2{,}500 / 12{,}000$ first pass; $4{,}000/1{,}109/4{,}000$ audit pass).
\end{remark}

\paragraph{The falsification-scale facts of the exchange ledger.}
The single-swap ledger statement --- that an uncovered position admits an
in-family companion at swap distance $2$ --- does \emph{not} hold
universally. On the fixed random pool ($900$ instances, $300$ each at
$N\in\{10,12,14\}$; $1{,}501$ uncovered-witness positions in the no-$V_7$
superset reading): a distance-$2$ companion in the family image exists for
$1{,}469/1{,}501$ ($97.87\%$) of witnesses, and within the five-family
domain for $1{,}447/1{,}501$ ($96.40\%$). Failing witnesses: $54/1{,}501$
($3.60\%$), concentrated in $12/900$ instances ($1.33\%$); all $12$ lie in
$N\in\{12,14\}$ with $|P|$-spectrum
$\{5,6,9,9,10,11,11,12,12,12,12,14\}$ --- $10$ of $12$ biased to high
$|P|$, $2$ outside the dense band. The cause split of the $54$ is
$32$ structural-side (no in-family companion exists) against $22$
ledger-side (companions exist but none satisfies the ledger constraints).
The distance-$\le4$ binary form also fails on the same pool: $26/54$ have
a distance-$4$ five-family companion, $28/54$ do not. Two accounting
disciplines bind every quotation: the instance-level ($12/900$) and
witness-level ($54/1{,}501$) figures are different measurement layers and
are never interchanged, and the $96.40\%$ figure is a
no-$V_7$-superset-reading statistic not directly comparable to
true-reading protocols.

\begin{remark}[the bridge guards]\label{rem:clabridge}
The bridge between the $\mu$-sufficiency question (whether satisfaction of
all $\mu$-inequalities implies feasibility on the difference board) and
the valuation criterion carries three explicit guards, all of which must
accompany any quotation. \emph{(a)} The equivalence ``abstract theorem
$\iff$ $\mu$-sufficiency on the difference board'' holds only under the
$(H2)/n\le p-1$ guard; the unguarded form is falsified by an explicit
$\mathbb{Z}_{14}$ counterexample ($P=\{0,1,2\}$, $T=(0,1)$, $d=(1,2)$:
all $\mu$-inequalities hold, the board is infeasible). \emph{(b)} The
identification of the abstract existence gap with the general-$N$
existence half above is valid only as a $k=3$-layer monotone-reduction
pointer; as a literal identity it fails on four axes (the $\beta$-band,
the so-slice, the $k$-domain, and the closure grade). \emph{(c)} The
$k\ge4$ layer is not touched by the criterion at all --- no criterion
exists in the current language for $k\ge4$, a language-level gap rather
than an unexecuted proof. Consequently: a min-val$\le1$-type criterion
verified by per-instance machine scans over finitely many $N$ is not an
existence proof for general $N$, and $\mu$-inequality checks over
finitely many $N$ do not constitute an existence theorem. Necessity of
the $\mu$-inequalities is strictly proved (feasible placements satisfy
all of them), the sufficiency direction is false on general bipartite
boards (an explicit $2\times2$ counterexample whose colouring violates
the difference rule), and sufficiency on the difference board is open
(zero violations in $N=6..14$ machine accounts; no general-$N$ paper
proof). Finally, the $\mu$ $k=2$ stratified existence question --- the
main open wall of this line --- is open in its paper face (machine face
closed on seven domains for the $k=2$ layer; recorded obstruction: no
base-to-step mechanism), and it has \emph{no logical dependency} on the
family-image characterization: characterizing the family image is not the
realize-layer existence.
\end{remark}

\begin{remark}[what the criterion layer does and does not give]
\label{rem:clacriterionlimits}
The criterion-level results of this subsection share one logical shape:
each is a necessary condition, or a family of necessary conditions, whose
completeness is established on explicitly enumerated domains, or whose
per-$N$ verification is theorem-form on the listed $N$. None of them
yields, unconditionally and for general parameters, the existence of the
object the criterion is meant to witness: the valuation criterion does
not give feasible pairs for general $N$ (the existence half is open), the
moment criteria do not give packings outside the census domain (and the
top-layer completeness is bounded at $|S|\le6$), and the covering identity
characterizes an image, not a realizability theorem. In particular the
conclusion of the conditional chain consuming these criteria remains
conditional: no unconditional existence statement, and hence no
unconditional progress on the underlying decomposition conjecture, is
claimed anywhere in this section.
\end{remark}

% ============================================================================
% sec-computation.tex — Computational methodology: exhaustive accounts,
% audits, and Lean trust boundaries. Label prefix: com: (all labels owned
% by this file).
% ============================================================================

\section{Computational methodology: exhaustive accounts, audits, and trust
boundaries}\label{sec:computation}

Every machine-backed statement in this paper is quoted at a measured
strength, and this section documents the apparatus that produces the
measurements. There are four layers. First, the exact-value pipeline behind
the constructive parts (Section~\ref{com:subsec-pipeline}), together with a
self-contained recheck script released with the paper
(Section~\ref{com:subsec-recheck}). Second, the re-verification account of
the conditional chain's machine layer
(Section~\ref{com:subsec-kline}). Third, the Lean development: what is
formalized, and where its trust boundaries lie
(Section~\ref{com:subsec-lean}). Fourth, the key-table ledger behind the
covering identity (Proposition~\ref{prop:clacover}) and the independent
audit it passed (Section~\ref{com:subsec-ledger}). Formalization status is
itemized declaration by declaration in Section~\ref{sec:open}.

\subsection{The exhaustive construction pipeline}\label{com:subsec-pipeline}

All exact values and all decompositions cited in the constructive parts were
produced and re-checked by one pipeline: a bitmask exact dynamic program
over edge subsets computes exact values of $\pn(G)$ (and of the endpoint
budget $\beta(G,x)$ of Definition~\ref{int:def-budget}, with an endpoint
constraint); a backtracking exact-cover solver with longest-path-first
enumeration handles larger instances; and an independent
\texttt{verify\_decomposition} routine asserts that a witness consists of
simple paths exactly partitioning the edge set. Randomized steps use fixed
seeds, and every number is reproducible from the scripts and certificate
files released with the paper. In the per-theorem verification table
(Appendix~\ref{com:app-comp}), ``exact'' means computed by an exact algorithm
(dynamic program or exhaustive cover), ``witness'' means a machine-verified
decomposition, and ``census'' means that all graphs of the stated class were
processed.

Three rows of that table indicate its granularity.

\begin{itemize}[leftmargin=1.8em]
\item The odd-column value $\pn(K_{4,n})=(n+1)/2$
(Proposition~\ref{prop:conk4large}) was re-derived independently for every
$n=3,\dots,31$ by a structured exact search over $K_{4,n}$, recovering
$(n+1)/2$ each time; the closed-form construction was verified for
$k=6,\dots,301$ (i.e.\ $n=13,\dots,603$), with witness lists
machine-verified. The printed small-case witnesses of
Proposition~\ref{prop:conk4small} are machine-verified by the same pipeline
(edge-exact partition, simplicity).
\item The witnesses behind Corollary~\ref{cor:concol6} of the $n_2=6$ column
($K_{7,6}$, $K_{9,6}$, $K_{11,6}$, $K_{13,6}$) were each verified by two
independent implementations (the pipeline verifier and an adjacency-matrix
checker); the ring construction was re-derived for $k=7..31$ from the base
$k_0=3$, with overlapping chains from $k_0=4,5$.
\item The negatives on the five open bases
(Observation~\ref{obs:concostwall}) are \emph{truncations, not exclusions}:
more than $14$ million search nodes (column $n_2=8$) and about $0.9$ million
nodes (column $n_2=10$) were spent over every ledger-feasible deficit
multiset, without a hit and without exhausting a single schedule. The
conjectural status of the five bases is therefore unchanged by these
searches, in either direction.
\end{itemize}

\subsection{A self-contained independent recheck script}\label{com:subsec-recheck}

The repository releases a recheck script, \texttt{checks.py}, written from
scratch during the independent review of the paper: it imports none of the
project's research code, consists of ten check functions built from
assertions only, writes its results to \texttt{check\_results.json}, and
exits nonzero if any assertion fails. Running
\texttt{python3 checks.py} from the repository root executes all ten checks
in about three seconds (wall clock on a laptop-class machine; measured
$3.1$\,s for the release version). The coverage is:

\begin{itemize}[leftmargin=1.8em]
\item \emph{Ring closure} (Theorem~\ref{thm:conringclosure}): all $280$
parameter pairs ($h=1,\dots,14$; per $h$ the $20$ even ring sizes $m$ from
$2h+2$ to $2h+41$), checking path simplicity and that the paths partition
the edge set exactly once.
\item \emph{The $K_{4,n}$ construction} (Proposition~\ref{prop:conk4large}):
all $95$ parameter values $k=6,\dots,100$, asserting $k+1$ paths, simplicity,
and the exact $A\times B$ edge partition with multiplicity $1$.
\item \emph{The odd-clique local engine} (Theorem~\ref{thm:conL0}): $61$ odd
values $j=1,3,\dots,121$, checking that the removed arcs form a Hamilton
path, that the remaining $j$ paths have the prescribed lengths, and that the
partition is exact. Scope: odd $j$ only; even $j$ is not verified by the
script (it is open; Remark~\ref{rem:conchaincomp}).
\item \emph{The non-infinity prime-family arithmetic}
(Theorem~\ref{thm:claprimefam}): $86$ primes $p\equiv3\pmod4$ up to $991$,
checking the sieve's allowed regimes and the perfect-square obstruction. This
is a finite recheck of the elementary steps; it does not substitute for the
paper proof and covers the non-infinity regimes only.
\item \emph{The $K_9@36$ census ledger}: the six recorded regime lines'
configuration counts against their closed combinatorial forms, the histogram
sums, the zero recorded hits, and the total-domain arithmetic
$18{,}385{,}772{,}680=\binom{71}{8}+\binom{36}{8}\cdot2^8$. This is ledger
and artifact arithmetic; the census itself is not rerun.
\item \emph{The non-infinity death table}: $28$ full-death points for
$n\le200$, of which $23$ have $n-1$ prime (matching
Theorem~\ref{thm:claprimefam}) and the composite points are
$\{112,116,124,160,188\}$. Again artifact accounting, not a rerun of the
table's feasibility searches.
\item \emph{The corrected Fourier identity} (Theorem~\ref{thm:clafourier}):
the $K_3$ configuration $\{\infty,1,3\}$ at $k=3$ with $\psi(a)=\omega^a$,
reproducing both the failure of the superseded form (residual $-2-\omega$)
and the vanishing of the corrected form \eqref{eq:clarho}, in exact
$\mathbb{Z}[\omega]$ integer arithmetic.
\item \emph{Block versus global degree}
(Section~\ref{int:subsec-blockvsglobal}): the $K_{1,5}$-with-pendant-leaves
witness, confirming that $\EV(G)\cong K_{1,5}$ has a vertex of E-degree $5$
while every block of $\EV(G)$ is a single edge (maximum degree $1$).
\item \emph{The Lean placeholder scan}: a comment-stripped scan of the Lean
sources for \texttt{sorry}, \texttt{admit} and \texttt{axiom}
declarations, finding exactly five occurrences, listed in
Section~\ref{com:subsec-lean}.
\item \emph{Snapshot integrity}: SHA-256 verification of the archived source
snapshot (20 files plus the pinned artifacts). In the review environment the
same check additionally reported the six post-review-edit live files (the
expected repair edits to the working notes and four section files, not frozen
deliverables); that comparison required the live working tree, which is not
part of the released repository, so the released script reports the field as
empty.
\end{itemize}

One reading rule binds the released result file: a PASS records that the
intended recheck succeeded, and for two of the ten checks success means
\emph{reproducing a recorded error or distinction} (the falsified Fourier
form; the block-versus-global-degree witness). A PASS is an execution
record, not a certificate of global correctness.

\paragraph{Census methodology note ($K_9@36$).} The largest
configuration-level account behind the classification programme is
conducted line by line: for each regime line of $K_9@36$, the
configurations are enumerated in full and summarized by a closed-form
count checked against the enumeration, with a histogram-sum completeness
assertion per line. The nine regime lines sum to
$\binom{71}{8}=10{,}639{,}125{,}640$ (the Vandermonde identity is
machine-checked against the per-line counts); the infinity line,
$\binom{36}{8}\cdot2^8=7{,}746{,}647{,}040$, is covered in four shards of
$1{,}936{,}661{,}760$ each under an exact non-overlapping-partition
assertion; every line has zero hits; and the total domain is
$18{,}385{,}772{,}680$ configurations. The released checks verify the
ledger arithmetic of this account (per-line closed forms, histogram sums,
recorded hits, and the total); the census itself is not rerun by them.
The tier of the verdict --- machine universal-closure candidate, not
independently audited --- is the one recorded in
Section~\ref{subsubsec:firstlayers}, and the layered-arithmetic verdict of
the same section provides a mechanism-level cross-check.

\paragraph{Complete-walk counts on the flicker line.} The GEN existence
programme contributes a counting account of a different kind: exact
enumeration of zero-dodge complete walks on the skeleton template family,
$w_1=r-5$, by depth-first search with node budgets and cap tracking,
cross-checked across independent implementations. On the odd tiers the
exact counts are $158$, $1{,}733$, $11{,}466$, $81{,}660$, $484{,}578$,
$3{,}670{,}846$ for $r=19,21,23,25,27,29$ --- six tiers, never presented
as one closed sequence; the $r=29$ count is exact but
class-bound-conditioned and is quoted only with that caveat. One further
exact single point, $23{,}933{,}468$ at $r=31$, stands apart at
exhaustive grade ($j=0$: $22{,}607{,}250$; $j=1$: $1{,}326{,}218$). At
$r=33$ there is no count: the search volume is a quantified hard wall,
and a segmented long run reports none. The $w_1=r-4$ side is zero for all
$r\ge19$ and the even-$r$ side with $w_1=r-5$ is zero by the parity
theorem --- both theorem-level, with the four theorem-relevant cells
($r\in\{20,22\}$, both $w_1$) genuinely exhausted at zero; the theorem
clears the zero-dodge side only, complete walks with a positive
zero-defect existing at $r=20,22$ in the record and remaining untouched.

\subsection{The re-verification account of the conditional chain's machine
layer}\label{com:subsec-kline}

The conditional chain of Section~\ref{int:subsec-budget} (Theorem~\ref{int:thm-ledger})
was subjected to an independent
re-derivation review: a from-scratch re-derivation of its load-bearing
implications, conducted before reading the recorded proofs and then
reconciled against them, accompanied by fresh machine checks re-implemented
from the definitions.

\paragraph{Implications.} Of the $34$ load-bearing implications of the chain,
$31$ were independently re-derived as valid. The remaining $3$ are
wording- and registration-level only: an overstated equivalence in the
density remark of the chain (the unconditional direction is one-way), an
under-registered residue inside the degree-two layer of the budget
conjecture (recorded in Section~\ref{sta:subsec-limitations}), and one
stale supporting statistic. None reverses an implication direction, and no
mathematical hole was found in any implication.

\paragraph{Machine checks, at the downgraded caliber.} Six independent
machine-check groups accompanied the review. \emph{Five of the six} closed
with zero failures (volumes: $180{,}070$ construction instances;
$550$ end-to-end pipeline instances; $876$ census graphs; $121$ types of the
placement family; $2{,}342$ buried-vertex instances; $5{,}880$ alignment
instances; plus a dedicated connected-domain cut-vertex check on $106$
instances). The sixth group closed with three of its four blocks at zero
failures and one block emitting $60$ alerts; each of the $60$ was examined
individually and all $60$ are \emph{disconnected} graphs, outside the
connected domain on which the budget conjecture is stated --- zero connected
counterexamples. The account is therefore never ``six groups with zero
failures'': the correct statement is five groups clean, and the sixth clean
up to a fully discharged false alarm caused by a missing connectivity
filter, with the connected domain covered by the dedicated check.

\paragraph{What was \emph{not} re-verified.} The same review registered its
limits, which travel with any citation of it: three combinational notes of
the chain were logic-checked but not separately machine-checked; machine
checks predating the review were not replayed; the reviewers'
derivation-before-reading process is not itself reproducible (only its
products were reconciled against the originals); the full proof
reconstruction of the all-odd branch of Theorem~\ref{int:thm-lovasz}
remains open (Section~\ref{sta:subsec-limitations}); and the domains $n\ge13$
and the general-graph domain beyond $12$ edges at $n\ge9$ were not
enumerated. (One gap registered by the review --- two of the five domains
below not freshly recomputed --- was subsequently closed by a full fresh
recomputation of all five domains, $178{,}860/178{,}860$ instances over
$16$ fields with zero differences.)

\paragraph{The five exhaustively enumerated domains.} The budget conjecture
$K$ (Definition~\ref{int:def-budget}) carries a machine account on five
exhaustively enumerated domains: the atlas of connected graphs on
$n\in\{4,6\}$ with $m\le13$ together with \texttt{geng -c 8 1:12}, and the
connected \texttt{geng} censuses \texttt{8/9/10/11/12 1:12} --- in total
$39{,}027$ graphs, $178{,}860$ even-degree point instances
(a vertex $x$ of even degree in a graph of the domain), of which $43{,}009$
have degree $d\ge4$. The boundary of the domain is part of the account: only
the sparse layers with $m\le12$ are enumerated (at $n=11$ the average degree
is $\le2.18$); the general-graph domain beyond $m\le12$ is not covered; and
odd and even $n$ are accounted separately throughout, never pooled (the
budget ceiling is $\lceil n/2\rceil$ and the analysis differs by parity).

\begin{table}[ht]
\centering\small
\caption{The five exhaustively enumerated domains of the budget conjecture.
Columns: graphs; even-degree point instances; with $d\ge4$; instances dead
at the first two insertion depths of the degree-insertion tool; instances
where the tool certifies failure; instances attaining the ceiling
$\beta=\lceil n/2\rceil$. In every domain: zero violations of
$\beta(G,x)\le\lceil n/2\rceil$, and zero mismatches of the insertion
criterion's equivalence against an independent tool-truth oracle.}
\label{com:tab-kdomains}
\begin{tabular}{lrrrrrr}
\hline
Domain & graphs & instances & $d\ge4$ & two-probe & tool-dead & tight\\
\hline
atlas $n\in\{4,6\}$, $m\le13$, $\cup$ \texttt{geng -c 8} & $2{,}933$ &
$11{,}328$ & $4{,}232$ & $306$ & $136$ & $63$\\
\texttt{geng -c 9} & $7{,}654$ & $33{,}296$ & $10{,}003$ & $227$ & $12$ &
$1{,}420$\\
\texttt{geng -c 10} & $11{,}989$ & $55{,}664$ & $13{,}471$ & $783$ & $60$ &
$72$\\
\texttt{geng -c 11} & $10{,}874$ & $52{,}126$ & $10{,}558$ & $1{,}217$ &
$0$ & $990$\\
\texttt{geng -c 12} & $5{,}577$ & $26{,}446$ & $4{,}745$ & $1{,}102$ & $0$ &
$12$\\
\hline
total & $39{,}027$ & $178{,}860$ & $43{,}009$ & $3{,}635$ & $208$ &
$2{,}557$\\
\hline
\end{tabular}
\end{table}

Two readings of Table~\ref{com:tab-kdomains} are fixed by the account.
First, \emph{tool-dead is not a counterexample}: the degree-insertion
criterion (Section~\ref{sta:subsec-limitations}) characterizes one
construction tool, and all $208$ tool-dead instances still satisfy
$\beta\le\lceil n/2\rceil$ --- indeed the zero-violation verdict of the
conjecture itself reads on all $178{,}860$ instances. Second, the witness
certificates attached to the account total $1{,}279$ ($484$ of $500$ in the
base domain plus $795$ of $800$ across the extended domains, $5$
tool-dead instances skipped by design, zero failures), each an explicit
decomposition verified three ways. The residual proof debts of this layer
are collected in Section~\ref{sta:subsec-limitations}.

\subsection{The Lean development and its trust boundaries}\label{com:subsec-lean}

The Lean development (toolchain \texttt{leanprover/lean4:4.34.0}) certifies
a precise subset of the paper's statements. We describe what is proved, how
each class of proof is trusted, and where the boundaries are; the
itemized list with file anchors is Table~\ref{sta:tab-leanitems} of
Section~\ref{sec:open}.

\paragraph{What is proved.} Two classification-line theorems are fully
formalized from first principles: for every odd $n\ge3$,
$\pn(K_{4,n})=(n+1)/2$ (with numeric corollaries $3,7,10$ at
$n=5,13,19$ following without independent gaps), and for every $k\ge1$,
$\pn(K_{2k+1})=k+1$ (via the edge-count lower bound and an explicit
Walecki-type construction; the endpoint $k=0$, where $\pn(K_1)=0$, is a
separate fact outside the theorem). Both carry the standard axiom profile
\texttt{[propext, Classical.choice, Quot.sound]} and no \texttt{sorryAx}.
On the criterion line, two declarations are formalized without
\texttt{sorryAx}: the criterion-level bound $g_{13}\le d_2+1$
(\texttt{thmA2663\_upgraded}, axioms \texttt{[propext, Quot.sound]}), where
$g_{13}$ is the criterion-layer count functional of the valuation line
(in the paper's notation, the number of ``leaks'' of $C_1$ --- points of
the shifted layer $F_{13}$ of Definition~\ref{int:def-legalval} outside
$C_1$; in the Lean sources it is defined as a count over
\texttt{List.range n} in the base module and quoted by
\texttt{ThmA2663.lean}), whose
hypotheses --- the target side has the prescribed size, and
$\mathrm{val}\le1$ --- are \emph{assumed, not proved} (it is a bound under a
hypothesis, not a feasibility-existence proof; see no-claim
\textup{(6)} of Section~\ref{int:subsec-claims}); and the case-$2$
structural discriminator (\texttt{case2\_struct\_alpha\_iff}, axioms
\texttt{[propext, Quot.sound]}), an if-and-only-if for the machine
predicate $C_1((t_1+g)\bmod N)=\texttt{true}$ under six explicit
hypotheses. Finally, the four-domain witness accounts of the termination
bound are kernel-checked at the level of the instance counts
$588/30{,}760/634{,}960/7{,}798{,}903$ and of two fully certified
five-column global ledgers:
$\mathrm{ledger}_6=(5{,}724;\,2{,}613;\,190;\,398;\,588)$
(legal pairs; witnesses; $\min\mathrm{val}=0$ bucket; $\min\mathrm{val}=1$
bucket; instances) and
$\mathrm{ledger}_8=(701{,}360;\,285{,}372;\,20{,}308;\,10{,}452;\,30{,}760)$.
The body of the termination theorem itself is proved on paper, not in Lean;
the accounts are its machine shadow.

\paragraph{The \texttt{native\_decide} boundary.} The finite certificates
(including both ledgers and the $N=10$/$N=12$ account modules) are closed by
\texttt{native\_decide}. Such theorems must be quoted as \emph{machine-checked
by kernel reduction}: they depend on the instantiation axiom of
\texttt{Lean.ofReduceBool} (the kernel-reflection axiom of the evaluation
mechanism), so they are not unproved assertions, but they do not carry the
trust level of pure term-level proofs. Three facts anchor the boundary. The
judging-layer semantics (the value functional, the board functionals, the
admission predicates) are inherited from the base module with zero
redefinition. The $N=8$ ledger is certified on two mutually pinned routes ---
whole-domain direct evaluation, and a $56$-bucket aggregation with per-bucket
certificates --- and cross-pinned against a second, independent Python
implementation, bucket by bucket. And the trust base is the honesty of the
Lean~4.34.0 kernel reflector itself.

\paragraph{The axiom audit.} The published declarations carry an axiom audit
(\texttt{\#print axioms} over the audit's probe file against the built
library). Re-running the audit against the current sources reproduces the
archived audit log entry for entry, and the six audited source files are
byte-identical to the review snapshot, so the log applies to the released
sources. The audit's complement is equally explicit: the probe covers the
declarations named in Table~\ref{sta:tab-leanitems} and the five Fan-line
declarations discussed next; anything outside the probe (the uniqueness
corollary of the Fan-type chain, the $N=10$/$N=12$ account modules, and the
auxiliary complete-graph and $K_{4,n}$ modules beyond the two named
theorems) is \emph{not} claimed fully formalized until individually probed.
A live comment-stripped scan finds exactly five \texttt{sorry} lines:
three placeholders of the Fan-type formalization chain
(\texttt{Astar.lean}, lines 2592, 2849, 2864), one \emph{refuted} research
lemma kept as a record only (\texttt{Fan.lean}, line 765 --- its statement
was disproved by a machine counterexample and is never cited; see
Remark~\ref{rem:claformal}), and one unproved research gap consumed by a
weakened variant (\texttt{Fan.lean}, line 4487). Accordingly the
Fan-type theorem chain
(Theorems~\ref{thm:claAstar} and~\ref{thm:lasep}) is cited as paper proofs
with formalization in progress, never as machine-verified
(Remark~\ref{rem:claformal}).

\paragraph{Certificate-grade example: the packing table.} The packing
residual of Section~\ref{subsubsec:criteria} carries its own Lean
certificates, which illustrate the certificate tier end to end. The main
certificate theorem checks, pointwise, all $79{,}200$ points of the table
domain ($120$ canonical orbits $\times$ $330$ demand vectors $\times$ the
two board sides): the packing predicate holds exactly when $\Sigma
d\le|S|-1$, or $\Sigma d=|S|$ with the demand outside the recorded
top-failure set. It is closed by \texttt{native\_decide} (zero
\texttt{sorryAx}, with the \texttt{ofReduceBool} disclosure attached), and
its data anchors are certified alongside it ($330$ demands; $120$ orbits;
$1{,}761$ top failures $=204+1{,}557$). Nine companion theorems over the
same table domain record: the downset property on both sides; the uniform
layer ($\Sigma d\le|S|-1$ implies feasibility on both sides); an explicit
no-threshold witness (two demands with equal $\Sigma d=4$, one feasible and
one not --- refuting any proposed threshold law of the form ``feasible
$\iff f(\Sigma d)$''); the concentration of all $1{,}761$ failures in the
top layer; the top-failure set being an upset; and the failure-orbit counts
$110$ and $2$. All are \texttt{native\_decide} certificates. Above them
sits the paper-level \emph{general} downset theorem
(\texttt{pkg\_downset}, axioms \texttt{[propext, Classical.choice,
Quot.sound]}, no \texttt{native\_decide}, no \texttt{sorryAx}), quantified
over arbitrary inputs under three slot-value hypotheses: if the slot
relation $\mathrm{slot}(d',i)+1=\mathrm{slot}(d,i)$ holds with all other
slots equal and $\mathrm{slot}(d,i)\ge1$, then feasibility of $d$ implies
feasibility of $d'$. Two boundaries travel with it: the hypotheses are
stated as slot-value relations (the subtraction form is not yet
formalized; instances are discharged through the slot relations), and the
Lean reimplementation of the packing predicate is not proven
program-level identical to the reference Python implementation --- that
identity rests on three-way machine cross-checks (the $79{,}200$-point
sweep with zero violations, bit-level reconstruction of the certificate
data, and a full Lean-side sweep).

\subsection{The key-table ledger for the covering identity}\label{com:subsec-ledger}

The machine account behind Proposition~\ref{prop:clacover} is organized as
\emph{key tables}: a key is the datum $(N,W,k,\delta_{12},\delta_{13})$ up
to rotation, and each key summarizes all instances sharing it. Across nine
domains $N\in\{6,\dots,14\}$ the key tables comprise $2{,}412{,}436$ keys
covering $98{,}602{,}396$ key-weighted admissible positions (the sum over
keys of $\binom{|W|}{d_1}$); the per-domain key counts are
$271/1{,}218/4{,}614/15{,}627/48{,}913/144{,}408/407{,}584/1{,}110{,}057/679{,}744$
for $N=6,\dots,14$. On every key: family$\setminus$prediction $=$
prediction$\setminus$family $=0$ (both inclusions of the covering identity
in key form), and the necessary-side boundary violations are $0$.

Four conventions bind every citation of these numbers; they are stated
here once and apply paper-wide.

\begin{enumerate}[label=\textup{(\arabic*)},leftmargin=2em]
\item \emph{Key-weighted versus instance-weighted.} The key tables count
each key once; instance-weighted accounts (each instance once) exist
separately for $N=8$ ($28{,}056$ instances, $221{,}480$ positions) and $N=9$
($128{,}016$ instances, $1{,}328{,}040$ positions) and are listed apart.
The two accountings are never mixed or compared.
\item \emph{Two domain conventions.} The domains $N\in\{6,\dots,13\}$ come
from exhaustive enumeration of the instance space, while $N=14$ comes from
an affine-normalized domain (orbit representatives under the normalized
action) whose equivalence to exhaustive enumeration is \emph{unproven}; the
nine-domain totals are therefore not a union under a single convention.
\item \emph{The no-$V_7$ superset reading.} All machine checks of the
covering identity are performed with the $V_7$ branch omitted
(Remark~\ref{rem:clacoverriders}), a conservative superset reading whose
zero-counterexample verdict transfers to the true face one-way. The residual
$V_7$-only mass is nonzero and precisely measured: $340$ positions at $N=8$
($1{,}266$ superset-uncovered minus $926$ true-face uncovered) and
$2{,}267{,}655$ instance-weighted at $N=14$.
\item \emph{No extrapolation.} Nothing in the ledger is extrapolated to
$N\ge15$, and the normalized domain is never read as the exhaustive domain.
\end{enumerate}

The ledger itself passed two independent controls: a re-aggregation with a
separate decoder and an independent combinatorial recomputation of the
position counts (matching the archived accounts digit for digit), and a
three-shard fresh replay with digit-identical accounting fields.

\paragraph{The independent audit.} Subsequent to, and separately from, the
verification passes above, an independent audit re-derived all twelve
supporting lemmas of the covering identity (the nine assembly lemmas and
the three micro-lemmas of Remark~\ref{rem:clacoverriders}) from the
definitional layer rather than by checking the write-up, probed each
lemma's plausible failure modes (none broke), probed each of the three
declared assembly gaps (none blocks the assembly), and re-ran the
domain-consistency check with a \emph{third} independent implementation
(exhaustive reference-form predicates, tag-level reverse direction) on
fresh key tables: four full domains ($N=6,7,8,10$; $55{,}016$ keys) plus
one $N=12$ shard ($50{,}948$ keys) --- $105{,}964$ keys in total, six
violation classes all zero, with independent aggregation matching the
headers and archived accounts digit for digit. The randomized machine
checks of the micro-lemmas were re-run in the audit pass with zero
violations ($4{,}000/1{,}109/4{,}000$ instances, after the first pass
$12{,}558/2{,}500/12{,}000$). The audit's verdict upgrades the package to
``paper candidate that has passed an independent audit'' --- which is
still not a mechanized or formalized theorem. The audit did \emph{not}
re-run, and these items are carried as archived accounts: the
single-swap falsification experiments of Section~\ref{subsubsec:criteria},
the transcription-fidelity anchor ($240$ instances, $97{,}848$ pairs), the
per-instance full-domain enumerations at $N=8,9$, and the per-instance
glue verification (the tool sources of the latter were lost with no
recoverable copy).

\appendix

\section{Computational anchors}\label{com:app-comp}

This appendix lists, for each theorem of the paper, the verification
status and the machine anchors. All computations used one pipeline
(Section~\ref{com:subsec-pipeline}); randomized steps used fixed seeds;
every number below is reproducible from the archived scripts released
with the paper. In the table, ``exact'' means computed by an exact
algorithm (dynamic program or exhaustive cover), ``witness'' means a
machine-verified decomposition, and ``census'' means all graphs in the
stated class were processed.

\begingroup\small
\begin{longtable}{p{4.6cm} p{6.2cm} p{4.2cm}}
\hline
Result & Verification & Anchors (script / certificate) \\
\hline
\endfirsthead
\hline
Result & Verification & Anchors (script / certificate) \\
\hline
\endhead
Theorem~\ref{thm:claAstar}, Theorem~\ref{thm:lasep}: witness instances &
exact $\pn\le\floor{n/2}$ on $101$ random (A1)--(A2) instances, $50$
instances with $\Delta(\EV)=5$, $15$ large witnesses; $0$ violations &
block-reload and witness-upgrade scripts\\
\hline
Theorem~\ref{thm:lasep}, Remark~\ref{rem:claCO},
Proposition~\ref{prop:claresidue}, Proposition~\ref{prop:claodddensity} &
$413$ p-anchored (A1)--(A2) instances ($13$ structured, $400$ random),
$0$ violations; exhaustive $n=5$ (all $1024$ graphs) and $n=6$ (all
$32768$ graphs) degree/triangle-layer checks behind
Proposition~\ref{prop:claodddensity} & structured-fix scripts\\
\hline
Theorem~\ref{thm:clatwohigh}, Corollary~\ref{cor:clatwohigh} (two-high
layer) & phantom-step audit $4/4$; triangle-forced exits of dirty shared
neighbours $34/34$; $34$ structured chain certificates; random two-high
census $144/144$ chains verified; exact DP cross-checks $43/43$; $0$
violations & two-high-layer scripts\\
\hline
Corollary~\ref{cor:clatwohigh}: the dirty corner & $207$ dirty-class
instances ($154$ legacy: $34$ structured $+\,120$ random; $53$ new), plus
$16$ LB1-tightness hardening runs and $60$ two-high regression chains; in
every instance a decomposition of $\tilde G$ into at most
$\floor{n/2}-1$ paths was produced (constructive refutation of the
corner); $0$ corner hits, $0$ violations & corner-narrowing scripts\\
\hline
Corollary~\ref{cor:clatwohigh}, Remark~\ref{rem:clagap}: band-free
tightness & $125$ decision-level instances of the constructed family:
all $66$ band-free instances sealed constructively (paired-decomposition
constructor output verified to partition the edge set into exactly
$\mathrm{odd}/2$ simple paths, with per-vertex endpoint assertions); all
$59$ banded instances rejected with band certificates; constructor
success $\iff$ band-free reconciled $125/125$; excess
$\mathrm{excess}\le\lfloor y/2\rfloor$ audited on these $125$ and the
$110$ earlier band-level instances of the same pipeline ($235$ in all),
zero violations & band-tightness scripts\\
\hline
Remark~\ref{rem:clagap}: the band-free statement beyond the constructed
family & $71$ decision-level instance runs: the general band-free
statement refuted by cycle-block instances certified three ways
(band-free by the criterion, untight by the decision procedure, and an
exact combinatorial witness at $\pn=\mathrm{odd}/2+1$ for the
cycle-with-one-pendant block and $\pn=\mathrm{odd}/2+2$ for the
double-triangle block); $0$ failures & relaxed-domain scripts\\
\hline
Proposition~\ref{prop:claWmbound}, Proposition~\ref{prop:claescape}:
$W(m)$ & parity certificate by hand; decompositions into exactly
$\floor{n/2}$ paths verified for $m\in\{5,7,9\}$; class memberships
asserted programmatically & witness-family scripts\\
\hline
Proposition~\ref{prop:claoscsharp}: exactness on tested $k$ & DP exact at
$k=2$; verified $(k+1)$-path witnesses for $k\in\{2,3,4,5,6,8,10\}$ &
witness-upgrade scripts\\
\hline
Theorem~\ref{thm:conk4odd} & structured exact search over $K_{4,n}$,
$n=3..31$, recovering $(n+1)/2$ independently for every $n$; closed-form
construction verified for $k=6..301$ ($n=13..603$); witness lists
machine-verified & $K_{4,n}$ scripts and witness files\\
\hline
Proposition~\ref{prop:conk4small} & witness verification (edge-exact
partition, simplicity) & small-case witness files\\
\hline
Even columns of Case~4 (Section~\ref{subsec:columns};
Theorem~\ref{thm:concapgen}, Theorem~\ref{thm:constructgen},
Theorem~\ref{thm:conringclosure}, Corollary~\ref{cor:concol6},
Observation~\ref{obs:concostwall}) & $n_2=6$: witnesses for
$K_{7,6},K_{9,6},K_{11,6},K_{13,6}$ each verified by two independent
implementations (pipeline verifier $+$ adjacency-matrix checker); ring
construction re-derived for $k=7..31$ from the base $k_0=3$, with
overlapping chains from $k_0=4,5$. $n_2=8$: bases
$K_{9,8},K_{11,8},K_{13,8}$ found at $5/7/43$ search nodes and verified
($K_{11,8},K_{13,8}$ re-derived byte-identically); closure sampled at
$60$ levels $k=9..30$; negatives on $K_{15,8},K_{17,8}$: $>14$M nodes
over all five deficit multisets, zero exhaustions. $n_2=10$: bases
$K_{11,10},K_{13,10},K_{15,10}$ found at $6/13/182$ nodes and
triple-verified; closure sampled $k=11..18$; negatives on
$K_{17,10},K_{19,10},K_{21,10}$, $\approx0.9$M nodes, zero exhaustions &
column scripts and witness files\\
\hline
Theorem~\ref{thm:conL0} (anchor) & Walecki partition checked edge by edge
for $j=2..16$; formula verified for all odd $j\le121$ & engine scripts\\
\hline
Theorem~\ref{thm:conchain} assembly & $14$ random chains
($b\in\{2..5\}$) assembled and verified; $b=4,5$ chains attain
$\sum j_i-1$ & engine scripts\\
\hline
Remark~\ref{rem:conchaincomp}(2): three $K_5$'s & the $5$-path
decomposition printed in the text is hand-checkable; contradicts the
archived exact value $6$ (re-verification requested) & this paper, project
archive\\
\hline
Conjecture~$K$ (Definition~\ref{int:def-budget}) evidence & the five
exhaustively enumerated domains of Table~\ref{com:tab-kdomains}; $0$
violations & domain scripts and certificate files\\
\hline
Proposition~\ref{prop:claincomp}: exact values & DP exact:
$\pn(K_{2,5}+2\ell)=3$, $\pn(G')=4$ (triangle plus two decorated
bridges), $n=9,11$ & witness-upgrade scripts\\
\hline
Lemma~\ref{int:lem-fan6}, both orientations (transcription and
verification note in Section~\ref{int:subsec-fanlemmas}) &
direction-corrected scan over both ordered (gains-end, budget-end)
orientations of every deleted edge: $n\le5$ complete, all $482{,}080$
ordered instances ($241{,}040$ per orientation), $0$ violations; $n=6$
complete, all $1{,}701{,}141{,}150$ ordered instances ($850{,}570{,}575$
per orientation; the fixed-orientation half exactly reproduces the
historical census) with $0$ violations & direction-scan scripts\\
\hline
\end{longtable}
\endgroup

\section{The corner analysis and the band-free tightness theorem}
\label{com:app-band}

This appendix records, with complete proofs, the theorem-level material
behind the band-free tightness statements used in
Remark~\ref{rem:clagap}, in the proof of Corollary~\ref{cor:clatwohigh},
and in open problem O2. Throughout the appendix, $G$ is a graph of the
\emph{dirty class} of Corollary~\ref{cor:clatwohigh}: $G$ is connected
and satisfies {\rm(A1)}--{\rm(A2)}, has exactly two high vertices $u$ and
$x$ with $k=\dEV(u)$ and $l=\dEV(x)$ (both odd and $\ge5$ by (A2)),
$t=t_{ux}\ge1$ common EV-neighbours $v_1,\dots,v_t$, every common
EV-neighbour of E-degree $3$, and no private EV-neighbour of E-degree $1$
on either side. We write $z_j$ for the third EV-neighbour of $v_j$, $F$
for the E-star at $u$, and $\tilde G_j=G-F-\{xv_j\}$, abbreviating
$\tilde G=\tilde G_1$. By (A1) the common EV-neighbours are pairwise
non-adjacent inside $\EV(G)$ and each $z_j$ is a \emph{far} even vertex
of $\tilde G_j$ (an even vertex outside $N_{\EV}(u)\cup N_{\EV}(x)$ that
is not a common EV-neighbour); also $\dEV(z_j)\le3$, since
$\dEV(z_j)\ge4$ would make $z_j$ a third high vertex.

\begin{definition}[corridors, terminal units, bands]\label{com:def-band}
Fix $j$ and write $\tilde G$ for $\tilde G_j$. For an even vertex $f$ of
$\tilde G$, let $Y(f)$ be its set of \emph{corridor} neighbours: the
degree-two even vertices $y$ with $N_{\tilde G}(y)=\{x,f\}$, where $N(f)$
denotes the neighbour set of $f$ in $\tilde G$. The remaining incident
edges of $f$ are its \emph{terminal units}, numbering $|N(f)|-|Y(f)|$.
Call $f$ a \emph{deadlock band} if $|Y(f)|>|N(f)|-|Y(f)|$. Let
$\mathrm{Far}$ be the set of far even vertices of $\tilde G$ and $Y$ the
total number of corridors (each corridor is a private EV-neighbour of
$x$, so $Y\le l-t$, with equality in the constructed family below).
\end{definition}

\begin{definition}[constructed family]\label{com:def-family}
$G$ belongs to the \emph{constructed family} if it lies in the dirty
class and: (i) every private vertex of E-degree $2$ (even in $G$, on
either the $u$- or the $x$-side) has its second EV-neighbour in
$\mathrm{Far}$, and no EV-edge of $\tilde G_1$ joins two far even
vertices; (ii) the odd-neighbour block of $u$ --- the substructure hanging
at $u$'s odd-degree neighbours --- is a set of $s$ pendant leaves at $u$
with $s$ odd, or a number of gadgets ($u$--$o$, $o$--$a$, $o$--$b$, with
two pendant leaves on each of $a$ and $b$), possibly mixed; and likewise
for $x$; (iii) every far vertex $f$ of $\tilde G_1$ has $\dEV(f)\le3$ and
carries at most one repair leaf, whence $\deg_{\tilde G_1}(f)\in\{2,4\}$;
and (iv) every common EV-neighbour of $G$ is adjacent to an odd positive
number of pendant leaves and to no odd-degree vertices other than these
leaves. (The band condition of Definition~\ref{com:def-band} does not
depend on $j$: no far--far EV-edge exists by (i), so deleting $xv_j$
touches no far vertex and none of its neighbours. The machine instances
of Appendix~\ref{com:app-comp} all lie in the family, with the
cross-private instances covered by the extension that follows.)
\emph{Extension.} Cross-private pairs are admitted: a $u$-private vertex
$w$ may be adjacent to an $x$-private vertex $y$, so that
$N_{\EV}(w)=\{u,y\}$ and $N_{\EV}(y)=\{x,w\}$; $y$ then acts as a
terminal corridor with $\phi(y)=w$, and $wy$ is a terminal unit of $w$.
The band-free tightness theorem below covers this extension by the
identical step (ii) of its proof; the load account and the counting
theorem are used only on the family proper, where every private even
vertex has its second EV-neighbour in $\mathrm{Far}$.
\end{definition}

\subsection{The counting gap and the counting theorem}

\begin{lemma}[counting-gap ledger]\label{com:lem-gapledger}
For any $j$, let $\mathrm{odd}:=\mathrm{odd}(\tilde G_j)$ and
$\mathrm{gap}:=\floor{n/2}-\mathrm{odd}/2$, where $n=|V(G)|$ and
$\mathrm{odd}_G$ is the number of odd-degree vertices of $G$. Then
$\mathrm{odd}=k+1+\mathrm{odd}_G$,
$n=2+t+(k-t)+(l-t)+|\mathrm{Far}|+\mathrm{odd}_G$, and
\[
\mathrm{gap}=\frac{l-t+|\mathrm{Far}|+1}{2}\ \ (n\text{ even}),
\qquad
\mathrm{gap}=\frac{l-t+|\mathrm{Far}|}{2}\ \ (n\text{ odd});
\]
in particular $\mathrm{gap}\ge1$.
\end{lemma}

\begin{proof}
The odd-degree vertices of $\tilde G_j$ are: $u$, which loses the $k$
star edges (an odd number); $x$, which loses $xv_j$; the common
EV-neighbours $v_i$ with $i\ne j$ ($t-1$ of them, each losing $uv_i$);
the private EV-neighbours of $u$ ($k-t$ of them, each losing $uw$); and
all odd-degree vertices of $G$ (untouched, since the star edges all join
even-degree vertices and $xv_j$ joins two even-degree vertices). Hence
$\mathrm{odd}=2+(t-1)+(k-t)+\mathrm{odd}_G=k+1+\mathrm{odd}_G$. The
vertex set splits into $\{u,x\}$, the common EV-neighbours, the private
EV-neighbours of $u$ and of $x$, the far even vertices (which include the
$z_i$), and the odd-degree vertices of $G$, giving the displayed count.
For odd $n$: the remaining terms $2$, $k+l$ and $\mathrm{odd}_G$ are even,
so $n$ is odd exactly when $t+|\mathrm{Far}|$ is odd, which makes
$l-t+|\mathrm{Far}|$ even and the second quotient an integer. Finally
$l\ge t$ and $|\mathrm{Far}|\ge1$ (as $z_j\in\mathrm{Far}$): for even $n$
the numerator is $\ge2$; for odd $n$ it is a positive even number, hence
$\ge2$. So $\mathrm{gap}\ge1$.
\end{proof}

\begin{lemma}[load account]\label{com:lem-load}
In the family proper (the constructed family without cross-private
pairs), $\sum_{f\in\mathrm{Far}}\dEV(f)=k+l-t$ (E-degrees read in $G$).
Consequently $\dEV(f)\le3$ for every $f\in\mathrm{Far}$, and
$|\mathrm{Far}|\ge2$.
\end{lemma}

\begin{proof}
Each common EV-neighbour $v_i$ has EV-neighbours $u$, $x$ and
$z_i\in\mathrm{Far}$, contributing $t$ EV-edges with one end in
$\mathrm{Far}$; each private even vertex (of E-degree $2$ in $G$, on
either side) has its second EV-neighbour in $\mathrm{Far}$ by
Definition~\ref{com:def-family}(i), contributing the $k-t$ edges at the
private EV-neighbours of $u$ and the $l-t$ edges at the corridors of $x$.
By (i) there are no EV-edges between far vertices, and odd-degree
vertices carry no EV-edges; so every EV-edge counted has exactly one end
in $\mathrm{Far}$, and summing the E-degrees over $\mathrm{Far}$ gives
$t+(k-t)+(l-t)=k+l-t$. If some $f$ had $\dEV(f)\ge4$, then $f$ would be a
third high vertex, impossible; so all E-degrees in $\mathrm{Far}$ are
$\le3$. If $|\mathrm{Far}|=1$, then $k+l-t\le3$, so
$t\ge k+l-3\ge5+5-3=7>\min(k,l)\ge t$, a contradiction; hence
$|\mathrm{Far}|\ge2$.
\end{proof}

\begin{proposition}[the counting theorem: small gap forces band-free]
\label{com:prop-gapnoband}
In the family proper (the constructed family without cross-private
pairs), if $\mathrm{gap}\le2$, then $\tilde G_j$ has no deadlock band,
for every $j$.
\end{proposition}

\begin{proof}
By Definition~\ref{com:def-family} the band condition is
$j$-independent, so we argue in $\tilde G_1$. Suppose $f$ is a band. Then
at least two corridors hang at $f$, and every corridor is a private
EV-neighbour of $x$; write $y:=l-t\ge Y\ge2$.
If $n$ is even, Lemma~\ref{com:lem-gapledger} gives
$y+|\mathrm{Far}|\le 2\,\mathrm{gap}-1\le3$, while $y\ge2$ and
$|\mathrm{Far}|\ge2$ (Lemma~\ref{com:lem-load}) give
$y+|\mathrm{Far}|\ge4$ --- a contradiction. If $n$ is odd,
$\mathrm{gap}\le2$ gives $y+|\mathrm{Far}|\le4$, so
$y\le4-|\mathrm{Far}|\le2$; equality throughout forces $|\mathrm{Far}|=2$
and $y=2$. The load account gives $k+l-t=k+y=k+2\le3|\mathrm{Far}|=6$, so
$k\le4$, contradicting $k$ odd and $\ge5$. So $y=2$ is impossible, and
the bound $y\le2$ then leaves $y\le1$, i.e.\ $l-t\le1$: no band exists (a
band needs two corridors at one far vertex).
\end{proof}

\subsection{The pairing framework and the band-free tightness theorem}

\begin{proposition}[pairing schemes produce trail decompositions]
\label{com:prop-pairframe}
Let $H$ be a finite simple graph with a local structure on the edges at
each vertex: if $\deg_H(v)$ is even, its incident edges are partitioned
into pairs; if $\deg_H(v)$ is odd, exactly one incident edge is marked
\emph{free} and the rest are paired. Then $E(H)$ decomposes into maximal
alternating trails --- starting from a free mark, at each internal visit
leaving through the mate of the arriving edge --- and every edge lies on
exactly one trail. Each trail either has two free ends or is closed; the
open trails number $\mathrm{odd}(H)/2$. If moreover no trail repeats a
vertex, then $H$ has a path decomposition $D$ with
$|D|=\mathrm{odd}(H)/2$, every odd-degree vertex an endpoint of exactly
one path of $D$ and no even-degree vertex an endpoint; in particular
$\pn(H)=\mathrm{odd}(H)/2$.
\end{proposition}

\begin{proof}
At each vertex the successor of an arriving edge is uniquely determined
by the pairing, so the trails are edge-disjoint and exhaust $E(H)$. A
free mark can occur only at a trail end (a free edge carries at most one
free mark at each of its ends); an odd vertex has exactly one free mark
and is therefore the end of exactly one trail; an even vertex has none
and is internal wherever it occurs. The number of open trails is the
number of free marks divided by two, i.e.\ $\mathrm{odd}(H)/2$. If all
trails are paths, they form a decomposition $D$ with $|D|=\mathrm{odd}/2$
and endpoint set exactly the odd vertices; the odd-degree lower bound of
Lemma~\ref{int:lem-LB1}(i) gives $\pn(H)\ge\mathrm{odd}/2$, so the value
is exact.
\end{proof}

\begin{lemma}[odd-neighbour blocks close by themselves]\label{com:lem-blocks}
In the constructed family, the odd-neighbour block of $c\in\{u,x\}$ can
be structured as Proposition~\ref{com:prop-pairframe} requires, with all
resulting trails staying inside the block, as follows. A leaf block of
$s$ leaves ($s$ odd): mark $c\ell_1$ free at $c$ (the single-edge path
$\ell_1$--$c$) and pair $(c\ell_{2i-1},c\ell_{2i})$ for $2\le 2i\le s$,
giving $(s+1)/2$ paths for the $s+1$ odd vertices involved. A gadget:
pair $c$--$o$ free at $c$ (path $c$--$o$), $(o$--$a,o$--$b)$ at $o$,
$(a\ell_1,a o)$ at $a$ with $a\ell_2$ free, and $(b o,b\ell_3)$ at $b$
with $b\ell_4$ free; the resulting paths $c$--$o$,
$\ell_1$--$a$--$o$--$b$--$\ell_3$, $\ell_2$--$a$, $\ell_4$--$b$ have
endpoint multiset covering $\{c,o,a,b,\ell_1,\dots,\ell_4\}$ exactly once
--- $8$ odd vertices, $4$ paths. A mixed block (gadgets and leaves at the
same $c$): the block edges at $c$ number odd in total, so exactly one is
free (chosen canonically, say the first), the rest pair arbitrarily among
leaf and gadget edges, and a mixed pair $(c\ell,c o)$ forms the legal
path $\ell$--$c$--$o$. All trails end at free marks inside the block, so
the block never interacts with the rest of the scheme.
\end{lemma}

\begin{proof}
This is a verification of the pairing tables just displayed: within a
gadget the pairs at $o,a,b$ are fixed as shown; leaves pair among the
edges at $c$; and in a mixed block the total number of block edges at $c$
is $s+g$ ($s$ leaves, $g$ gadgets, each gadget contributing a single
attachment edge $co$), which is odd because $c$ has even degree in $G$
and an odd number of EV-neighbours, so the free/paired structure exists;
the endpoint count of each displayed path list matches the odd vertices
it involves, each exactly once. No trail leaves the block, since free
marks sit at block vertices only.
\end{proof}

\begin{lemma}[forced fusions]\label{com:lem-fuse}
In any decomposition of $\tilde G$ as in the last sentence of
Proposition~\ref{com:prop-pairframe} (a \emph{paired decomposition}), the
edges at every even vertex are traversed in internal pairs. Hence a
corridor $y$ (degree $2$) forces its two edges $xy$ and $y\phi(y)$ into
one path; and the unit at $v_1$ is forced: $v_1z_1$ pairs with
$v_1\ell_1$, where $\ell_1$ is the leaf at $v_1$ (any further leaves at
$v_1$ pair among themselves, their number being even by
Definition~\ref{com:def-family}(iv)).
\end{lemma}

\begin{proof}
At an even vertex $v$ of a paired decomposition, $D(v)=0$, so no path has
$v$ as an endpoint; each visit is internal and uses exactly two edges, so
the incident edges split into internal pairs by path. A degree-two vertex
therefore has its whole edge set as a single pair on a single path.
\end{proof}

\begin{proposition}[injection at far vertices]\label{com:prop-inject}
Fix a far vertex $f$ of $\tilde G$ with corridor set $Y(f)$, $c(f):=|Y(f)|$,
and terminal units $T(f)$, $|T(f)|=\deg(f)-c(f)$. The edges at $f$ can be
paired so that every corridor edge $fy$ is paired with a \emph{distinct}
terminal unit if and only if $f$ is not a band, i.e.\ $c(f)\le|T(f)|$ (the
pairing at $f$ is otherwise unconstrained, so this is a matching condition
in a complete bipartite graph between $Y(f)$ and $T(f)$). In the
constructed family $c(f)\le2$.
\end{proposition}

\begin{proof}
Only the bound remains: by Definition~\ref{com:def-family}(iii)
$\dEV(f)\le3$ and $f$ carries at most one repair leaf, so
$\deg(f)\in\{2,4\}$, and $c(f)\le\deg(f)/2\le2$.
\end{proof}

\begin{lemma}[sorted offset matching at $x$]\label{com:lem-offset}
In the constructed family, $x$ is odd in $\tilde G$, with exactly one
free edge; its non-block edges are the $m$ corridor edges (each labelled
by the far vertex it hangs at, with multiplicity $c(f)\le2$ per label; a
cross-private corridor counts as a label of its own, since no two of them
hang at the same $w$) and the $s=t-1$ bridge edges $xv_i$ ($i\ge2$), with
$m+s=l-1$ even. There is a pairing of these edges in which no two paired
corridor edges carry the same label:
if $m$ is even, sort the corridors $p_1,\dots,p_m$ by (label, index) and
pair $(p_i,p_{i+m/2})$ for $i=1,\dots,m/2$, pairing the bridges among
themselves --- if $p_i,p_{i+m/2}$ shared the label $f$, the label $f$
would occupy at least $m/2+1$ of the sorted positions, so
$c(f)\ge m/2+1\ge3$ for $m\ge4$, contradicting
Proposition~\ref{com:prop-inject};
if $m=2$ with equal labels, pair each of the two corridors with a bridge
and the remaining bridges among themselves ($s-2$ is even since
$s\equiv m\pmod2$); the failing configuration $m=2$, equal labels, $s=0$
is unreachable, since $s=0$ forces $t=1$ and $m=l-1$, i.e.\ $l=3$,
contradicting $l\ge5$;
if $m$ is odd (possible only when $t$ is even), extract one corridor of a
label of maximum multiplicity together with one bridge into a mixed pair,
and apply the even case to the rest.
\end{lemma}

\begin{proof}
The counting arguments displayed are the whole proof: for the even case,
if both $p_i$ and $p_{i+m/2}$ carried label $f$, then all positions from
$i$ to $i+m/2$ would carry label $f$ (the sort is by label first), giving
$m/2+1$ corridors at $f$. For $m=2$ with equal labels and $s\ge2$, the
corridor--bridge pairs leave $s-2\ge0$ bridges, even in number. For odd
$m$, extracting one corridor of a maximum-multiplicity label and one
bridge leaves an even number of each, and the extracted pair mixes a
corridor with a bridge, so the no-equal-label condition is preserved.
\end{proof}

\begin{theorem}[band-free tightness in the constructed family]
\label{com:thm-bandfree}
Let $G$ belong to the constructed family and fix $j$. If $\tilde G_j$ has
no deadlock band, then $\tilde G_j$ admits a paired decomposition;
consequently $\pn(\tilde G_j)=\mathrm{odd}(\tilde G_j)/2$.
\end{theorem}

\begin{proof}
We define the local structure of Proposition~\ref{com:prop-pairframe} on
$\tilde G_j$.
(i) The odd-neighbour blocks of $u$ and of $x$ are structured as in
Lemma~\ref{com:lem-blocks}, $x$'s free edge being the canonical first
edge of its block.
(ii) The degree-two even vertices are fused as in Lemma~\ref{com:lem-fuse}:
each corridor $y$ pairs $xy$ with $y\phi(y)$, where
$\phi(y)\in\mathrm{Far}$ is the far vertex at its other end (for the
cross-private \emph{extension} of Definition~\ref{com:def-family} the
same fusion applies with $\phi(y)=w$ a $u$-private vertex, $wy$ being a
terminal unit of $w$, so the extension is covered by the identical
argument); the unit at $v_j$ pairs $v_jz_j$ with $v_j\ell_j$. A shared
vertex $v_i$, $i\ne j$, takes the pair $(v_ix,v_i\ell_i)$, marks $v_iz_i$
free, and self-pairs its remaining leaf edges.
(iii) At each far vertex $f$, Proposition~\ref{com:prop-inject} ---
applicable because $\tilde G_j$ is band-free --- pairs every corridor
edge $fy$ with a distinct terminal unit; the remaining terminal units
number $\deg(f)-2c(f)$, which is even, and pair among themselves.
(iv) At $x$, the free edge is the block edge of (i), and the remaining
edges pair according to Lemma~\ref{com:lem-offset}.
(v) Every degree-one vertex carries the free mark at its leaf end. The
free-mark audit is immediate: each odd vertex of $\tilde G_j$ carries
exactly one free mark and each even vertex none, by construction.

\emph{No trail closes.} Consider any trail of the resulting scheme. At a
far vertex a corridor end is always paired with a terminal unit, so a
trail contains at most two corridors and, having left a far vertex, never
returns to it; at $x$ each trail passes exactly once (a single pair); the
fusions at $v_j$ and the bridges at the $v_i$ are one-way terminations.
Hence every trail has one of the forms
\[
\text{terminal}--f--y--x--y'--f'--\text{terminal}',\qquad
\text{terminal}--f--y--x--v_i--\ell_i,
\]
\[
\ell_i--v_i--x--v_{i'}--\ell_{i'},\qquad
\text{terminal}--f--\text{terminal}',
\]
where in the first form $y$ and $y'$ are corridors of \emph{distinct} far
vertices $f\ne f'$, by Lemma~\ref{com:lem-offset}; a ``terminal'' means a
terminal unit followed to its free end --- the whole segment from $f$
along a repair-leaf edge, along a hook edge to a $u$-private vertex $w$,
or along the $z$-role edge $v_iz_i$ read from $z_i$ to its free mark at
$v_i$. When the corridor $y'$ of the first form is cross-private, its
right-hand segment $y'$--$f'$--terminal$'$ is replaced by $y'$--$w$: the
trail ends at the free mark of $w$, an odd vertex of $\tilde G$ (so no
closed trail arises). In each form the vertices are pairwise distinct, so
every trail is a simple path. Proposition~\ref{com:prop-pairframe} yields
a paired decomposition, and the value of $\pn$ follows.
\end{proof}

\begin{lemma}[a band blocks paired decompositions]\label{com:lem-bandblocks}
If $\tilde G_j$ has a deadlock band $f$, then $\tilde G_j$ admits no
paired decomposition; consequently
$\pn(\tilde G_j)\ge\mathrm{odd}(\tilde G_j)/2+1$.
\end{lemma}

\begin{proof}
Two corridors $y_1\ne y_2$ of the same band $f$ cannot lie on one path of
any decomposition: if a path contained both $fy_1$ and $fy_2$, then
$y_1$, being of degree two and internal (even vertices are never
endpoints in a paired decomposition), would force $xy_1$ onto the same
path, likewise $xy_2$, so the path would contain all four edges of the
cycle $x$--$y_1$--$f$--$y_2$--$x$ --- impossible for a simple path. Now in
a paired decomposition the $\deg(f)$ edges at $f$ split into
$\deg(f)/2$ internal pairs. The edge $fy$ of a corridor $y\in Y(f)$
cannot pair with another corridor edge $fy'$ (previous sentence), so it
pairs with an edge at a terminal unit; the resulting assignment
$y\mapsto$ its paired edge is injective, whence
$c(f)\le|N(f)|-c(f)$, contradicting the band condition. So no paired
decomposition exists, i.e.\ no decomposition with $\mathrm{odd}/2$ paths
and endpoint set the odd vertices; with fewer than $\mathrm{odd}/2$ paths
excluded by Lemma~\ref{int:lem-LB1}(i), this gives
$\pn\ge\mathrm{odd}/2+1$.
\end{proof}

\begin{proposition}[one band costs exactly one path]\label{com:prop-patch}
Let $G$ belong to the constructed family and suppose $\tilde G_j$ carries
exactly one deadlock band $f$. Then
$\pn(\tilde G_j)=\mathrm{odd}(\tilde G_j)/2+1$.
\end{proposition}

\begin{proof}
For a single band the pair $(c(f),|T(f)|)$ is $(2,0)$ or $(3,1)$: the
band condition $c(f)>|T(f)|$ with $c(f)+|T(f)|=\deg(f)\in\{2,4\}$ and
$c(f)\ge2$ leaves exactly these two cases, and in both the surplus of
corridor ends over terminal units has size $|S(f)|=2$; write
$S(f)=\{a,a'\}$. Modify the scheme of Theorem~\ref{com:thm-bandfree} in
three places only:
(1) at $f$, self-pair the surplus edges $(fa,fa')$ instead of injecting
them;
(2) at $x$, spend the free edge on $xa$ rather than on a block edge,
marking $xa$ free at both ends (the single-edge path $x$--$a$);
(3) the remaining injections at $f$ continue as before (in the case
$(3,1)$ the one remaining corridor pairs with the one terminal unit);
every other edge follows Theorem~\ref{com:thm-bandfree}, whose offset
matching at $x$ keeps distinct labels apart automatically, including among
the remaining corridors of the band.
The resulting trails are $a$--$f$--$a'$--$x$--\dots--terminal (with $a'$
fused, $x$ passed once, $f$ passed once) and the single edge $x$--$a$;
all are simple, and the $4$-cycle $x$--$a$--$f$--$a'$--$x$ is cut at both
$x$ and $a$ by the free marks of $xa$, so it does not arise as a trail.
The endpoint count: $a$ gains two ends (one on the $x$ side, one on the
$f$ side) and every other vertex is unchanged, so the number of path-ends
is $\mathrm{odd}+2$ and the number of paths is $\mathrm{odd}/2+1$. The
reverse inequality is Lemma~\ref{com:lem-bandblocks}.
\end{proof}

\begin{proposition}[several bands]\label{com:prop-excb}
If $\tilde G_j$ carries $B$ deadlock bands, then
$\pn(\tilde G_j)\le\mathrm{odd}(\tilde G_j)/2+2B-1$; moreover
$B\le\lfloor Y/2\rfloor$, where $Y$ is the total number of corridors.
\end{proposition}

\begin{proof}
Start from the scheme of Theorem~\ref{com:thm-bandfree} and repair the
bands one at a time. The first band is repaired as in
Proposition~\ref{com:prop-patch} (cost $+1$, spending $x$'s free edge and
one ordinary matching partner). Each further band is repaired by the
double-endpoint modification: self-pair the band's surplus pair at its
far vertex $f$ and self-pair two $x$-edges at $x$, freeing four ends; the
two resulting trails $a$--$x$--$a'$ and $a$--$f$--$a'$ add two ends, that
is, one path, per band --- an exchange costing two endpoint pairs for one
path, relying on no matching partner. Hence
$p\le\mathrm{odd}/2+1+2(B-1)=\mathrm{odd}/2+2B-1$. Finally, distinct
bands have disjoint corridor sets (each corridor hangs at exactly one far
vertex) and each band has $c\ge2$, so
$Y=\sum_{f\text{ band}}c(f)\ge2B$.
\end{proof}

\begin{corollary}[bands are exactly the obstructions]\label{com:cor-bandiff}
In the constructed family: $\tilde G_j$ has a deadlock band if and only
if $\pn(\tilde G_j)>\mathrm{odd}(\tilde G_j)/2$. In particular, on the
family proper (the constructed family without cross-private pairs), if
the counting gap is at most $2$, then
$\pn(\tilde G_j)=\mathrm{odd}(\tilde G_j)/2=\floor{n/2}-\mathrm{gap}$ for
every $j$; the cross-private instances at gap at most $2$ are covered at
probe level by the machine certificates of Appendix~\ref{com:app-comp}.
\end{corollary}

\begin{proof}
One direction is Lemma~\ref{com:lem-bandblocks}, the other
Theorem~\ref{com:thm-bandfree}. If $\mathrm{gap}\le2$,
Proposition~\ref{com:prop-gapnoband} gives band-free for every $j$ on the
family proper, and Theorem~\ref{com:thm-bandfree} with
Lemma~\ref{com:lem-gapledger} gives
$\pn=\mathrm{odd}/2=\floor{n/2}-\mathrm{gap}$.
\end{proof}

\subsection{The five escape routes at a common EV-neighbour}

\begin{theorem}[escape routes]\label{com:thm-escape}
Let $G$ be a dirty-class graph, fix a common EV-neighbour $v_1$ with
third EV-neighbour $z_1$, write $B=\floor{n/2}$ and $\tilde G=\tilde G_1$,
and let $Y$ be the set of private EV-neighbours of $x$. If any of the
following holds, then $\pn(G)\le B$:
\begin{enumerate}[label=\rm(\alph*),leftmargin=2em]
\item $\pn(\tilde G)\le B-1$;
\item $\tilde G$ has a decomposition $\tilde D$ with $|\tilde D|\le B$ and
$\tilde D(x)\ge3$;
\item $\tilde G$ has a decomposition $\tilde D$ with $|\tilde D|\le B$,
$\tilde D(x)=1$ and $\tilde D(z_1)\ge1$;
\item $\tilde G$ has a decomposition $\tilde D$ with $|\tilde D|\le B$,
$\tilde D(v_1)\ge2$ and $\tilde D(v_1)>\#\{y\in Y:\tilde D(y)=0\}$;
\item $\tilde G$ has at least two components of odd order.
\end{enumerate}
\end{theorem}

\begin{proof}
The reinsertion skeleton is that of the proof of
Theorem~\ref{thm:clatwohigh}. (s1) Reinsert $xv_1$ towards $v_1$ by
Lemma~\ref{int:lem-fan6}; the condition to check is
$\tilde D(x)>\#\{w\in N_{\tilde G}(v_1):\tilde D(w)=0\}$. Among the
neighbours of $v_1$ in $\tilde G$, the only even-degree vertex is $z_1$:
$x$ is odd in $\tilde G$ (it lost $xv_1$), and every other neighbour of
$v_1$ is an odd-degree vertex of $G$ (the edge $uv_1$ is deleted, and
$v_1$ has E-degree $3$ with EV-neighbours $u,x,z_1$). So the passing
count at $v_1$ is at most $1$. (batch) Lemma~\ref{int:lem-fan7} then
absorbs the whole star $F$ towards $u$, returning $A\subseteq F$ with
$|A|\ge\ceil{k/2}$; its hypothesis holds read over $G$ because $v_1$ is
now active and every other star vertex is odd in $\tilde G$, hence has
$\tilde D\ge1$ by Lemma~\ref{int:lem-LB1}(ii). (L) Each edge of
$F\setminus A$, $|F\setminus A|=\floor{k/2}$, is reinserted towards its
star vertex as in the proof of Theorem~\ref{thm:claAstar}: dead
candidates are confined to $N_{\EV}(v)\setminus\{u\}$, of size at most
$\dEV(v)-1\le2$ by (A2), against a budget of at least $3$ at the last
step.

(b) $\tilde D(x)\ge3>1\ge$ passing$(v_1)$, so (s1) applies and activates
$v_1$; (batch) and (L) follow.

(a) Let $D^*$ be optimal with $|D^*|\le B-1$. If $D^*(x)\ge3$, route (b)
applies. Otherwise $D^*(x)=1$ (it is odd by Lemma~\ref{int:lem-LB1}(iii),
$x$ being odd in $\tilde G$). Now
$\deg_{\tilde G}(x)=\deg_G(x)-1\ge l+1-1\ge5$ (as $x$ is high,
$\deg_G(x)\ge l+1$), so $x$ is internal to some path of $D^*$; cutting
that path at $x$ raises $|D^*|$ by one --- still at most $B$ --- and
raises the endpoint count at $x$ to $3$: route (b).

(c) $\tilde D(z_1)\ge1$ makes the passing count at $v_1$ equal to $0$, so
(s1) applies unconditionally; (batch) and (L) as in (b).

(d) Reverse the order of (s1): reinsert $xv_1$ \emph{towards $x$} first.
The Lemma~\ref{int:lem-fan6} condition is now applied at $v_1$ as the
donor: $\tilde D(v_1)>\#\{w\in N_{\tilde G}(x):\tilde D(w)=0\}$, and the
dead candidates at $x$ are precisely the dead private EV-neighbours of
$x$ (the other neighbours of $x$, namely $v_2,\dots,v_t$ and the
odd-degree neighbours, all have $\tilde D\ge1$); hypothesis (d) is
exactly this inequality. Afterwards $v_1$ is still active
($\tilde D(v_1)-1\ge1$), so Lemma~\ref{int:lem-fan7} absorbs the
\emph{full} star $F$, including $uv_1$, towards $u$; and the (L) step
runs as above, $u$'s budget being untouched by the first step. The
endpoint count at $x$ ends at $\tilde D(x)+1$, which is even, consistent
with $x$ being even in $G$; all $k+1$ deleted edges are restored and
$|D|=|\tilde D|\le B$.

(e) Every component of $\tilde G$ lies in the class of
Theorem~\ref{int:thm-fan}: it is triangle-free (an induced subgraph of
$\EV(G)$) and every E-degree is at most $3$ (the private EV-neighbours of
$x$ by (A2) at $x$, the far vertices because $u$ and $x$ are the only
high vertices, $v_1$ because $\dEV(v_1)=3$). Component-wise
Theorem~\ref{int:thm-fan} with Lemma~\ref{int:lem-sum} gives
$\pn(\tilde G)\le\sum_i\floor{n_i/2}=(n-c_{\mathrm{odd}})/2$, where
$c_{\mathrm{odd}}\ge2$ is the number of odd-order components (it is even
when $n$ is even, and odd and $\ge3$ when $n$ is odd, giving
$(n-3)/2=(n-1)/2-1=B-1$). In both cases $\pn(\tilde G)\le B-1$: route
(a).
\end{proof}

\begin{corollary}[the corner, formally]\label{com:cor-corner}
Let $G$ be a dirty-class graph with $\pn(G)>\floor{n/2}$. Then for
\emph{every} common EV-neighbour $v_j$ (with third EV-neighbour $z_j$ and
$\tilde G_j=G-F-\{xv_j\}$), all of the following hold:
$\pn(\tilde G_j)=\floor{n/2}$; $\beta(\tilde G_j,x)=\floor{n/2}+1$; $z_j$
is buried (endpoint count $0$) in every optimal decomposition of
$\tilde G_j$; in every optimal decomposition $\tilde D$, either
$\tilde D(v_j)=0$, or $\tilde D(v_j)\ge2$ and at least $\tilde D(v_j)$
private EV-neighbours of $x$ have $\tilde D=0$; and $\tilde G_j$ has at
most one component of odd order.
\end{corollary}

\begin{proof}
This is the contrapositive of Theorem~\ref{com:thm-escape}, route by
route, for each $j$. First $\pn(\tilde G_j)\le\floor{n/2}$ always holds:
every E-degree in $\tilde G_j$ is at most $3$ (as in the proof of route
(e)), so component-wise Theorem~\ref{int:thm-fan} with
Lemma~\ref{int:lem-sum} applies. Route (a) failing then gives
$\pn(\tilde G_j)=\floor{n/2}$. Route (b) failing, with endpoint counts at
$x$ odd (Lemma~\ref{int:lem-LB1}(iii), $x$ being odd in $\tilde G_j$),
gives $\tilde D(x)=1$ in every optimal decomposition; cutting a path of
an optimal decomposition at the internal vertex $x$ (legal since
$\deg_{\tilde G_j}(x)=\deg_G(x)-1\ge5$) yields $\floor{n/2}+1$ paths with
endpoint count $3$ at $x$, so $\beta(\tilde G_j,x)=\floor{n/2}+1$. Route
(c) failing gives $\tilde D(z_j)=0$ in every optimal decomposition.
Route (d) failing gives the fourth condition; when $l-t\le1$ it
degenerates to $\tilde D(v_j)=0$, since at most one private EV-neighbour
of $x$ can be dead while $\tilde D(v_j)\ge2$ would need at least two.
Route (e) failing gives at most one odd-order component.

(A note on conventions: the $\beta$ here is the $D(x)\ge2$ variant of
Definition~\ref{int:def-budget}. The $D(x)\ge3$ variant takes the same
value $\floor{n/2}+1$: its lower bound comes from route (b) failing
exactly as above, and its upper bound from the same cutting construction,
which raises the endpoint count at $x$ to $3$.)
\end{proof}

% end sec-computation.tex

% ============================================================================
% sec-status.tex — Formalization status, limitations, and open problems.
% Label prefix: sta: (all labels owned by this file).
% ============================================================================
\section{Formalization status, limitations, and open problems}\label{sec:open}

\subsection{The Lean formalization, item by item}\label{sta:subsec-leanitems}

Table~\ref{sta:tab-leanitems} lists every Lean declaration this paper
cites that is free of \texttt{sorryAx}, with its exact statement, source
anchor and axiom profile; every entry of the axiom column was audited by
\texttt{\#print axioms} against the built library, and the audit was
re-run against the current sources (identical output; the audited files
byte-identical to the review snapshot; Section~\ref{com:subsec-lean}).
Declarations whose proofs consume \texttt{native\_decide} carry the
kernel-reflection axiom \texttt{Lean.ofReduceBool} explicitly and are
quoted as ``machine-checked by kernel reduction''
(Section~\ref{com:subsec-lean}); they are not equated with pure proofs.

\begin{table}[ht]
\centering\small
\caption{The nine \texttt{sorryAx}-free Lean items cited by this paper.
Axiom profiles: $[\mathrm{S}]:=$ \texttt{[propext, Classical.choice,
Quot.sound]} (standard), $[\mathrm{P},\mathrm{Q}]:=$
\texttt{[propext, Quot.sound]}, and $\mathrm{nd}:=$ the
\texttt{Lean.ofReduceBool} kernel-reflection axiom instance. Line anchors
refer to the released sources where the file is released;
\texttt{CompleteOdd.lean} and \texttt{Cert8.lean} are not part of the
released snapshot and their anchors refer to the audited development
snapshot of the project.}
\label{sta:tab-leanitems}
\begin{tabular}{p{3.55cm} p{4.55cm} p{2.7cm} p{2.0cm}}
\hline
Declaration & Statement (exactly as proved) & Anchor & Axioms\\
\hline
\texttt{pNumK4n\_odd\_all} & for every odd $n\ge3$:
$\pn(K_{4,n})=(n+1)/2$; hypotheses $3\le n$, $n$ odd & \texttt{T273.lean}
:976--978 & $[\mathrm{S}]$\\
\texttt{cor\_pNumK4n\_5/13/19} & $\pn(K_{4,5})=3$,
$\pn(K_{4,13})=7$, $\pn(K_{4,19})=10$ & \texttt{T273.lean} :981--987 &
$[\mathrm{S}]$\\
\texttt{pCompleteOdd} & for $k\ge1$: $\pn(K_{2k+1})=k+1$; hypothesis
$k\ne0$ & \texttt{CompleteOdd.lean} :1590--1596 & $[\mathrm{S}]$\\
\texttt{thmA2663\_upgraded} & under hypotheses
$\mathrm{hd2}$ (target side has the prescribed size) and
$\mathrm{val}\le1$: $g_{13}\le d_2+1$ --- the $\mathrm{val}\le1$
hypothesis is \emph{not} proved by the theorem &
\texttt{ThmA2663.lean} :197--200 & $[\mathrm{P},\mathrm{Q}]$\\
\texttt{case2\_struct\_alpha\_iff} & under six explicit hypotheses:
$C_1((t_1+g)\bmod N)=\texttt{true}\iff t_1=s$ & \texttt{Struct2663.lean}
:260--265 & $[\mathrm{P},\mathrm{Q}]$\\
\texttt{ledger6\_ok} & $\mathrm{ledger}_6=(5{,}724;\,2{,}613;\,190;\,398;\,588)$
(N=6 five-column global account) & \texttt{Cert2663.lean} :113 &
$[\mathrm{P},\mathrm{nd}]$\\
\texttt{ledger8\_ok} & $\mathrm{ledger}_8=(701{,}360;\,285{,}372;\,20{,}308;\,10{,}452;\,30{,}760)$
(N=8; dual-route, second Python implementation pinned) &
\texttt{Cert8.lean} :52 & $[\mathrm{P},\mathrm{nd}]$\\
\texttt{l2b\_toplevel\_certificate} & the $79{,}200$-point packing-table
certificate ($120$ orbits $\times$ $330$ demands $\times$ two sides),
pointwise iff & \texttt{Pi3L2b.lean} :248--249 & $[\mathrm{P},\mathrm{nd}]$\\
\texttt{Pi3Uni} (nine theorems) & downset tables (both sides); uniform
layer; the no-threshold witness refuting any ``feasible $\iff
f(\Sigma d)$'' law; top-only failures ($1{,}761$); top-failure set an
upset; orbit counts $110/2$ & \texttt{Pi3Uni.lean} :57--115 &
$[\mathrm{P},\mathrm{nd}]$\\
\hline
\end{tabular}
\end{table}

\emph{Not formalized.} Three boundary entries complete the picture.
Theorem~\ref{thm:conchain} (odd clique chains) and its engine
Theorem~\ref{thm:conL0} have \emph{no} Lean statement: they are paper
proofs, and the machine anchors cover the underlying decompositions, not
the chaining argument (Section~\ref{subsec:conformal}). The Fan-type chain
(Theorems~\ref{thm:claAstar},~\ref{thm:lasep}) carries
\texttt{sorryAx}-dependent declarations --- formalization in progress, and
never quoted as machine-verified (Remark~\ref{rem:claformal}). The
research lemma \texttt{fan6\_nf} is never cited: its statement was
refuted by a machine counterexample and it survives in the sources as a
record only.

\emph{The four-domain account layer and its honest debt list.} The
termination bound's four-domain witness accounts
($588/30{,}760/634{,}960/7{,}798{,}903$ instances at $N=6,8,10,12$;
Section~\ref{com:subsec-lean}) are kernel-checked equalities at the level
of the domain counts and of the two certified ledgers; the theorem body
itself is a paper proof with no Lean statement. Five residuals are
registered against the account layer, none claimed closed:
(R1) the fast$\equiv$slow enumeration equivalence has no general-$N$
theorem (it is pinned by double-scale agreement at the $N=6/8$ anchors and
across all bucket tables); (R2) the full witness tables at $N=10$
($12{,}211{,}196$ rows) and $N=12$ ($313{,}315{,}956$ rows) are not
certified --- bucket-first sampling and the aggregate bucketed ledgers
cover them; (R3) the trailing four digits of the $N=10$/$N=12$ tuples have
no kernel-level \emph{third} implementation (the Lean pins and the Python
second implementation make two); (R4) the $h\bmod 2g$ keys bridge is
proved at zero axioms for its instances, but its general-$N$ derivation is
open; (R5) the declaration-listing verifier is a text-level instrument
whose statement-level boundary is carried as such.

\subsection{Limitations}\label{sta:subsec-limitations}

\subsubsection{The budget-conjecture residuals}\label{sta:subsubsec-kres}

\paragraph{The degree-two cut-vertex branch.} The degree-two layer of the
budget conjecture assembles through a case analysis whose cut-vertex
branch ($G-x$ disconnected) is closed at the \emph{computational} level
only: $1{,}870$ cut-vertex even-degree instances over the base domain of
Table~\ref{com:tab-kdomains} ($492$ with $d=2$, $1{,}165$ with $d=4$,
$213$ with $d=6$) show zero violations, and a dedicated connected
$n\le7$ check covers all $106$ degree-$2$ cut-vertex instances, again
with zero violations. The theorization is pending, and the drafted
splicing argument has a genuine gap: its parity arithmetic assumes
$\lceil n_1/2\rceil+\lceil n_2/2\rceil=n/2$ for the two components of
$G-x$, which holds only when $n$ is even (components of opposite
parity). For odd $n$ the components share parity: the same-odd split
gives $(n+1)/2$ and the same-even split gives $(n-1)/2$, so the drafted
equality fails by one --- the missing parity case must be supplied
before the branch is a theorem. The unconditional ceiling survives
unchanged: $\lceil n_1/2\rceil+\lceil n_2/2\rceil\le\lceil n/2\rceil$
for all parity splits, verified exhaustively for $n\le30$.

\paragraph{The degree-$\ge4$ layer.} The general degree-$\ge4$ layer of
the budget conjecture is \emph{not} discharged. The exact criterion
(Section~\ref{int:subsec-budget}) characterizes one construction tool ---
the degree-insertion operations --- and the alternative-construction
theory for the $208$ tool-dead instances of
Table~\ref{com:tab-kdomains} ($180$ of $208$ cut-vertex; only $7$ with
the reduced case $\sigma=0$) is not written; the enumeration is confined
to the $m\le12$ sparse layers. The residual proof obligation is booked,
and no closure language is used for this layer.

\paragraph{The recorded dead zone.} The $42$ recorded dead-zone instances
of the base domain are discharged at \emph{witness} level: each admits a
machine-verified sharp decomposition with $\beta\le\lceil n/2\rceil$
(equality in $41$ of $42$; the sole strict case has $\beta=3<4$),
attributed as $13$ two-probe artifacts, $23$ cut-vertex witnesses and $6$
tight non-cut instances. This is a finite-instance disposition and does
not eliminate the corresponding premise of the conditional chain, which
remains open as stated (Section~\ref{int:subsec-conditional}).

\paragraph{Corrected paper-level values of the extended account.} The
extended-domain verification records corrected values, which supersede the
first-pass report: $1{,}279$ witness certificates ($484+795$; Section~\ref{com:subsec-kline}); the tight-layer validation
splits as $147$ even-$n$ instances for the tight-layer corollary proper
and $2{,}410$ odd-$n$ instances for its companion characterization
(the literal odd-$n$ proposition has $216$ counterexamples and is
superseded by the companion form ``success $\iff M(u)=d$ or
$M(u-1)\ge d-2$''); the insertion-depth account volume is $175{,}244$
under the corrected depth formula; and in the base domain's $d=6$
slice, $157$ instances are rescued ($156$ at insertion depth $2$, $1$
at depth $3$), zero otherwise.

\paragraph{The $K_8@28$ registered prediction.} For the placeability
programme: the phrase ``surviving regimes only $\{q=0,q=7\}$'' at
$K_8@28$ is a \emph{registered prediction}, not a measurement
(Remark~\ref{rem:claregistered}); the actual census found zero surviving
regimes on every line.

\paragraph{The Lov\'asz 1968 external step.} The all-odd support consumed
by Theorem~\ref{int:thm-Kgallai} rests on Theorem~\ref{int:thm-lovasz},
verified at statement level by five concordant published sources and
carried by the entire subsequent literature. Its offline proof
reconstruction was, however, judged \emph{not achieved}: the skeleton's
deletion step inflates the path count and its re-merging is blocked by
shared internal vertices, a failure pinned mechanically by $22$ witness
graphs on $n\le5$. Two transcription blemishes are on record alongside
(the equality reading ``$|D|=\lfloor n/2\rfloor$'' must be ``at
most''; the ``$\le n-1$ paths'' corollary is loose for even $n$), and
the consumed upper-bound-inflation line is cited in its corrected
upper-bound form, not as an identity. Page-by-page reading of the 1968
original remains mandatory before the external flag is lifted
(Section~\ref{int:subsec-conditional}).

\subsubsection{The generic-account limitations}\label{sta:subsubsec-genres}

The key-table ledger of Section~\ref{com:subsec-ledger} carries its four
conventions as limitations: the $N=14$ affine-normalized domain
equivalence is unproven; key-weighted and instance-weighted accountings
are separate; the no-$V_7$ reading is a one-way superset; and the
machine face of the omitted branch is frozen (its legacy
candidate-generation sources are lost, with the residual mass precisely
measured at $340$ positions for $N=8$ and $2{,}267{,}655$ at $N=14$).
Two further accounting rules bind quotations: per-instance accounts at
$N=8,9$ are listed separately from the key accounts, and the two
accounting layers of the falsification experiments ($12/900$ failing
instances versus $54/1{,}501$ failing witnesses;
Section~\ref{subsubsec:criteria}) are never interchanged. No ledger
result is extrapolated to $N\ge15$ or read across the two domain
conventions.

\subsubsection{The classification-programme boundaries}\label{sta:subsubsec-clares}

The placeability classification programme
(Section~\ref{subsec:echar}) is bounded as follows. The evidence
structure of the exception-set candidate \eqref{eq:claechar} is closed to
$k\le16$ on the empty-side and to $n=13$ on the tower side, with the
$k=17$ full-type sweep at $53{,}394{,}755$ types; the layers $k\ge18$
and the uniqueness side (closed only at $k=15$, refuted as a literal
statement at $k=21$) are open (Remark~\ref{rem:claregistry}). The
prime-family theorem covers the non-infinity regimes only; the infinity
segment of the prime family is open for every member
(Remark~\ref{rem:claprimscope}). The corrected Fourier identity
(Theorem~\ref{thm:clafourier}) is the only form of the character equation
cited anywhere in this paper, and all machine results of the programme
are independent of it (Remark~\ref{rem:clafourierfix}).

\subsubsection{The existence half of the criterion}\label{sta:subsubsec-gapex}

The general-$N$ existence half of the valuation criterion in its so1
form (Conjecture~\ref{int:conj-gapex}) is open: the per-$N$ closures at
$N\in\{6,8,10,12\}$ are theorem-form, the general statement is not, and
no sentence of this paper states otherwise. Three empirical obstacles are
registered against it --- the $\lambda=1$ dead corners ($2{,}528$
instances at $N=10$ with no one-leak witness), the failed-completion
triple $31/1{,}094/13{,}488$, and four falsified discriminators --- and
the candidate mechanisms (near-saturation normal forms; the leak-repair
channel; per-line board normal forms) each carry its own falsification
experiment. The associated $\mu$ $k=2$ stratified existence wall is open
on its paper face (machine face closed on seven domains; no logical
dependency on the family-image characterization;
Remark~\ref{rem:clabridge}).

\subsubsection{The GEN existence programme boundary}\label{sta:subsubsec-flicker}

The GEN existence programme carries one further measured boundary. Its
zero-dodge complete-walk counts on the $w_1=r-5$ flicker line are the
six exact values
\[
158,\;1{,}733,\;11{,}466,\;81{,}660,\;484{,}578,\;3{,}670{,}846
\]
for $r=19,21,23,25,27,29$ (the $r=29$ value exact but
class-bound-conditioned, quoted with that caveat) --- six tiers, not a
closed sequence. One further exact single point, $23{,}933{,}468$ at
$r=31$, stands on its own at exhaustive grade; at $r=33$ there is
\emph{no} count (a quantified hard wall; a segmented long run was started
and reports none). The even-$r$ blocking companion clears the
zero-dodge side only: complete walks with a positive zero-defect exist at
$r=20,22$ in the record and are untouched by it. The existence theorem
itself (registry grade, $r\ge43$) is supported by paper certificates and
step-feasibility, with its machine anchors corroborating only.

\subsection{Open problems}\label{sta:subsec-open}

The open problems below supersede the numbering of any earlier working
list; each is stated with what the present account gives and what it
withholds.

\begin{enumerate}[label=\textbf{O\arabic*.},leftmargin=2.4em]
\item \textbf{The remaining even columns of Case~4}
(Conjecture~\ref{conj:con27A}). Is
$\pn(K_{n_1,n_2})=\lceil n_1/2\rceil$ for all odd $n_1$ and even $n_2$,
$2\le n_2<n_1$? The ring-closure theorem
(Theorem~\ref{thm:conringclosure}) reduces each column to $h+1$ base
instances; the columns $n_2=2,4,6$ are closed, and five named bases are
open: $K_{15,8}$, $K_{17,8}$ (column $n_2=8$) and $K_{17,10}$,
$K_{19,10}$, $K_{21,10}$ (column $n_2=10$). The recorded negatives on
these bases are truncations, not exclusions
(Observation~\ref{obs:concostwall}): the decompositions may lie deeper in
the search tree, or the conjecture may fail. The next methodological step
suggested by the data is a constraint-propagation search driven by the
deficit ledger of Theorem~\ref{thm:constructgen} and the endpoint parity,
rather than deeper node budgets.
\item \textbf{Removing the structural conditions of the Fan-type
theorems} (Conjecture~\ref{conj:claAstar}). The alignment and separation
hypotheses fail exactly at \emph{remote high vertices}; by
Proposition~\ref{prop:claresidue}, a minimal counterexample must have a
bad partner at every high vertex. The remaining corner of the two-high
layer is blocked by the excess bound (zero violations on all $235$
computed instances), and the residue of the conjecture is that corner
together with graphs of at least three high vertices that are neither
separated nor covered by the profile parity. Note that the band-free
phenomenon is family-specific: outside the constructed family, the
general band-free statement is refuted by certified counterexamples
(Section~\ref{subsec:astar}).
\item \textbf{The even-$j$ Walecki residual.} Theorem~\ref{thm:conL0} is
for odd $j$; for even $j$ (blocks of order $1\bmod4$), the chain
construction would extend to all odd-order complete-block chains. In-chain
existence is machine-verified at the block values $K_5$ and $K_9$, but no
general construction is known. The natural route is a Hamilton
decomposition of $K_{2j+1}-C^*$ for a suitable cycle $C^*$ through the
prescribed path, which touches the Oberwolfach problem for odd orders.
\item \textbf{The general-$N$ existence half of the criterion}
(Conjecture~\ref{int:conj-gapex}). For all $N$ and all so1 instances with
capacity slack $\beta\ge2$, is there an admissible pair with
$\mathrm{val}\le1$? The $C_1$ side is closed for all $N$ by the
constructive one-leak theorem; the obstruction is $C_2$-side
reachability, with the three registered empirical obstacles and the
candidate mechanisms of Section~\ref{sta:subsubsec-gapex}, each carrying
its falsification experiment.
\item \textbf{The $\mu$ $k=2$ stratified existence wall.} The paper core
is the existence of $C_1$ with $q\le d_3+1$ and realized. The machine
face is closed on seven domains for the $k=2$ layer (zero violations);
the paper face has a recorded obstruction (no base-to-step mechanism);
and the problem has no logical dependency on the family-image
characterization (Remark~\ref{rem:clabridge}).
\item \textbf{The degree-$\ge4$ layer of the budget conjecture
$K$} (Definition~\ref{int:def-budget}). Beyond degree~$2$ the
deletion--reinsertion scheme meets its true wall; the computational
record shows the wall is an artefact of the method, not of $K$
(Section~\ref{sta:subsubsec-kres}). What is needed is a construction that
produces a $\lceil n/2\rceil$-decomposition with $D(x)\ge2$ directly,
without deleting $x$.
\item \textbf{The exception-set programme.} Beyond the closed faces of
Remark~\ref{rem:claregistry}: the layers $k\ge18$; uniqueness of the
minimal core for $k\ge21$; the prime-family infinity segment
(Remark~\ref{rem:claprimscope}); and the counting-side target of the
$\alpha''$ bridge (one-way and rigidity-free, machine base eight odd
points, $p\ge19$ unenumerated; Section~\ref{subsec:echar}).
\item \textbf{The $4$-degenerate residual.} The theorem of
\cite{ZLH24} covers $3$-degenerate graphs; the class of
$4$-degenerate graphs is, as far as we know, entirely open.
$K_{4,n}$ shows the interesting behaviour starts immediately: it is
$4$-degenerate and, by Theorem~\ref{thm:conk4odd}, beats its published
bounds by one.
\end{enumerate}

\subsection{Concluding remarks}\label{sta:subsec-conclusion}

\paragraph{What is proved.} On the constructive side: the exact value
$\pn(K_{4,n})=(n+1)/2$ for odd $n\ge3$ (two of whose instances are
formalized circumstances of one Lean theorem), the full $n_2=6$ column
and the structural reduction of the remaining even columns of Case~4 of
\cite{CE18}; the odd clique chain theorem
(Theorem~\ref{thm:conchain}) at its all-odd-$j$ domain; and the Fan-type
class theorems admitting a block of $\EV(G)$ of maximum degree at
least~$4$ (Theorems~\ref{thm:claAstar} and~\ref{thm:lasep}), at paper
level with formalization in progress. On the conditional side: the
budget framework, the implication $K\Rightarrow$ Gallai
(Theorem~\ref{int:thm-Kgallai}), the conditional chain, and the
descriptive registration of the eight-item premise ledger
(Theorem~\ref{int:thm-ledger}).

\paragraph{What is conditional.} The conditional route is
restated one last time in its honest form: the criterion's first
conjunct is a machine-closed candidate on explicitly tested domains with
a paper status open; every one of the eight premise conjuncts is an open
conjecture, an unverified external step, or a paper-open item with
machine backing; and no premise of the ledger is claimed, here or
elsewhere, to be within reach of the machine layer except as recorded.
The Lean layer sharpens checkability, not the claim surface: its
\texttt{sorryAx}-free theorems are those of
Table~\ref{sta:tab-leanitems}, its \texttt{native\_decide} certificates
are quoted with the kernel-reduction disclosure, and everything outside
the axiom probe is explicitly not claimed formalized
(Section~\ref{com:subsec-lean}).

\paragraph{What is measured.} The measured layer is reported at its
registered strengths throughout: the five exhaustively enumerated domains
of the budget conjecture (Table~\ref{com:tab-kdomains}; zero violations,
$208$ tool-dead instances all satisfying the conjecture); the nine-domain
key ledger under its four conventions
(Section~\ref{com:subsec-ledger}); the per-$N$ criterion closures with
zero residue on $588/30{,}760/634{,}960/7{,}798{,}903$ instances at
$N=6,8,10,12$; and the classification programme's four complete-graph
layers at their stated tiers (Section~\ref{subsubsec:firstlayers}). The
overall conjecture --- Gallai's path decomposition conjecture --- remains
open. What this paper provides is the verified portion of the landscape
around it, a conditional route whose every edge is checkable, and a
measurement apparatus whose every citation carries its reading
conventions (the four ledger conventions of Section~\ref{com:subsec-ledger}
and the tier qualifiers of Sections~\ref{sec:computation}
and~\ref{sta:subsec-leanitems}), so that the paper remains checkable
against exactly the evidence it quotes.

\paragraph{The no-claim constraints, restated.} The six no-claim
constraints of Section~\ref{int:subsec-claims} bind this closing section
as they bind the whole account; in the order stated there, and matching
the introduction's wording in content: no
closure at $L\ge13$ (candidate level only); no general-$N$ paper closure
of the covering identity (the existence half open); no completeness
claim for the four-piece (E)+(R) landing; no paper closure of the
general-$L$ uniform-construction account (the capacity-starved branch of
the channel-or-menu dichotomy); no full closure of that dichotomy
programme; and no
unconditional theorem-form statement of the valuation criterion --- its
Lean formalization is a bound under an unproved hypothesis, and a
criterion verified over finitely many $N$ is not an existence proof for
general $N$.

% end sec-status.tex


\bibliographystyle{plain}
\bibliography{refs}

\end{document}
